%% signal.tex
%%
%% Copyright (C) 2000-2018 Simone Piccardi.  Permission is granted to
%% copy, distribute and/or modify this document under the terms of the GNU Free
%% Documentation License, Version 1.1 or any later version published by the
%% Free Software Foundation; with the Invariant Sections being "Un preambolo",
%% with no Front-Cover Texts, and with no Back-Cover Texts.  A copy of the
%% license is included in the section entitled "GNU Free Documentation
%% License".
%%

\chapter{I segnali}
\label{cha:signals}

I segnali sono il primo e più semplice meccanismo di comunicazione nei
confronti dei processi. Nella loro versione originale essi portano con sé
nessuna informazione che non sia il loro tipo; si tratta in sostanza di
un'interruzione software portata ad un processo.

In genere essi vengono usati dal kernel per riportare ai processi situazioni
eccezionali (come errori di accesso, eccezioni aritmetiche, ecc.) ma possono
anche essere usati come forma elementare di comunicazione fra processi (ad
esempio vengono usati per il controllo di sessione), per notificare eventi
(come la terminazione di un processo figlio), ecc.

In questo capitolo esamineremo i vari aspetti della gestione dei segnali,
partendo da una introduzione relativa ai concetti base con cui essi vengono
realizzati, per poi affrontarne la classificazione a secondo di uso e modalità
di generazione fino ad esaminare in dettaglio le funzioni e le metodologie di
gestione avanzate e le estensioni fatte all'interfaccia classica nelle nuovi
versioni dello standard POSIX.


\section{Introduzione}
\label{sec:sig_intro}

In questa sezione esamineremo i concetti generali relativi ai segnali, vedremo
le loro caratteristiche di base, introdurremo le nozioni di fondo relative
all'architettura del funzionamento dei segnali e alle modalità con cui il
sistema gestisce l'interazione fra di essi ed i processi.


\subsection{I concetti base}
\label{sec:sig_base}

Come il nome stesso indica i segnali sono usati per notificare ad un processo
l'occorrenza di un qualche evento. Gli eventi che possono generare un segnale
sono vari; un breve elenco di possibili cause per l'emissione di un segnale è
il seguente:

\begin{itemize*}
\item un errore del programma, come una divisione per zero o un tentativo di
  accesso alla memoria fuori dai limiti validi;
\item la terminazione di un processo figlio;
\item la scadenza di un timer o di un allarme;
\item il tentativo di effettuare un'operazione di input/output che non può
  essere eseguita;
\item una richiesta dell'utente dal terminale di terminare o fermare il
  programma.
\item l'invio esplicito da parte del processo stesso o di un altro.
\end{itemize*}

Ciascuno di questi eventi, compresi gli ultimi due che pure sono controllati
dall'utente o da un altro processo, comporta l'intervento diretto da parte del
kernel che causa la generazione di un particolare tipo di segnale.

Quando un processo riceve un segnale, invece del normale corso del programma,
viene eseguita una azione predefinita o una apposita funzione di gestione che
può essere stata specificata dall'utente, nel qual caso si dice che si
\textsl{intercetta} il segnale. Riprendendo la terminologia originale da qui
in avanti faremo riferimento a questa funzione come al \textsl{gestore} del
segnale, traduzione approssimata dell'inglese \textit{signal handler}.


\subsection{Le \textsl{semantiche} del funzionamento dei segnali}
\label{sec:sig_semantics}

Negli anni il comportamento del sistema in risposta ai segnali è stato
modificato in vari modi nelle differenti implementazioni di Unix.  Si possono
individuare due tipologie fondamentali di comportamento dei segnali (dette
\textsl{semantiche}) che vengono chiamate rispettivamente \textsl{semantica
  affidabile} (o \textit{reliable}) e \textsl{semantica inaffidabile} (o
\textit{unreliable}).

Nella \textsl{semantica inaffidabile}, che veniva implementata dalle prime
versioni di Unix, la funzione di gestione del segnale specificata dall'utente
non restava attiva una volta che era stata eseguita; era perciò compito
dell'utente ripetere l'installazione dello stesso all'interno del
\textsl{gestore} del segnale in tutti quei casi in cui si voleva che esso
restasse attivo.

\begin{figure}[!htbp]
  \footnotesize \centering
  \begin{minipage}[c]{\codesamplewidth}
    \includecodesample{listati/unreliable_sig.c}
  \end{minipage} 
  \normalsize 
  \caption{Esempio di codice di un gestore di segnale per la semantica
    inaffidabile.} 
  \label{fig:sig_old_handler}
\end{figure}

In questo caso però è possibile una situazione in cui i segnali possono essere
perduti. Si consideri il segmento di codice riportato in
fig.~\ref{fig:sig_old_handler}: nel programma principale viene installato un
gestore (\texttt{\small 5}), la cui prima operazione (\texttt{\small 11}) è
quella di reinstallare se stesso. Se nell'esecuzione del gestore fosse
arrivato un secondo segnale prima che esso abbia potuto eseguire la
reinstallazione di se stesso per questo secondo segnale verrebbe eseguito il
comportamento predefinito, il che può comportare, a seconda dei casi, la
perdita del segnale (se l'impostazione predefinita è quella di ignorarlo) o la
terminazione immediata del processo; in entrambi i casi l'azione prevista dal
gestore non verrebbe eseguita.

Questa è la ragione per cui l'implementazione dei segnali secondo questa
semantica viene chiamata \textsl{inaffidabile}: infatti la ricezione del
segnale e la reinstallazione del suo gestore non sono operazioni atomiche, e
sono sempre possibili delle \textit{race condition} (si ricordi
sez.~\ref{sec:proc_multi_prog}).  Un altro problema è che in questa semantica
non esiste un modo per bloccare i segnali quando non si vuole che arrivino; i
processi possono ignorare il segnale, ma non è possibile istruire il sistema a
non fare nulla in occasione di un segnale, pur mantenendo memoria del fatto
che è avvenuto.

Nella semantica \textsl{affidabile} (quella utilizzata da Linux e da ogni Unix
moderno) il gestore una volta installato resta attivo e non si hanno tutti i
problemi precedenti. In questa semantica i segnali vengono \textsl{generati}
dal kernel per un processo all'occorrenza dell'evento che causa il segnale. In
genere questo viene fatto dal kernel impostando un apposito campo della
\struct{task\_struct} del processo nella \textit{process table} (si veda
fig.~\ref{fig:proc_task_struct}).

Si dice che il segnale viene \textsl{consegnato} al processo (dall'inglese
\textit{delivered}) quando viene eseguita l'azione per esso prevista, mentre
per tutto il tempo che passa fra la generazione del segnale e la sua consegna
esso è detto \textsl{pendente} (o \textit{pending}). In genere questa
procedura viene effettuata dallo \textit{scheduler} quando, riprendendo
l'esecuzione del processo in questione, verifica la presenza del segnale nella
\struct{task\_struct} e mette in esecuzione il gestore.

In questa semantica un processo ha la possibilità di bloccare la consegna dei
segnali, in questo caso, se l'azione per il suddetto segnale non è quella di
ignorarlo, il segnale resta \textsl{pendente} fintanto che il processo non lo
sblocca (nel qual caso viene consegnato) o imposta l'azione corrispondente per
ignorarlo.

Si tenga presente che il kernel stabilisce cosa fare con un segnale che è
stato bloccato al momento della consegna, non quando viene generato; questo
consente di cambiare l'azione per il segnale prima che esso venga consegnato,
e si può usare la funzione \func{sigpending} (vedi sez.~\ref{sec:sig_sigmask})
per determinare quali segnali sono bloccati e quali sono pendenti.

Infine occorre precisare che i segnali predatano il supporto per i
\textit{thread} e vengono sempre inviati al processo come insieme, cosa che
può creare incertezza nel caso questo sia multi-\textit{thread}. In tal caso
quando è possibile determinare quale è il \textit{thread} specifico che deve
ricevere il segnale, come avviene per i segnali di errore, questo sarà inviato
solo a lui, altrimenti sarà inviato a discrezione del kernel ad uno qualunque
dei \textit{thread} del processo che possa riceverlo (che cioè non blocchi il
segnale), torneremo sull'argomento in sez.~\ref{sec:thread_signal}.

\subsection{Tipi di segnali}
\label{sec:sig_types}

In generale si tende a classificare gli eventi che possono generare dei
segnali in tre categorie principali: errori, eventi esterni e richieste
esplicite.

Un errore significa che un programma ha fatto qualcosa di sbagliato e non può
continuare ad essere eseguito. Non tutti gli errori causano dei segnali, in
genere le condizioni di errore più comuni comportano la restituzione di un
codice di errore da parte di una funzione di libreria. Sono gli errori che
possono avvenire nell'esecuzione delle istruzioni di un programma, come le
divisioni per zero o l'uso di indirizzi di memoria non validi, che causano
l'emissione di un segnale.

Un evento esterno ha in genere a che fare con le operazioni di lettura e
scrittura su file, o con l'interazione con dispositivi o con altri processi;
esempi di segnali di questo tipo sono quelli legati all'arrivo di dati in
ingresso, scadenze di un timer, terminazione di processi figli, la pressione
dei tasti di stop o di suspend su un terminale.

Una richiesta esplicita significa l'uso da parte di un programma delle
apposite funzioni di sistema, come \func{kill} ed affini (vedi
sez.~\ref{sec:sig_kill_raise}) per la generazione ``\textsl{manuale}'' di un
segnale.

Si dice poi che i segnali possono essere \textsl{asincroni} o
\textsl{sincroni}. Un segnale \textsl{sincrono} è legato ad una azione
specifica di un programma ed è inviato (a meno che non sia bloccato) durante
tale azione. Molti errori generano segnali \textsl{sincroni}, così come la
richiesta esplicita da parte del processo tramite le chiamate al sistema.
Alcuni errori come la divisione per zero non sono completamente sincroni e
possono arrivare dopo qualche istruzione.

I segnali \textsl{asincroni} sono generati da eventi fuori dal controllo del
processo che li riceve, e arrivano in tempi impredicibili nel corso
dell'esecuzione del programma. Eventi esterni come la terminazione di un
processo figlio generano segnali \textsl{asincroni}, così come le richieste di
generazione di un segnale effettuate da altri processi.

In generale un tipo di segnale o è sincrono o è asincrono, salvo il caso in
cui esso sia generato attraverso una richiesta esplicita tramite chiamata al
sistema, nel qual caso qualunque tipo di segnale (quello scelto nella
chiamata) può diventare sincrono o asincrono a seconda che sia generato
internamente o esternamente al processo.


\subsection{La notifica dei segnali}
\label{sec:sig_notification}

Come accennato quando un segnale viene generato, se la sua azione predefinita
non è quella di essere ignorato, il kernel prende nota del fatto nella
\struct{task\_struct} del processo; si dice così che il segnale diventa
\textsl{pendente} (o \textit{pending}), e rimane tale fino al momento in cui
verrà notificato al processo o verrà specificata come azione quella di
ignorarlo.

Normalmente l'invio al processo che deve ricevere il segnale è immediato ed
avviene non appena questo viene rimesso in esecuzione dallo \textit{scheduler}
che esegue l'azione specificata. Questo a meno che il segnale in questione non
sia stato bloccato prima della notifica, nel qual caso l'invio non avviene ed
il segnale resta \textsl{pendente} indefinitamente.

Quando lo si sblocca un segnale \textsl{pendente} sarà subito notificato. Si
tenga presente però che tradizionalmente i segnali \textsl{pendenti} non si
accodano, alla generazione infatti il kernel marca un flag nella
\struct{task\_struct} del processo, per cui se prima della notifica ne vengono
generati altri il flag è comunque marcato, ed il gestore viene eseguito sempre
una sola volta. In realtà questo non vale nel caso dei cosiddetti segnali
\textit{real-time}, che vedremo in sez.~\ref{sec:sig_real_time}, ma questa è
una funzionalità avanzata che per ora tralasceremo.

Si ricordi inoltre che se l'azione specificata per un segnale è quella di
essere ignorato questo sarà scartato immediatamente al momento della sua
generazione, e questo anche se in quel momento il segnale è bloccato, perché
bloccare su un segnale significa bloccarne la notifica. Per questo motivo un
segnale, fintanto che viene ignorato, non sarà mai notificato, anche se prima
è stato bloccato ed in seguito si è specificata una azione diversa, nel qual
caso solo i segnali successivi alla nuova specificazione saranno notificati.

Una volta che un segnale viene notificato, che questo avvenga subito o dopo
una attesa più o meno lunga, viene eseguita l'azione specificata per il
segnale. Per alcuni segnali (per la precisione \signal{SIGKILL} e
\signal{SIGSTOP}) questa azione è predeterminata dal kernel e non può essere
mai modificata, ma per tutti gli altri si può selezionare una delle tre
possibilità seguenti:

\begin{itemize*}
\item ignorare il segnale;
\item intercettare il segnale, ed utilizzare il gestore specificato;
\item accettare l'azione predefinita per quel segnale.
\end{itemize*}

Un programma può specificare queste scelte usando le due funzioni
\func{signal} e \func{sigaction}, che tratteremo rispettivamente in
sez.~\ref{sec:sig_signal} e sez.~\ref{sec:sig_sigaction}. Se si è installato
un gestore sarà quest'ultimo ad essere eseguito alla notifica del segnale.
Inoltre il sistema farà si che mentre viene eseguito il gestore di un segnale,
quest'ultimo venga automaticamente bloccato, così si possono evitare alla
radice possibili \textit{race condition}.

Nel caso non sia stata specificata un'azione, viene utilizzata la cosiddetta
azione predefinita che, come vedremo in sez.~\ref{sec:sig_standard}, è propria
di ciascun segnale. Nella maggior parte dei casi questa azione comporta la
terminazione immediata del processo, ma per alcuni segnali che rappresentano
eventi innocui l'azione predefinita è di essere ignorati. Inoltre esistono
alcuni segnali la cui azione è semplicemente quella di fermare l'esecuzione
del programma, vale a dire portarlo nello stato di \textit{stopped} (lo stato
\texttt{T}, si ricordi tab.~\ref{tab:proc_proc_states} e quanto illustrato in
sez.~\ref{sec:proc_sched}).

Quando un segnale termina un processo il padre può determinare la causa della
terminazione esaminandone lo stato di uscita così come viene riportato dalle
funzioni \func{wait} e \func{waitpid} (vedi sez.~\ref{sec:proc_wait}). Questo
ad esempio è il modo in cui la shell determina i motivi della terminazione di
un programma e scrive un eventuale messaggio di errore.

\itindbeg{core~dump}

I segnali che rappresentano errori del programma (divisione per zero o
violazioni di accesso) hanno come ulteriore caratteristica della loro azione
predefinita, oltre a terminare il processo, quella di scrivere nella directory
di lavoro corrente del processo di un file \file{core} su cui viene salvata
l'immagine della memoria del processo.

Questo file costituisce il cosiddetto \textit{core dump}, e contenendo
l'immagine della memoria del processo, consente di risalire allo stato dello
\textit{stack} (vedi sez.~\ref{sec:proc_mem_layout}) prima della
terminazione. Questo permette di esaminare il contenuto del file un secondo
tempo con un apposito programma (un \textit{debugger} come \cmd{gdb}) per
investigare sulla causa dell'errore, ed in particolare, grazie appunto ai dati
dello \textit{stack}, consente di identificare quale funzione ha causato
l'errore.

Si tenga presente che il \textit{core dump} viene creato non solo in caso di
errore effettivo, ma anche se il segnale per cui la sua creazione è prevista
nell'azione dell'azione predefinita viene inviato al programma con una delle
funzioni \func{kill}, \func{raise}, ecc.

\itindend{core~dump}


\section{La classificazione dei segnali}
\label{sec:sig_classification}

Esamineremo in questa sezione quali sono i vari segnali definiti nel sistema,
quali sono le loro caratteristiche e la loro tipologia, tratteremo le varie
macro e costanti che permettono di identificarli, e illustreremo le funzioni
che ne stampano la descrizione.


\subsection{I segnali standard}
\label{sec:sig_standard}

Ciascun segnale è identificato dal kernel con un numero, ma benché per alcuni
segnali questi numeri siano sempre gli stessi, tanto da essere usati come
sinonimi, l'uso diretto degli identificativi numerici da parte dei programmi è
comunque da evitare, in quanto essi non sono mai stati standardizzati e
possono variare a seconda dell'implementazione del sistema, e nel caso di
Linux anche a seconda della architettura hardware e della versione del kernel.

Quelli che invece sono stati, almeno a grandi linee, standardizzati, sono i
nomi dei segnali e le costanti di preprocessore che li identificano, che sono
tutte nella forma \texttt{SIG\textsl{nome}}, e sono queste che devono essere
usate nei programmi. Come tutti gli altri nomi e le funzioni che concernono i
segnali, esse sono definite nell'header di sistema \headfile{signal.h}.

\begin{table}[!htb]
  \footnotesize
  \centering
  \begin{tabular}[c]{|l|c|c|l|}
    \hline
    \textbf{Segnale} &\textbf{Standard}&\textbf{Azione}&\textbf{Descrizione} \\
    \hline
    \hline
    \signal{SIGHUP}  &P & T & Hangup o terminazione del processo di 
                              controllo.\\
    \signal{SIGINT}  &PA& T & Interrupt da tastiera (\cmd{C-c}).\\
    \signal{SIGQUIT} &P & C & Quit da tastiera (\cmd{C-y}).\\
    \signal{SIGILL}  &PA& C & Istruzione illecita.\\
    \signal{SIGTRAP} &S & C & Trappole per un Trace/breakpoint.\\
    \signal{SIGABRT} &PA& C & Segnale di abort da \func{abort}.\\
    \signald{SIGIOT}  &B & C & Trappola di I/O. Sinonimo di \signal{SIGABRT}.\\
    \signal{SIGBUS}  &BS& C & Errore sul bus (bad memory access).\\
    \signal{SIGFPE}  &AP& C & Errore aritmetico.\\
    \signal{SIGKILL} &P & T& Segnale di terminazione forzata.\\
    \signal{SIGUSR1} &P & T & Segnale utente numero 1.\\
    \signal{SIGSEGV} &AP& C & Errore di accesso in memoria.\\
    \signal{SIGUSR2} &P & T & Segnale utente numero 2.\\
    \signal{SIGPIPE} &P & T & \textit{Pipe} spezzata.\\
    \signal{SIGALRM} &P & T & Segnale del timer da \func{alarm}.\\
    \signal{SIGTERM} &AP& T & Segnale di terminazione (\texttt{C-\bslash}).\\
    \signal{SIGCHLD} &P & I & Figlio terminato o fermato.\\
    \signal{SIGCONT} &P &-- & Continua se fermato.\\
    \signal{SIGSTOP} &P & S & Ferma il processo.\\
    \signal{SIGTSTP} &P & S & Pressione del tasto di stop sul terminale.\\
    \signal{SIGTTIN} &P & S & Input sul terminale per un processo 
                              in background.\\
    \signal{SIGTTOU} &P & S & Output sul terminale per un processo          
                              in background.\\
    \signal{SIGURG}  &BS& I & Ricezione di una \textit{urgent condition} su 
                              un socket.\\
    \signal{SIGXCPU} &BS& C & Ecceduto il limite sul tempo di CPU.\\
    \signal{SIGXFSZ} &BS& C & Ecceduto il limite sulla dimensione dei file.\\
    \signal{SIGVTALRM}&BS& T& Timer di esecuzione scaduto.\\
    \signal{SIGPROF} &BS& T & Timer del \textit{profiling} scaduto.\\
    \signal{SIGWINCH}&B & I & Finestra ridimensionata (4.3BSD, Sun).\\
    \signal{SIGIO}   &B & T & L'I/O è possibile.\\
    \signal{SIGPOLL} &VS& T & \textit{Pollable event}, sinonimo di
                              \signal{SIGIO}.\\
    \signal{SIGPWR}  &V & T & Fallimento dell'alimentazione.\\
    \signal{SIGSYS}  &VS& C & \textit{system call} sbagliata.\\
    \hline
    \signal{SIGSTKFLT}&?& T & Errore sullo stack del coprocessore (inusato).\\
    \signald{SIGUNUSED}&?& C & Segnale inutilizzato (sinonimo di
                               \signal{SIGSYS}).\\
    \hline
    \signal{SIGCLD}  &V & I & Sinonimo di \signal{SIGCHLD}.\\
    \signal{SIGEMT}  &V & C & Trappola di emulatore.\\
    \signal{SIGINFO} &B & T & Sinonimo di \signal{SIGPWR}.\\
    \signal{SIGLOST} &? & T & Perso un lock sul file, sinonimo
                              di \signal{SIGIO} (inusato).\\
    \hline
  \end{tabular}
  \caption{Lista dei segnali ordinari in Linux.}
  \label{tab:sig_signal_list}
\end{table}

In tab.~\ref{tab:sig_signal_list} si è riportato l'elenco completo dei segnali
ordinari definiti su Linux per tutte le possibili architetture (tratteremo
quelli \textit{real-time} in sez.~\ref{sec:sig_real_time}). Ma si tenga
presente che solo quelli elencati nella prima sezione della tabella sono
presenti su tutte le architetture. Nelle sezioni successive si sono riportati
rispettivamente quelli che esistono solo sull'architettura PC e quelli che non
esistono sull'architettura PC, ma sono definiti sulle architetture
\textit{alpha} o \textit{mips}.

Alcuni segnali erano previsti fin dallo standard ANSI C, ed i segnali sono
presenti in tutti i sistemi unix-like, ma l'elenco di quelli disponibili non è
uniforme, ed alcuni di essi sono presenti solo su alcune implementazioni o
architetture hardware, ed anche il loro significato può variare. Per questo si
sono riportati nella seconda colonna della tabella riporta gli standard in cui
ciascun segnale è stato definito, indicati con altrettante lettere da
interpretare secondo la legenda di tab.~\ref{tab:sig_standard_leg}. Si tenga
presente che il significato dei segnali è abbastanza indipendente dalle
implementazioni solo per quelli definiti negli standard POSIX.1-1990 e
POSIX.1-2001. 

\begin{table}[htb]
  \footnotesize
  \centering
  \begin{tabular}[c]{|c|l|}
    \hline
    \textbf{Sigla} & \textbf{Standard} \\
    \hline
    \hline
    P & POSIX.1-1990.\\
    B & BSD (4.2 BSD e Sun).\\
    A & ANSI C.\\
    S & SUSv2 (e POSIX.1-2001).\\
    V & System V.\\
    ? & Ignoto.\\
    \hline
  \end{tabular}
  \caption{Legenda dei valori degli standard riportati nella seconda colonna
    di tab.~\ref{tab:sig_signal_list}.} 
  \label{tab:sig_standard_leg}
\end{table}

Come accennato in sez.~\ref{sec:sig_notification} a ciascun segnale è
associata una specifica azione predefinita che viene eseguita quando nessun
gestore è installato. Le azioni predefinite possibili, che abbiamo già
descritto in sez.~\ref{sec:sig_notification}, sono state riportate in
tab.~\ref{tab:sig_signal_list} nella terza colonna, e di nuovo sono state
indicate con delle lettere la cui legenda completa è illustrata in
tab.~\ref{tab:sig_action_leg}).

\begin{table}[htb]
  \footnotesize
  \centering
  \begin{tabular}[c]{|c|l|}
    \hline
    \textbf{Sigla} & \textbf{Significato} \\
    \hline
    \hline
    T & L'azione predefinita è terminare il processo.\\
    C & L'azione predefinita è terminare il processo e scrivere un 
        \textit{core dump}.\\
    I & L'azione predefinita è ignorare il segnale.\\
    S & L'azione predefinita è fermare il processo.\\
    \hline
  \end{tabular}
  \caption{Legenda delle azioni predefinite dei segnali riportate nella terza
    colonna di tab.~\ref{tab:sig_signal_list}.}
  \label{tab:sig_action_leg}
\end{table}


Si inoltre noti come \signal{SIGCONT} sia l'unico segnale a non avere
l'indicazione di una azione predefinita nella terza colonna di
tab.~\ref{tab:sig_signal_list}, questo perché il suo effetto è sempre quello
di far ripartire un programma in stato \texttt{T} fermato da un segnale di
stop. Inoltre i segnali \signal{SIGSTOP} e \signal{SIGKILL} si distinguono da
tutti gli altri per la specifica caratteristica di non potere essere né
intercettati, né bloccati, né ignorati.

Il numero totale di segnali presenti è dato dalla macro \macrod{NSIG} (e tiene
conto anche di quelli \textit{real-time}) e dato che i numeri dei segnali sono
allocati progressivamente, essa corrisponde anche al successivo del valore
numerico assegnato all'ultimo segnale definito.  La descrizione dettagliata
del significato dei precedenti segnali, raggruppati per tipologia, verrà
affrontata nei paragrafi successivi.


\subsection{I segnali di errore}
\label{sec:sig_prog_error}

Questi segnali sono generati quando il sistema, o in certi casi direttamente
l'hardware (come per i \textit{page fault} non validi o le eccezioni del
processore) rileva un qualche errore insanabile nel programma in
esecuzione. In generale la generazione di questi segnali significa che il
programma ha dei gravi problemi (ad esempio ha dereferenziato un puntatore non
valido o ha eseguito una operazione aritmetica proibita) e l'esecuzione non
può essere proseguita.

In genere si intercettano questi segnali per permettere al programma di
terminare in maniera pulita, ad esempio per ripristinare le impostazioni della
console o eliminare i file di lock prima dell'uscita.  In questo caso il
gestore deve concludersi ripristinando l'azione predefinita e rialzando il
segnale, in questo modo il programma si concluderà senza effetti spiacevoli,
ma riportando lo stesso stato di uscita che avrebbe avuto se il gestore non ci
fosse stato.

L'azione predefinita per tutti questi segnali è causare la terminazione del
processo che li ha causati. In genere oltre a questo il segnale provoca pure
la registrazione su disco di un file di \textit{core dump}, che un debugger
può usare per ricostruire lo stato del programma al momento della
terminazione.  Questi segnali sono:
\begin{basedescript}{\desclabelwidth{2.0cm}}
\item[\signald{SIGFPE}] Riporta un errore aritmetico fatale. Benché il nome
  derivi da \textit{floating point exception} si applica a tutti gli errori
  aritmetici compresa la divisione per zero e l'overflow.  Se il gestore
  ritorna il comportamento del processo è indefinito, ed ignorare questo
  segnale può condurre ad un ciclo infinito.

%   Per questo segnale le cose sono complicate dal fatto che possono esserci
%   molte diverse eccezioni che \signal{SIGFPE} non distingue, mentre lo
%   standard IEEE per le operazioni in virgola mobile definisce varie eccezioni
%   aritmetiche e richiede che esse siano notificate.
% TODO trovare altre info su SIGFPE e trattare la notifica delle eccezioni 
  
\item[\signald{SIGILL}] Il nome deriva da \textit{illegal instruction},
  significa che il programma sta cercando di eseguire una istruzione
  privilegiata o inesistente, in generale del codice illecito. Poiché il
  compilatore del C genera del codice valido si ottiene questo segnale se il
  file eseguibile è corrotto o si stanno cercando di eseguire dei dati.
  Quest'ultimo caso può accadere quando si passa un puntatore sbagliato al
  posto di un puntatore a funzione, o si eccede la scrittura di un vettore di
  una variabile locale, andando a corrompere lo \textit{stack}. Lo stesso
  segnale viene generato in caso di overflow dello \textit{stack} o di
  problemi nell'esecuzione di un gestore.  Se il gestore ritorna il
  comportamento del processo è indefinito.

\item[\signald{SIGSEGV}] Il nome deriva da \textit{segment violation}, e
  significa che il programma sta cercando di leggere o scrivere in una zona di
  memoria protetta al di fuori di quella che gli è stata riservata dal
  sistema. In genere è il meccanismo della protezione della memoria che si
  accorge dell'errore ed il kernel genera il segnale.  È tipico ottenere
  questo segnale dereferenziando un puntatore nullo o non inizializzato
  leggendo al di là della fine di un vettore.  Se il gestore ritorna il
  comportamento del processo è indefinito.

\item[\signald{SIGBUS}] Il nome deriva da \textit{bus error}. Come
  \signal{SIGSEGV} questo è un segnale che viene generato di solito quando si
  dereferenzia un puntatore non inizializzato, la differenza è che
  \signal{SIGSEGV} indica un accesso non permesso su un indirizzo esistente
  (al di fuori dallo \textit{heap} o dallo \textit{stack}), mentre
  \signal{SIGBUS} indica l'accesso ad un indirizzo non valido, come nel caso
  di un puntatore non allineato.

\item[\signald{SIGABRT}] Il nome deriva da \textit{abort}. Il segnale indica
  che il programma stesso ha rilevato un errore che viene riportato chiamando
  la funzione \func{abort}, che genera questo segnale.

\item[\signald{SIGTRAP}] È il segnale generato da un'istruzione di breakpoint o
  dall'attivazione del tracciamento per il processo. È usato dai programmi per
  il debugging e un programma normale non dovrebbe ricevere questo segnale.

\item[\signald{SIGSYS}] Sta ad indicare che si è eseguita una istruzione che
  richiede l'esecuzione di una \textit{system call}, ma si è fornito un codice
  sbagliato per quest'ultima. 

\item[\signald{SIGEMT}] Il nome sta per \textit{emulation trap}. Il segnale non
  è previsto da nessuno standard ed è definito solo su alcune architetture che
  come il vecchio PDP11 prevedono questo tipo di interruzione, non è presente
  sui normali PC.
\end{basedescript}


\subsection{I segnali di terminazione}
\label{sec:sig_termination}

Questo tipo di segnali sono usati per terminare un processo; hanno vari nomi a
causa del differente uso che se ne può fare, ed i programmi possono
trattarli in maniera differente. 

La ragione per cui può essere necessario intercettare questi segnali è che il
programma può dover eseguire una serie di azioni di pulizia prima di
terminare, come salvare informazioni sullo stato in cui si trova, cancellare
file temporanei, o ripristinare delle condizioni alterate durante il
funzionamento (come il modo del terminale o le impostazioni di una qualche
periferica). L'azione predefinita di questi segnali è di terminare il
processo, questi segnali sono:
\begin{basedescript}{\desclabelwidth{2.0cm}}
\item[\signald{SIGTERM}] Il nome sta per \textit{terminate}. È un segnale
  generico usato per causare la conclusione di un programma. È quello che
  viene generato di default dal comando \cmd{kill}.  Al contrario di
  \signal{SIGKILL} può essere intercettato, ignorato, bloccato. In genere lo
  si usa per chiedere in maniera ``\textsl{educata}'' ad un processo di
  concludersi.

\item[\signald{SIGINT}] Il nome sta per \textit{interrupt}. È il segnale di
  interruzione per il programma. È quello che viene generato di default dal
  dall'invio sul terminale del carattere di controllo ``\textit{INTR}'',
  \textit{interrupt} appunto, che viene generato normalmente dalla sequenza
  \cmd{C-c} sulla tastiera.

\item[\signald{SIGQUIT}] È analogo a \signal{SIGINT} con la differenza che è
  controllato da un altro carattere di controllo, ``\textit{QUIT}'',
  corrispondente alla sequenza \texttt{C-\bslash} sulla tastiera. A differenza
  del precedente l'azione predefinita, oltre alla terminazione del processo,
  comporta anche la creazione di un \textit{core dump}.  In genere lo si può
  pensare come corrispondente ad una condizione di errore del programma
  rilevata dall'utente. Per questo motivo non è opportuno fare eseguire al
  gestore di questo segnale le operazioni di pulizia normalmente previste
  (tipo la cancellazione di file temporanei), dato che in certi casi esse
  possono eliminare informazioni utili nell'esame dei \textit{core dump}.

\item[\signald{SIGKILL}] Il nome è utilizzato per terminare in maniera immediata
  qualunque programma. Questo segnale non può essere né intercettato, né
  ignorato, né bloccato, per cui causa comunque la terminazione del processo.
  In genere esso viene generato solo per richiesta esplicita dell'utente dal
  comando (o tramite la funzione) \cmd{kill}. Dato che non lo si può
  intercettare è sempre meglio usarlo come ultima risorsa quando metodi meno
  brutali, come \signal{SIGTERM} o \cmd{C-c} non funzionano. 

  Se un processo non risponde a nessun altro segnale \signal{SIGKILL} ne causa
  sempre la terminazione (in effetti il fallimento della terminazione di un
  processo da parte di \signal{SIGKILL} costituirebbe un malfunzionamento del
  kernel). Talvolta è il sistema stesso che può generare questo segnale quando
  per condizioni particolari il processo non può più essere eseguito neanche
  per eseguire un gestore.

\item[\signald{SIGHUP}] Il nome sta per \textit{hang-up}. Segnala che il
  terminale dell'utente si è disconnesso, ad esempio perché si è interrotta la
  rete. Viene usato anche per riportare la terminazione del processo di
  controllo di un terminale a tutti i processi della sessione (vedi
  sez.~\ref{sec:sess_job_control}), in modo che essi possano disconnettersi
  dal relativo terminale.  Viene inoltre usato in genere per segnalare ai
  programmi di servizio (i cosiddetti \textsl{demoni}, vedi
  sez.~\ref{sec:sess_daemon}), che non hanno un terminale di controllo, la
  necessità di reinizializzarsi e rileggere il file (o i file) di
  configurazione.
\end{basedescript}


\subsection{I segnali di allarme}
\label{sec:sig_alarm}

Questi segnali sono generati dalla scadenza di un timer (vedi
sez.~\ref{sec:sig_alarm_abort}). Il loro comportamento predefinito è quello di
causare la terminazione del programma, ma con questi segnali la scelta
predefinita è irrilevante, in quanto il loro uso presuppone sempre la
necessità di un gestore.  Questi segnali sono:
\begin{basedescript}{\desclabelwidth{2.0cm}}
\item[\signald{SIGALRM}] Il nome sta per \textit{alarm}. Segnale la scadenza di
  un timer misurato sul tempo reale o sull'orologio di sistema. È normalmente
  usato dalla funzione \func{alarm}.

\item[\signald{SIVGTALRM}] Il nome sta per \textit{virtual alarm}. È analogo al
  precedente ma segnala la scadenza di un timer sul tempo di CPU usato dal
  processo. 

\item[\signald{SIGPROF}] Il nome sta per \textit{profiling}. Indica la scadenza
  di un timer che misura sia il tempo di CPU speso direttamente dal processo
  che quello che il sistema ha speso per conto di quest'ultimo. In genere
  viene usato dagli strumenti che servono a fare la profilazione dell'utilizzo
  del tempo di CPU da parte del processo.
\end{basedescript}


\subsection{I segnali di I/O asincrono}
\label{sec:sig_asyncio}

Questi segnali operano in congiunzione con le funzioni di I/O asincrono. Per
questo occorre comunque usare \func{fcntl} per abilitare un file descriptor a
generare questi segnali.  L'azione predefinita è di essere ignorati. Questi
segnali sono:
\begin{basedescript}{\desclabelwidth{2.0cm}}
\item[\signald{SIGIO}] Questo segnale viene inviato quando un file descriptor è
  pronto per eseguire dell'input/output. In molti sistemi solo i socket e i
  terminali possono generare questo segnale, in Linux questo può essere usato
  anche per i file, posto che la chiamata a \func{fcntl} che lo attiva abbia
  avuto successo.

\item[\signald{SIGURG}] Questo segnale è inviato quando arrivano dei dati
  urgenti o \textit{out-of-band} su di un socket; per maggiori dettagli al
  proposito si veda sez.~\ref{sec:TCP_urgent_data}.

\item[\signald{SIGPOLL}] Questo segnale è definito nella standard POSIX.1-2001,
  ed è equivalente a \signal{SIGIO} che invece deriva da BSD. Su Linux è
  definito per compatibilità con i sistemi System V.
\end{basedescript}


\subsection{I segnali per il controllo di sessione}
\label{sec:sig_job_control}

Questi sono i segnali usati dal controllo delle sessioni e dei processi, il
loro uso è specializzato e viene trattato in maniera specifica nelle sezioni
in cui si trattano gli argomenti relativi.  Questi segnali sono:
\begin{basedescript}{\desclabelwidth{2.0cm}}
\item[\signald{SIGCHLD}] Questo è il segnale mandato al processo padre quando un
  figlio termina o viene fermato. L'azione predefinita è di ignorare il
  segnale, la sua gestione è trattata in sez.~\ref{sec:proc_wait}.

\item[\signald{SIGCLD}] Per Linux questo è solo un segnale identico al
  precedente e definito come sinonimo. Il nome è obsoleto, deriva dalla
  definizione del segnale su System V, ed oggi deve essere evitato.

\item[\signald{SIGCONT}] Il nome sta per \textit{continue}. Il segnale viene
  usato per fare ripartire un programma precedentemente fermato da
  \signal{SIGSTOP}. Questo segnale ha un comportamento speciale, e fa sempre
  ripartire il processo prima della sua consegna. Il comportamento predefinito
  è di fare solo questo; il segnale non può essere bloccato. Si può anche
  installare un gestore, ma il segnale provoca comunque il riavvio del
  processo.
  
  La maggior pare dei programmi non hanno necessità di intercettare il
  segnale, in quanto esso è completamente trasparente rispetto all'esecuzione
  che riparte senza che il programma noti niente. Si possono installare dei
  gestori per far sì che un programma produca una qualche azione speciale
  se viene fermato e riavviato, come per esempio riscrivere un prompt, o
  inviare un avviso. 

\item[\signald{SIGSTOP}] Il segnale ferma l'esecuzione di un processo, lo porta
  cioè nello stato \textit{stopped} (vedi sez.~\ref{sec:proc_sched}). Il
  segnale non può essere né intercettato, né ignorato, né bloccato.

\item[\signald{SIGTSTP}] Il nome sta per \textit{interactive stop}. Il segnale
  ferma il processo interattivamente, ed è generato dal carattere
  ``\textit{SUSP}'', prodotto dalla combinazione di tasti \cmd{C-z}, ed al
  contrario di \signal{SIGSTOP} può essere intercettato e ignorato. In genere
  un programma installa un gestore per questo segnale quando vuole lasciare il
  sistema o il terminale in uno stato definito prima di fermarsi; se per
  esempio un programma ha disabilitato l'eco sul terminale può installare un
  gestore per riabilitarlo prima di fermarsi.

\item[\signald{SIGTTIN}] Un processo non può leggere dal terminale se esegue
  una sessione di lavoro in \textit{background}. Quando un processo in
  \textit{background} tenta di leggere da un terminale viene inviato questo
  segnale a tutti i processi della sessione di lavoro. L'azione predefinita è
  di fermare il processo.  L'argomento è trattato in
  sez.~\ref{sec:sess_job_control_overview}.

\item[\signald{SIGTTOU}] Segnale analogo al precedente \signal{SIGTTIN}, ma
  generato quando si tenta di scrivere sul terminale o modificarne uno dei
  modi con un processo in \textit{background}. L'azione predefinita è di
  fermare il processo, l'argomento è trattato in
  sez.~\ref{sec:sess_job_control_overview}.
\end{basedescript}


\subsection{I segnali di operazioni errate}
\label{sec:sig_oper_error}

Questi segnali sono usati per riportare al programma errori generati da
operazioni da lui eseguite; non indicano errori del programma quanto errori
che impediscono il completamento dell'esecuzione dovute all'interazione con il
resto del sistema.  L'azione predefinita di questi segnali è normalmente
quella di terminare il processo, questi segnali sono:
\begin{basedescript}{\desclabelwidth{2.0cm}}
\item[\signald{SIGPIPE}] Sta per \textit{Broken pipe}. Se si usano delle
  \textit{pipe}, (o delle FIFO o dei socket) è necessario, prima che un
  processo inizi a scrivere su una di esse, che un altro l'abbia aperta in
  lettura (si veda sez.~\ref{sec:ipc_pipes}). Se il processo in lettura non è
  partito o è terminato inavvertitamente alla scrittura sulla \textit{pipe} il
  kernel genera questo segnale. Se il segnale è bloccato, intercettato o
  ignorato la chiamata che lo ha causato fallisce, restituendo l'errore
  \errcode{EPIPE}.

\item[\signald{SIGXCPU}] Sta per \textit{CPU time limit exceeded}. Questo
  segnale è generato quando un processo eccede il limite impostato per il
  tempo di CPU disponibile, vedi sez.~\ref{sec:sys_resource_limit}. Fino al
  kernel 2.2 terminava semplicemente il processo, a partire dal kernel 2.4,
  seguendo le indicazioni dello standard POSIX.1-2001 viene anche generato un
  \textit{core dump}.

\item[\signald{SIGXFSZ}] Sta per \textit{File size limit exceeded}. Questo
  segnale è generato quando un processo tenta di estendere un file oltre le
  dimensioni specificate dal limite impostato per le dimensioni massime di un
  file, vedi sez.~\ref{sec:sys_resource_limit}.  Fino al kernel 2.2 terminava
  semplicemente il processo, a partire dal kernel 2.4, seguendo le indicazioni
  dello standard POSIX.1-2001 viene anche generato un \textit{core dump}.

\item[\signald{SIGLOST}] Sta per \textit{Resource lost}. Tradizionalmente è il
  segnale che viene generato quando si perde un advisory lock su un file su
  NFS perché il server NFS è stato riavviato. Il progetto GNU lo utilizza per
  indicare ad un client il crollo inaspettato di un server. In Linux è
  definito come sinonimo di \signal{SIGIO} e non viene più usato.
\end{basedescript}


\subsection{Ulteriori segnali}
\label{sec:sig_misc_sig}

Raccogliamo qui infine una serie di segnali che hanno scopi differenti non
classificabili in maniera omogenea. Questi segnali sono:
\begin{basedescript}{\desclabelwidth{2.0cm}}
\item[\signald{SIGUSR1}] Insieme a \signal{SIGUSR2} è un segnale a disposizione
  dell'utente che lo può usare per quello che vuole. Viene generato solo
  attraverso l'invocazione della funzione \func{kill}. Entrambi i segnali
  possono essere utili per implementare una comunicazione elementare fra
  processi diversi, o per eseguire a richiesta una operazione utilizzando un
  gestore. L'azione predefinita è di terminare il processo.
\item[\signald{SIGUSR2}] È il secondo segnale a disposizione degli utenti. Per
  il suo utilizzo vale esattamente quanto appena detto per \signal{SIGUSR1}.
\item[\signald{SIGWINCH}] Il nome sta per \textit{window (size) change} e viene
  generato in molti sistemi (GNU/Linux compreso) quando le dimensioni (in
  righe e colonne) di un terminale vengono cambiate. Viene usato da alcuni
  programmi testuali per riformattare l'uscita su schermo quando si cambia
  dimensione a quest'ultimo. L'azione predefinita è di essere ignorato.
\item[\signald{SIGINFO}] Il segnale indica una richiesta di informazioni. È
  usato con il controllo di sessione, causa la stampa di informazioni da parte
  del processo leader del gruppo associato al terminale di controllo, gli
  altri processi lo ignorano. Su Linux però viene utilizzato come sinonimo di
  \signal{SIGPWR} e l'azione predefinita è di terminare il processo.
\item[\signald{SIGPWR}] Il segnale indica un cambio nello stato di
  alimentazione di un eventuale gruppo di continuità e viene usato
  principalmente per segnalare l'assenza ed il ritorno della corrente. Viene
  usato principalmente con \cmd{init} per attivare o fermare le procedure di
  spegnimento automatico all'esaurimento delle batterie. L'azione predefinita
  è di terminare il processo.
\item[\signald{SIGSTKFLT}] Indica un errore nello stack del coprocessore
  matematico, è definito solo per le architetture PC, ma è completamente
  inusato. L'azione predefinita è di terminare il processo.
\end{basedescript}


\subsection{Le funzioni \func{strsignal} e \func{psignal}}
\label{sec:sig_strsignal}

Per la descrizione dei segnali il sistema mette a disposizione due funzioni
che stampano un messaggio di descrizione specificando il numero del segnale
con una delle costanti di tab.~\ref{tab:sig_signal_list}.  In genere si usano
quando si vuole notificare all'utente il segnale ricevuto, ad esempio nel caso
di terminazione di un processo figlio o di un gestore che gestisce più
segnali.

La prima funzione, \funcd{strsignal}, è una estensione GNU fornita dalla
\acr{glibc}, ed è accessibile solo avendo definito la macro
\macro{\_GNU\_SOURCE}, il suo comportamento è analogo a quello della funzione
\func{strerror} (si veda sez.~\ref{sec:sys_strerror}) usata per notificare gli
errori:

\begin{funcproto}{
\fhead{string.h}
\fdecl{char *strsignal(int signum)} 
\fdesc{Ottiene la descrizione di un segnale.} 
}

{La funzione ritorna puntatore ad una stringa che descrive il segnale, non
  sono previste condizioni di errore ed \var{errno} non viene modificata.}
\end{funcproto}


La funzione ritorna sempre il puntatore ad una stringa che contiene la
descrizione del segnale indicato dall'argomento \param{signum}, se questo non
indica un segnale valido viene restituito il puntatore ad una stringa che
segnale che il valore indicato non è valido.  Dato che la stringa è allocata
staticamente non se ne deve modificare il contenuto, che resta valido solo
fino alla successiva chiamata di \func{strsignal}. Nel caso si debba mantenere
traccia del messaggio sarà necessario copiarlo.

La seconda funzione, \funcd{psignal}, deriva da BSD ed è analoga alla funzione
\func{perror} descritta in sez.~\ref{sec:sys_strerror}, il suo prototipo è:

\begin{funcproto}{
\fhead{signal.h}
\fdecl{void psignal(int sig, const char *s)}
\fdesc{Stampa un messaggio di descrizione di un segnale.} 
}
{La funzione non ritorna nulla e non prevede errori.}  
\end{funcproto}

La funzione stampa sullo \textit{standard error} un messaggio costituito dalla
stringa passata nell'argomento \param{s}, seguita dal carattere di due punti
ed una descrizione del segnale indicato dall'argomento \param{sig}. 

Una modalità alternativa per utilizzare le descrizioni restituite da
\func{strsignal} e \func{psignal} è quello di usare la variabile globale
\var{sys\_siglist}, che è definita in \headfile{signal.h} e può essere
acceduta con la dichiarazione:
\includecodesnip{listati/siglist.c}

L'array \var{sys\_siglist} contiene i puntatori alle stringhe di descrizione,
indicizzate per numero di segnale, per cui una chiamata del tipo di \code{char
  *decr = strsignal(SIGINT)} può essere sostituita dall'equivalente \code{char
  *decr = sys\_siglist[SIGINT]}.



\section{La gestione di base dei segnali}
\label{sec:sig_management}

I segnali sono il primo e più classico esempio di eventi asincroni, cioè di
eventi che possono accadere in un qualunque momento durante l'esecuzione di un
programma. Per questa loro caratteristica la loro gestione non può essere
effettuata all'interno del normale flusso di esecuzione dello stesso, ma è
delegata appunto agli eventuali gestori che si sono installati.

In questa sezione vedremo come si effettua la gestione dei segnali, a partire
dalla loro interazione con le \textit{system call}, passando per le varie
funzioni che permettono di installare i gestori e controllare le reazioni di
un processo alla loro occorrenza.


\subsection{Il comportamento generale del sistema}
\label{sec:sig_gen_beha}

Abbiamo già trattato in sez.~\ref{sec:sig_intro} le modalità con cui il
sistema gestisce l'interazione fra segnali e processi, ci resta da esaminare
però il comportamento delle \textit{system call}; in particolare due di esse,
\func{fork} ed \func{exec}, dovranno essere prese esplicitamente in
considerazione, data la loro stretta relazione con la creazione di nuovi
processi.

Come accennato in sez.~\ref{sec:proc_fork} quando viene creato un nuovo
processo esso eredita dal padre sia le azioni che sono state impostate per i
singoli segnali, che la maschera dei segnali bloccati (vedi
sez.~\ref{sec:sig_sigmask}).  Invece tutti i segnali pendenti e gli allarmi
vengono cancellati; essi infatti devono essere recapitati solo al padre, al
figlio dovranno arrivare solo i segnali dovuti alle sue azioni.

Quando si mette in esecuzione un nuovo programma con \func{exec} (si ricordi
quanto detto in sez.~\ref{sec:proc_exec}) tutti i segnali per i quali è stato
installato un gestore vengono reimpostati a \constd{SIG\_DFL}. Non ha più
senso infatti fare riferimento a funzioni definite nel programma originario,
che non sono presenti nello spazio di indirizzi del nuovo programma.

Si noti che questo vale solo per le azioni per le quali è stato installato un
gestore, viene mantenuto invece ogni eventuale impostazione dell'azione a
\constd{SIG\_IGN}. Questo permette ad esempio alla shell di impostare ad
\const{SIG\_IGN} le risposte per \signal{SIGINT} e \signal{SIGQUIT} per i
programmi eseguiti in background, che altrimenti sarebbero interrotti da una
successiva pressione di \texttt{C-c} o \texttt{C-y}.

Per quanto riguarda il comportamento di tutte le altre \textit{system call} si
danno sostanzialmente due casi, a seconda che esse siano \textsl{lente}
(\textit{slow}) o \textsl{veloci} (\textit{fast}). La gran parte di esse
appartiene a quest'ultima categoria, che non è influenzata dall'arrivo di un
segnale. Esse sono dette \textsl{veloci} in quanto la loro esecuzione è
sostanzialmente immediata. La risposta al segnale viene sempre data dopo che
la \textit{system call} è stata completata, in quanto attendere per eseguire
un gestore non comporta nessun inconveniente.

\index{system~call~lente|(}

In alcuni casi però alcune \textit{system call} possono bloccarsi
indefinitamente e per questo motivo vengono chiamate \textsl{lente} o
\textsl{bloccanti}. In questo caso non si può attendere la conclusione della
\textit{system call}, perché questo renderebbe impossibile una risposta pronta
al segnale, per cui il gestore viene eseguito prima che la \textit{system
  call} sia ritornata.  Un elenco dei casi in cui si presenta questa
situazione è il seguente:
\begin{itemize*}
\item la lettura da file che possono bloccarsi in attesa di dati non ancora
  presenti (come per certi file di dispositivo, i socket o le \textit{pipe});
\item la scrittura sugli stessi file, nel caso in cui dati non possano essere
  accettati immediatamente (di nuovo comune per i socket);
\item l'apertura di un file di dispositivo che richiede operazioni non
  immediate per una risposta (ad esempio l'apertura di un nastro che deve
  essere riavvolto);
\item le operazioni eseguite con \func{ioctl} che non è detto possano essere
  eseguite immediatamente;
\item l'uso di funzioni di intercomunicazione fra processi (vedi
  cap.~\ref{cha:IPC}) che si bloccano in attesa di risposte da altri processi;
\item l'uso della funzione \func{pause} (vedi sez.~\ref{sec:sig_pause_sleep})
  e le analoghe \func{sigsuspend}, \func{sigtimedwait}, e \func{sigwaitinfo}
  (vedi sez.~\ref{sec:sig_real_time}), usate appunto per attendere l'arrivo di
  un segnale;
\item l'uso delle funzioni associate al \textit{file locking} (vedi
  sez.~\ref{sec:file_locking})
\item l'uso della funzione \func{wait} e le analoghe funzioni di attesa se
  nessun processo figlio è ancora terminato.
\end{itemize*}

In questo caso si pone il problema di cosa fare una volta che il gestore sia
ritornato. La scelta originaria dei primi Unix era quella di far ritornare
anche la \textit{system call} restituendo l'errore di \errcode{EINTR}. Questa
è a tutt'oggi una scelta corrente, ma comporta che i programmi che usano dei
gestori controllino lo stato di uscita delle funzioni che eseguono una system
call lenta per ripeterne la chiamata qualora l'errore fosse questo.

Dimenticarsi di richiamare una \textit{system call} interrotta da un segnale è
un errore comune, tanto che la \acr{glibc} provvede una macro
\code{TEMP\_FAILURE\_RETRY(expr)} che esegue l'operazione automaticamente,
ripetendo l'esecuzione dell'espressione \var{expr} fintanto che il risultato
non è diverso dall'uscita con un errore \errcode{EINTR}.

La soluzione è comunque poco elegante e BSD ha scelto un approccio molto
diverso, che è quello di fare ripartire automaticamente una \textit{system
  call} interrotta invece di farla fallire. In questo caso ovviamente non c'è
bisogno di preoccuparsi di controllare il codice di errore; si perde però la
possibilità di eseguire azioni specifiche all'occorrenza di questa particolare
condizione.

Linux e la \acr{glibc} consentono di utilizzare entrambi gli approcci,
attraverso una opportuna opzione di \func{sigaction} (vedi
sez.~\ref{sec:sig_sigaction}). È da chiarire comunque che nel caso di
interruzione nel mezzo di un trasferimento parziale di dati, le \textit{system
  call} ritornano sempre indicando i byte trasferiti.

Si tenga presente però che alcune \textit{system call} vengono comunque
interrotte con un errore di \errcode{EINTR} indipendentemente dal fatto che ne
possa essere stato richiesto il riavvio automatico, queste funzioni sono:

\begin{itemize*}
\item le funzioni di attesa di un segnale: \func{pause} (vedi
  sez.~\ref{sec:sig_pause_sleep}) o \func{sigsuspend}, \func{sigtimedwait}, e
  \func{sigwaitinfo} (vedi sez.~\ref{sec:sig_real_time}).
\item le funzioni di attesa dell'\textit{I/O multiplexing} (vedi
  sez.~\ref{sec:file_multiplexing}) come \func{select}, \func{pselect},
  \func{poll}, \func{ppoll}, \func{epoll\_wait} e \func{epoll\_pwait}.
\item le funzioni del System V IPC che prevedono attese: \func{msgrcv},
  \func{msgsnd} (vedi sez.~\ref{sec:ipc_sysv_mq}), \func{semop} e
  \func{semtimedop} (vedi sez.~\ref{sec:ipc_sysv_sem}).
\item le funzioni per la messa in attesa di un processo come \func{usleep},
  \func{nanosleep} (vedi sez.~\ref{sec:sig_pause_sleep}) e
  \func{clock\_nanosleep} (vedi sez.~\ref{sec:sig_timer_adv}).
\item le funzioni che operano sui socket quando è stato impostato un
  \textit{timeout} sugli stessi con \func{setsockopt} (vedi
  sez.~\ref{sec:sock_generic_options}) ed in particolare \func{accept},
  \func{recv}, \func{recvfrom}, \func{recvmsg} per un \textit{timeout} in
  ricezione e \func{connect}, \func{send}, \func{sendto} e \func{sendmsg} per
  un \textit{timeout} in trasmissione.
%\item la funzione \func{io\_getevents} per l'I/O asincrono (vedi sez.??)
\end{itemize*}


\index{system~call~lente|)}


\subsection{L'installazione di un gestore}
\label{sec:sig_signal}

L'interfaccia più semplice per la gestione dei segnali è costituita dalla
funzione di sistema \funcd{signal} che è definita fin dallo standard ANSI C.
Quest'ultimo però non considera sistemi multitasking, per cui la definizione è
tanto vaga da essere del tutto inutile in un sistema Unix. Per questo motivo
ogni implementazione successiva ne ha modificato e ridefinito il
comportamento, pur mantenendone immutato il prototipo\footnote{in realtà in
  alcune vecchie implementazioni (SVr4 e 4.3+BSD in particolare) vengono usati
  alcuni argomenti aggiuntivi per definire il comportamento della funzione,
  vedremo in sez.~\ref{sec:sig_sigaction} che questo è possibile usando la
  funzione \func{sigaction}.}  che è:

\begin{funcproto}{
\fhead{signal.h}
\fdecl{sighandler\_t signal(int signum, sighandler\_t handler)}
\fdesc{Installa un gestore di segnale (\textit{signal handler}).} 
}

{La funzione ritorna il precedente gestore in caso di successo in caso di
  successo e \constd{SIG\_ERR} per un errore, nel qual caso \var{errno}
  assumerà il valore:
  \begin{errlist}
  \item[\errcode{EINVAL}] il numero di segnale \param{signum} non è valido.
  \end{errlist}
}  
\end{funcproto}

In questa definizione per l'argomento \param{handler} che indica il gestore da
installare si è usato un tipo di dato, \type{sighandler\_t}, che è una
estensione GNU, definita dalla \acr{glibc}, che permette di riscrivere il
prototipo di \func{signal} nella forma appena vista, molto più leggibile di
quanto non sia la versione originaria, che di norma è definita come:
\includecodesnip{listati/signal.c}
questa infatti, per la poca chiarezza della sintassi del C quando si vanno a
trattare puntatori a funzioni, è molto meno comprensibile.  Da un confronto
con il precedente prototipo si può dedurre la definizione di
\typed{sighandler\_t} che è:
\includecodesnip{listati/sighandler_t.c}
e cioè un puntatore ad una funzione \ctyp{void} (cioè senza valore di ritorno)
e che prende un argomento di tipo \ctyp{int}. Si noti come si devono usare le
parentesi intorno al nome della funzione per via delle precedenze degli
operatori del C, senza di esse si sarebbe definita una funzione che ritorna un
puntatore a \ctyp{void} e non un puntatore ad una funzione \ctyp{void}.

La funzione \func{signal} quindi restituisce e prende come secondo argomento
un puntatore a una funzione di questo tipo, che è appunto la funzione che
verrà usata come gestore del segnale.  Il numero di segnale passato
nell'argomento \param{signum} può essere indicato direttamente con una delle
costanti definite in sez.~\ref{sec:sig_standard}.  

L'argomento \param{handler} che indica il gestore invece, oltre all'indirizzo
della funzione da chiamare all'occorrenza del segnale, può assumere anche i
due valori costanti \const{SIG\_IGN} e \const{SIG\_DFL}. Il primo indica che
il segnale deve essere ignorato. Il secondo ripristina l'azione predefinita, e
serve a tornare al comportamento di default quando non si intende più gestire
direttamente un segnale. Si ricordi però che i due segnali \signal{SIGKILL} e
\signal{SIGSTOP} non possono essere né ignorati né intercettati e per loro
l'uso di \func{signal} non ha alcun effetto, qualunque cosa si specifichi
per \param{handler}.

La funzione restituisce l'indirizzo dell'azione precedente, che può essere
salvato per poterlo ripristinare (con un'altra chiamata a \func{signal}) in un
secondo tempo. Si ricordi che se si imposta come azione \const{SIG\_IGN} o si
imposta \const{SIG\_DFL} per un segnale la cui azione predefinita è di essere
ignorato, tutti i segnali pendenti saranno scartati, e non verranno mai
notificati.

L'uso di \func{signal} è soggetto a problemi di compatibilità, dato che essa
si comporta in maniera diversa per sistemi derivati da BSD o da System V. In
questi ultimi infatti la funzione è conforme al comportamento originale dei
primi Unix in cui il gestore viene disinstallato alla sua chiamata, secondo la
semantica inaffidabile; anche Linux seguiva questa convenzione con le vecchie
librerie del C come la \acr{libc4} e la \acr{libc5}.\footnote{nelle
  \acr{libc5} esiste però la possibilità di includere \file{bsd/signal.h} al
  posto di \headfile{signal.h}, nel qual caso la funzione \func{signal} viene
  ridefinita per seguire la semantica affidabile usata da BSD.}

Al contrario BSD segue la semantica affidabile, non disinstallando il gestore
e bloccando il segnale durante l'esecuzione dello stesso. Con l'utilizzo della
\acr{glibc} dalla versione 2 anche Linux è passato a questo comportamento.  Il
comportamento della versione originale della funzione, il cui uso è deprecato
per i motivi visti in sez.~\ref{sec:sig_semantics}, può essere ottenuto
chiamando \funcm{sysv\_signal}, una volta che si sia definita la macro
\macro{\_XOPEN\_SOURCE}.  In generale, per evitare questi problemi, l'uso di
\func{signal}, che tra l'altro ha un comportamento indefinito in caso di
processo multi-\textit{thread}, è da evitare: tutti i nuovi programmi devono
usare \func{sigaction}.

È da tenere presente che, seguendo lo standard POSIX, il comportamento di un
processo che ignora i segnali \signal{SIGFPE}, \signal{SIGILL}, o
\signal{SIGSEGV}, qualora questi non originino da una chiamata ad una
\func{kill} o altra funzione affine, è indefinito. Un gestore che ritorna da
questi segnali può dare luogo ad un ciclo infinito.


\subsection{Le funzioni per l'invio di segnali}
\label{sec:sig_kill_raise}

Come accennato in sez.~\ref{sec:sig_types} un segnale può anche essere
generato direttamente nell'esecuzione di un programma, attraverso la chiamata
ad una opportuna \textit{system call}. Le funzioni che si utilizzano di solito
per inviare un segnale generico ad un processo sono \func{raise} e
\func{kill}.

La funzione \funcd{raise}, definita dallo standard ANSI C, serve per inviare
un segnale al processo corrente,\footnote{non prevedendo la presenza di un
  sistema multiutente lo standard ANSI C non poteva che definire una funzione
  che invia il segnale al programma in esecuzione, nel caso di Linux questa
  viene implementata come funzione di compatibilità.}  il suo prototipo è:

\begin{funcproto}{
\fhead{signal.h}
\fdecl{int raise(int sig)}
\fdesc{Invia un segnale al processo corrente.} 
}

{La funzione ritorna $0$ in caso di successo e $-1$ per un errore, nel qual
  caso \var{errno} assumerà il valore: 
  \begin{errlist}
  \item[\errcode{EINVAL}] il segnale \param{sig} non è valido.
  \end{errlist}
}
\end{funcproto}

Il valore di \param{sig} specifica il segnale che si vuole inviare e può
essere specificato con una delle costanti illustrate in
tab.~\ref{tab:sig_signal_list}.  In genere questa funzione viene usata per
riprodurre il comportamento predefinito di un segnale che sia stato
intercettato. In questo caso, una volta eseguite le operazioni volute, il
gestore dovrà prima reinstallare l'azione predefinita, per poi attivarla
chiamando \func{raise}.

In realtà \func{raise} è una funzione di libreria, che per i processi ordinari
viene implementata attraverso la funzione di sistema \funcd{kill} che è quella
che consente effettivamente di inviare un segnale generico ad un processo, il
 suo prototipo è:

\begin{funcproto}{
\fhead{sys/types.h}
\fhead{signal.h}
\fdecl{int kill(pid\_t pid, int sig)}
\fdesc{Invia un segnale ad uno o più processi.} 
}

{La funzione ritorna $0$ in caso di successo e $-1$ per un errore, nel qual
  caso \var{errno} assumerà uno dei valori: 
  \begin{errlist}
    \item[\errcode{EINVAL}] il segnale specificato non esiste.
    \item[\errcode{EPERM}] non si hanno privilegi sufficienti ad inviare il
      segnale.
    \item[\errcode{ESRCH}] il processo o il gruppo di processi indicato non
      esiste.
  \end{errlist}
}
\end{funcproto}

La funzione invia il segnale specificato dall'argomento \param{sig} al
processo o ai processi specificati con l'argomento \param{pid}.  Lo standard
POSIX prevede che il valore 0 per \param{sig} sia usato per specificare il
segnale nullo.  Se la funzione viene chiamata con questo valore non viene
inviato nessun segnale, ma viene eseguito il controllo degli errori, in tal
caso si otterrà un errore \errcode{EPERM} se non si hanno i permessi necessari
ed un errore \errcode{ESRCH} se il processo o i processi specificati
con \param{pid} non esistono.

\begin{table}[htb]
  \footnotesize
  \centering
  \begin{tabular}[c]{|r|p{8cm}|}
    \hline
    \textbf{Valore} & \textbf{Significato} \\
    \hline
    \hline
    $>0$ & Il segnale è mandato al processo con \ids{PID} uguale
           a \param{pid}.\\
    0    & Il segnale è mandato ad ogni processo del \textit{process group}
           del chiamante.\\  
    $-1$ & Il segnale è mandato ad ogni processo (eccetto \cmd{init}).\\
    $<-1$& Il segnale è mandato ad ogni processo del \textit{process group} 
           con \ids{PGID} uguale a $|\param{pid}|$.\\
    \hline
  \end{tabular}
  \caption{Valori dell'argomento \param{pid} per la funzione
    \func{kill}.}
  \label{tab:sig_kill_values}
\end{table}

A seconda del valore dell'argomento \param{pid} si può inviare il segnale ad
uno specifico processo, ad un \textit{process group} (vedi
sez.~\ref{sec:sess_proc_group}) o a tutti i processi, secondo quanto
illustrato in tab.~\ref{tab:sig_kill_values} che riporta i valori possibili
per questo argomento. Si tenga conto però che il sistema ricicla i \ids{PID}
(come accennato in sez.~\ref{sec:proc_pid}) per cui l'esistenza di un processo
non significa che esso sia realmente quello a cui si intendeva mandare il
segnale.

Indipendentemente dalla funzione specifica che viene usata solo
l'amministratore può inviare un segnale ad un processo qualunque, in tutti gli
altri casi l'\ids{UID} reale o l'\ids{UID} effettivo del processo chiamante
devono corrispondere all'\ids{UID} reale o all'\ids{UID} salvato della
destinazione. Fa eccezione il caso in cui il segnale inviato sia
\signal{SIGCONT}, nel quale occorre anche che entrambi i processi appartengano
alla stessa sessione.

Si tenga presente che, per il ruolo fondamentale che riveste nel sistema, non
è possibile inviare al processo 1 (cioè a \cmd{init}) segnali per i quali esso
non abbia un gestore installato.  Infine, seguendo le specifiche POSIX
1003.1-2001, l'uso della chiamata \code{kill(-1, sig)} comporta che il segnale
sia inviato (con la solita eccezione di \cmd{init}) a tutti i processi per i
quali i permessi lo consentano. Lo standard permette comunque alle varie
implementazioni di escludere alcuni processi specifici: nel caso in questione
Linux non invia il segnale al processo che ha effettuato la chiamata.

Si noti pertanto che la funzione \code{raise(sig)} può essere definita in
termini di \func{kill}, ed è sostanzialmente equivalente ad una
\code{kill(getpid(), sig)}. Siccome \func{raise}, che è definita nello
standard ISO C, non esiste in alcune vecchie versioni di Unix, in generale
l'uso di \func{kill} finisce per essere più portabile.  Una seconda funzione
che può essere definita in termini di \func{kill} è \funcd{killpg}, il suo
prototipo è:

\begin{funcproto}{
\fhead{signal.h}
\fdecl{int killpg(pid\_t pidgrp, int signal)}
\fdesc{Invia un segnale ad un \textit{process group}.} 
}

{ La funzione ritorna $0$ in caso di successo e $-1$ per un errore, e gli
  errori sono gli stessi di \func{kill}.
}
\end{funcproto}


La funzione invia il segnale \param{signal} al \textit{process group} il cui
\acr{PGID} (vedi sez.~\ref{sec:sess_proc_group}) è indicato
dall'argomento \param{pidgrp}, che deve essere un intero positivo. Il suo
utilizzo è sostanzialmente equivalente all'esecuzione di \code{kill(-pidgrp,
  signal)}.

Oltre alle precedenti funzioni di base, vedremo più avanti che esistono altre
funzioni per inviare segnali generici, come \func{sigqueue} per i segnali
\textit{real-time} (vedi sez.~\ref{sec:sig_real_time}) e le specifiche
funzioni per i \textit{thread} che tratteremo in sez.~\ref{sec:thread_signal}.

Esiste però un'ultima funzione che permette l'invio diretto di un segnale che
vale la pena di trattare a parte per le sue peculiarità. La funzione in
questione è \funcd{abort} che, come accennato in
sez.~\ref{sec:proc_termination}, permette di abortire l'esecuzione di un
programma tramite l'invio del segnale \signal{SIGABRT}. Il suo prototipo è:

\begin{funcproto}{
\fhead{stdlib.h}
\fdecl{void abort(void)}
\fdesc{Abortisce il processo corrente.} 
}

{La funzione non ritorna, il processo viene terminato.}
\end{funcproto}

La differenza fra questa funzione e l'uso di \func{raise} o di un'altra
funzione per l'invio di \signal{SIGABRT} è che anche se il segnale è bloccato
o ignorato, la funzione ha effetto lo stesso. Il segnale può però essere
intercettato per effettuare eventuali operazioni di chiusura prima della
terminazione del processo.

Lo standard ANSI C richiede inoltre che anche se il gestore ritorna, la
funzione non ritorni comunque. Lo standard POSIX.1 va oltre e richiede che se
il processo non viene terminato direttamente dal gestore sia la stessa
\func{abort} a farlo al ritorno dello stesso. Inoltre, sempre seguendo lo
standard POSIX, prima della terminazione tutti i file aperti e gli stream
saranno chiusi ed i buffer scaricati su disco. Non verranno invece eseguite le
eventuali funzioni registrate con \func{atexit} e \func{on\_exit}.

% TODO trattare pidfd_send_signal, aggiunta con il kernel 5.1 (vedi
% https://lwn.net/Articles/783052/) per mandare segnali a processi senza dover
% usare un PID, vedi anche https://lwn.net/Articles/773459/,
% https://git.kernel.org/linus/3eb39f47934f 


\subsection{Le funzioni di allarme ed i \textit{timer}}
\label{sec:sig_alarm_abort}

Un caso particolare di segnali generati a richiesta è quello che riguarda i
vari segnali usati per la temporizzazione, per ciascuno di essi infatti sono
previste delle funzioni specifiche che ne effettuino l'invio. La più comune, e
la più semplice, delle funzioni usate per la temporizzazione è la funzione di
sistema \funcd{alarm}, il cui prototipo è:

\begin{funcproto}{
\fhead{unistd.h}
\fdecl{unsigned int alarm(unsigned int seconds)}
\fdesc{Predispone l'invio di un allarme.} 
}

{La funzione ritorna il numero di secondi rimanenti ad un precedente allarme,
  o $0$ se non c'erano allarmi pendenti, non sono previste condizioni di
  errore.}
\end{funcproto}

La funzione fornisce un meccanismo che consente ad un processo di predisporre
un'interruzione nel futuro, ad esempio per effettuare una qualche operazione
dopo un certo periodo di tempo, programmando l'emissione di un segnale (nel
caso in questione \signal{SIGALRM}) dopo il numero di secondi specificato
dall'argomento \param{seconds}. Se si specifica per \param{seconds} un valore
nullo non verrà inviato nessun segnale. Siccome alla chiamata viene cancellato
ogni precedente allarme, questo valore può essere usato per cancellare una
programmazione precedente.

La funzione inoltre ritorna il numero di secondi rimanenti all'invio
dell'allarme programmato in precedenza. In questo modo è possibile controllare
se non si è cancellato un precedente allarme e predisporre eventuali misure
che permettano di gestire il caso in cui servono più interruzioni.

In sez.~\ref{sec:sys_unix_time} abbiamo visto che ad ogni processo sono
associati tre tempi diversi: il \textit{clock time}, l'\textit{user time} ed
il \textit{system time}.  Per poterli calcolare il kernel mantiene per ciascun
processo tre diversi timer:
\begin{itemize*}
\item un \textit{real-time timer} che calcola il tempo reale trascorso (che
  corrisponde al \textit{clock time}). La scadenza di questo timer provoca
  l'emissione di \signal{SIGALRM};
\item un \textit{virtual timer} che calcola il tempo di processore usato dal
  processo in \textit{user space} (che corrisponde all'\textit{user time}). La
  scadenza di questo timer provoca l'emissione di \signal{SIGVTALRM};
\item un \textit{profiling timer} che calcola la somma dei tempi di processore
  utilizzati direttamente dal processo in \textit{user space}, e dal kernel
  nelle \textit{system call} ad esso relative (che corrisponde a quello che in
  sez.~\ref{sec:sys_unix_time} abbiamo chiamato \textit{processor time}). La
  scadenza di questo timer provoca l'emissione di \signal{SIGPROF}.
\end{itemize*}

Il timer usato da \func{alarm} è il \textit{clock time}, e corrisponde cioè al
tempo reale. La funzione come abbiamo visto è molto semplice, ma proprio per
questo presenta numerosi limiti: non consente di usare gli altri timer, non
può specificare intervalli di tempo con precisione maggiore del secondo e
genera il segnale una sola volta.

Per ovviare a questi limiti Linux deriva da BSD la funzione \funcd{setitimer}
che permette di usare un timer qualunque e l'invio di segnali periodici, al
costo però di una maggiore complessità d'uso e di una minore portabilità. Il
suo prototipo è:

\begin{funcproto}{
\fhead{sys/time.h}
\fdecl{int setitimer(int which, const struct itimerval *value, struct
  itimerval *ovalue)}
  
\fdesc{Predispone l'invio di un segnale di allarme.} 
}

{La funzione ritorna $0$ in caso di successo e $-1$ per un errore, nel qual
  caso \var{errno} assumerà uno dei valori \errval{EINVAL} o \errval{EFAULT}
  nel loro significato generico.}
\end{funcproto}


La funzione predispone l'invio di un segnale di allarme alla scadenza
dell'intervallo indicato dall'argomento \param{value}.  Il valore
dell'argomento \param{which} permette di specificare quale dei tre timer
illustrati in precedenza usare; i possibili valori sono riportati in
tab.~\ref{tab:sig_setitimer_values}.
\begin{table}[htb]
  \footnotesize
  \centering
  \begin{tabular}[c]{|l|l|}
    \hline
    \textbf{Valore} & \textbf{Timer} \\
    \hline
    \hline
    \constd{ITIMER\_REAL}    & \textit{real-time timer}\\
    \constd{ITIMER\_VIRTUAL} & \textit{virtual timer}\\
    \constd{ITIMER\_PROF}    & \textit{profiling timer}\\
    \hline
  \end{tabular}
  \caption{Valori dell'argomento \param{which} per la funzione
    \func{setitimer}.}
  \label{tab:sig_setitimer_values}
\end{table}

Il valore della struttura specificata \param{value} viene usato per impostare
il timer, se il puntatore \param{ovalue} non è nullo il precedente valore
viene salvato qui. I valori dei timer devono essere indicati attraverso una
struttura \struct{itimerval}, definita in fig.~\ref{fig:file_stat_struct}.

La struttura è composta da due membri, il primo, \var{it\_interval} definisce
il periodo del timer; il secondo, \var{it\_value} il tempo mancante alla
scadenza. Entrambi esprimono i tempi tramite una struttura \struct{timeval} che
permette una precisione fino al microsecondo.

Ciascun timer decrementa il valore di \var{it\_value} fino a zero, poi invia
il segnale e reimposta \var{it\_value} al valore di \var{it\_interval}, in
questo modo il ciclo verrà ripetuto; se invece il valore di \var{it\_interval}
è nullo il timer si ferma.

\begin{figure}[!htb]
  \footnotesize \centering
  \begin{minipage}[c]{0.8\textwidth}
    \includestruct{listati/itimerval.h}
  \end{minipage} 
  \normalsize 
  \caption{La struttura \structd{itimerval}, che definisce i valori dei timer
    di sistema.}
  \label{fig:sig_itimerval}
\end{figure}

L'uso di \func{setitimer} consente dunque un controllo completo di tutte le
caratteristiche dei timer, ed in effetti la stessa \func{alarm}, benché
definita direttamente nello standard POSIX.1, può a sua volta essere espressa
in termini di \func{setitimer}, come evidenziato dal manuale della \acr{glibc}
\cite{GlibcMan} che ne riporta la definizione mostrata in
fig.~\ref{fig:sig_alarm_def}.\footnote{questo comporta anche che non è il caso
  di mescolare chiamate ad \func{abort} e a \func{setitimer}.}

\begin{figure}[!htb]
  \footnotesize \centering
  \begin{minipage}[c]{0.8\textwidth}
    \includestruct{listati/alarm_def.c}
  \end{minipage} 
  \normalsize 
  \caption{Definizione di \func{alarm} in termini di \func{setitimer}.} 
  \label{fig:sig_alarm_def}
\end{figure}

Si deve comunque tenere presente che fino al kernel 2.6.16 la precisione di
queste funzioni era limitata dalla frequenza del timer di sistema, determinato
dal valore della costante \texttt{HZ} di cui abbiamo già parlato in
sez.~\ref{sec:proc_hierarchy}, in quanto le temporizzazioni erano calcolate in
numero di interruzioni del timer (i cosiddetti ``\textit{jiffies}''), ed era
assicurato soltanto che il segnale non sarebbe stato mai generato prima della
scadenza programmata (l'arrotondamento cioè era effettuato per
eccesso).\footnote{questo in realtà non è del tutto vero a causa di un bug,
  presente fino al kernel 2.6.12, che in certe circostanze causava l'emissione
  del segnale con un arrotondamento per difetto.}

L'uso del contatore dei \textit{jiffies}, un intero a 32 bit nella maggior
parte dei casi, comportava inoltre l'impossibilità di specificare tempi molto
lunghi. superiori al valore della costante \constd{MAX\_SEC\_IN\_JIFFIES},
pari, nel caso di default di un valore di \const{HZ} di 250, a circa 99 giorni
e mezzo. Con il cambiamento della rappresentazione effettuato nel kernel
2.6.16 questo problema è scomparso e con l'introduzione dei timer ad alta
risoluzione (vedi sez.~\ref{sec:sig_timer_adv}) nel kernel 2.6.21 la
precisione è diventata quella fornita dall'hardware disponibile.

Una seconda causa di potenziali ritardi è che il segnale viene generato alla
scadenza del timer, ma poi deve essere consegnato al processo; se quest'ultimo
è attivo (questo è sempre vero per \const{ITIMER\_VIRTUAL}) la consegna è
immediata, altrimenti può esserci un ulteriore ritardo che può variare a
seconda del carico del sistema.

Questo ha una conseguenza che può indurre ad errori molto subdoli, si tenga
conto poi che in caso di sistema molto carico, si può avere il caso patologico
in cui un timer scade prima che il segnale di una precedente scadenza sia
stato consegnato. In questo caso, per il comportamento dei segnali descritto
in sez.~\ref{sec:sig_sigchld}, un solo segnale sarà consegnato. Per questo
oggi l'uso di questa funzione è deprecato a favore degli
\textit{high-resolution timer} e della cosiddetta \textit{POSIX Timer API},
che tratteremo in sez.~\ref{sec:sig_timer_adv}.

Dato che sia \func{alarm} che \func{setitimer} non consentono di leggere il
valore corrente di un timer senza modificarlo, è possibile usare la funzione
\funcd{getitimer}, il cui prototipo è:

\begin{funcproto}{
\fhead{sys/time.h}
\fdecl{int getitimer(int which, struct itimerval *value)}
\fdesc{Legge il valore di un timer.} 
}

{ La funzione ritorna $0$ in caso di successo e $-1$ per un errore, nel qual
  caso \var{errno} assumerà gli stessi valori di \func{getitimer}.  }
\end{funcproto}

La funzione legge nella struttura \struct{itimerval} puntata da \param{value}
il valore del timer specificato da \param{which} ed i suoi argomenti hanno lo
stesso significato e formato di quelli di \func{setitimer}.


\subsection{Le funzioni di pausa e attesa}
\label{sec:sig_pause_sleep}

Sono parecchie le occasioni in cui si può avere necessità di sospendere
temporaneamente l'esecuzione di un processo. Nei sistemi più elementari in
genere questo veniva fatto con un ciclo di attesa in cui il programma ripete
una operazione un numero sufficiente di volte per far passare il tempo
richiesto.

Ma in un sistema multitasking un ciclo di attesa è solo un inutile spreco di
tempo di processore, dato che altri programmi possono essere eseguiti nel
frattempo, per questo ci sono delle apposite funzioni che permettono di
mantenere un processo in attesa per il tempo voluto, senza impegnare il
processore. In pratica si tratta di funzioni che permettono di portare
esplicitamente il processo nello stato di \textit{sleep} (si ricordi quanto
illustrato in tab.~\ref{tab:proc_proc_states}) per un certo periodo di tempo.

La prima di queste è la funzione di sistema \funcd{pause}, che viene usata per
mettere un processo in attesa per un periodo di tempo indefinito, fino
all'arrivo di un segnale, il suo prototipo è:

\begin{funcproto}{
\fhead{unistd.h}
\fdecl{int pause(void)}
\fdesc{Pone il processo in pausa fino al ricevimento di un segnale.} 
}

{La funzione ritorna solo dopo che un segnale è stato ricevuto ed il relativo
  gestore è ritornato, nel qual caso restituisce $-1$ e \var{errno} assume il
  valore \errval{EINTR}.}
\end{funcproto}

La funzione ritorna sempre con una condizione di errore, dato che il successo
sarebbe quello di continuare ad aspettare indefinitamente. In genere si usa
questa funzione quando si vuole mettere un processo in attesa di un qualche
evento specifico che non è sotto il suo diretto controllo, ad esempio la si
può usare per interrompere l'esecuzione del processo fino all'arrivo di un
segnale inviato da un altro processo.

Quando invece si vuole fare attendere un processo per un intervallo di tempo
già noto in partenza, lo standard POSIX.1 prevede una funzione di attesa
specifica, \funcd{sleep}, il cui prototipo è:

\begin{funcproto}{

\fhead{unistd.h}
\fdecl{unsigned int sleep(unsigned int seconds)}
\fdesc{Pone il processo in pausa per un tempo in secondi.} 
}

{La funzione ritorna $0$ se l'attesa viene completata  o il
  numero di secondi restanti se viene interrotta da un segnale, non sono
  previsti codici di errore.}
\end{funcproto}

La funzione pone il processo in stato di \textit{sleep} per il numero di
secondi specificato dall'argomento \param{seconds}, a meno di non essere
interrotta da un segnale. Alla terminazione del periodo di tempo indicato la
funzione ritorna riportando il processo in stato \textit{runnable} così che
questo possa riprendere l'esecuzione.

In caso di interruzione della funzione non è una buona idea ripetere la
chiamata per il tempo rimanente restituito dalla stessa, in quanto la
riattivazione del processo può avvenire in un qualunque momento, ma il valore
restituito sarà sempre arrotondato al secondo. Questo può avere la conseguenza
che se la successione dei segnali è particolarmente sfortunata e le differenze
si accumulano, si possono avere ritardi anche di parecchi secondi rispetto a
quanto programmato inizialmente. In genere la scelta più sicura in questo caso
è quella di stabilire un termine per l'attesa, e ricalcolare tutte le volte il
numero di secondi che restano da aspettare.

Si tenga presente che alcune implementazioni l'uso di \func{sleep} può avere
conflitti con quello di \signal{SIGALRM}, dato che la funzione può essere
realizzata con l'uso di \func{pause} e \func{alarm}, in una maniera analoga a
quella dell'esempio che vedremo in sez.~\ref{sec:sig_example}. In tal caso
mescolare chiamate di \func{alarm} e \func{sleep} o modificare l'azione
associata \signal{SIGALRM}, può portare a dei risultati indefiniti. Nel caso
della \acr{glibc} è stata usata una implementazione completamente indipendente
e questi problemi non ci sono, ma un programma portabile non può fare questa
assunzione.

La granularità di \func{sleep} permette di specificare attese soltanto in
secondi, per questo sia sotto BSD4.3 che in SUSv2 è stata definita un'altra
funzione con una precisione teorica del microsecondo. I due standard hanno
delle definizioni diverse, ma la \acr{glibc} segue (secondo la pagina di
manuale almeno dalla versione 2.2.2)  quella di SUSv2 per cui la
funzione \funcd{usleep} (dove la \texttt{u} è intesa come sostituzione di
$\mu$), ha il seguente prototipo:

\begin{funcproto}{
\fhead{unistd.h}
\fdecl{int usleep(unsigned long usec)}
\fdesc{Pone il processo in pausa per un tempo in microsecondi.} 
}

{La funzione ritorna $0$ se l'attesa viene completata e $-1$ per un errore,
  nel qual caso \var{errno} assumerà uno dei valori:
  \begin{errlist}
  \item[\errcode{EINTR}] la funzione è stata interrotta da un segnale.
  \item[\errcode{EINVAL}] si è indicato un valore di \param{usec} maggiore di
    1000000.
  \end{errlist}
}
\end{funcproto}

Anche questa funzione, a seconda delle implementazioni, può presentare
problemi nell'interazione con \func{alarm} e \signal{SIGALRM}, per questo
motivo, pur essendovi citata, nello standard POSIX.1-2001 viene deprecata in
favore della nuova funzione di sistema \funcd{nanosleep}, il cui prototipo è:

\begin{funcproto}{
\fhead{unistd.h}
\fdecl{int nanosleep(const struct timespec *req, struct timespec *rem)}
\fdesc{Pone il processo in pausa per un intervallo di tempo.} 
}

{La funzione ritorna $0$ se l'attesa viene completata e $-1$ per un errore,
  nel qual caso \var{errno} assumerà uno dei valori:
  \begin{errlist}
    \item[\errcode{EINTR}] la funzione è stata interrotta da un segnale.
    \item[\errcode{EINVAL}] si è specificato un numero di secondi negativo o un
      numero di nanosecondi maggiore di 999.999.999.
  \end{errlist}
}
\end{funcproto}

La funzione pone il processo in pausa portandolo nello stato di \textit{sleep}
per il tempo specificato dall'argomento \param{req}, ed in caso di
interruzione restituisce il tempo restante nell'argomento \param{rem}.  Lo
standard richiede che la funzione sia implementata in maniera del tutto
indipendente da \func{alarm}, e nel caso di Linux questo è fatto utilizzando
direttamente il timer del kernel. Lo standard richiede inoltre che la funzione
sia utilizzabile senza interferenze con l'uso di \signal{SIGALRM}. La funzione
prende come argomenti delle strutture di tipo \struct{timespec}, la cui
definizione è riportata in fig.~\ref{fig:sys_timespec_struct}, il che permette
di specificare un tempo con una precisione teorica fino al nanosecondo.

La funzione risolve anche il problema di proseguire l'attesa dopo
l'interruzione dovuta ad un segnale; infatti in tal caso in \param{rem} viene
restituito il tempo rimanente rispetto a quanto richiesto
inizialmente,\footnote{con l'eccezione, valida solo nei kernel della serie
  2.4, in cui, per i processi riavviati dopo essere stati fermati da un
  segnale, il tempo passato in stato \texttt{T} non viene considerato nel
  calcolo della rimanenza.} e basta richiamare la funzione per completare
l'attesa.

Anche qui però occorre tenere presente che i tempi sono arrotondati, per cui
la precisione, per quanto migliore di quella ottenibile con \func{sleep}, è
relativa e in caso di molte interruzioni si può avere una deriva, per questo
esiste la funzione \func{clock\_nanosleep} (vedi sez.~\ref{sec:sig_timer_adv})
che permette di specificare un tempo assoluto anziché un tempo relativo.

Chiaramente, anche se il tempo può essere specificato con risoluzioni fino al
nanosecondo, la precisione di \func{nanosleep} è determinata dalla risoluzione
temporale del timer di sistema. Perciò la funzione attenderà comunque il tempo
specificato, ma prima che il processo possa tornare ad essere eseguito
occorrerà almeno attendere la successiva interruzione del timer di sistema,
cioè un tempo che a seconda dei casi può arrivare fino a 1/\const{HZ}, (sempre
che il sistema sia scarico ed il processa venga immediatamente rimesso in
esecuzione). Per questo motivo il valore restituito in \param{rem} è sempre
arrotondato al multiplo successivo di 1/\const{HZ}. 

Con i kernel della serie 2.4 in realtà era possibile ottenere anche pause più
precise del centesimo di secondo usando politiche di \textit{scheduling}
\textit{real-time} come \const{SCHED\_FIFO} o \const{SCHED\_RR} (vedi
sez.~\ref{sec:proc_real_time}); in tal caso infatti il calcolo sul numero di
interruzioni del timer veniva evitato utilizzando direttamente un ciclo di
attesa con cui si raggiungevano pause fino ai 2~ms con precisioni del
$\mu$s. Questa estensione è stata rimossa con i kernel della serie 2.6, che
consentono una risoluzione più alta del timer di sistema; inoltre a partire
dal kernel 2.6.21, \func{nanosleep} può avvalersi del supporto dei timer ad
alta risoluzione, ottenendo la massima precisione disponibile sull'hardware
della propria macchina.


\subsection{Un esempio elementare}
\label{sec:sig_sigchld}

Un semplice esempio per illustrare il funzionamento di un gestore di segnale è
quello della gestione di \signal{SIGCHLD}. Abbiamo visto in
sez.~\ref{sec:proc_termination} che una delle azioni eseguite dal kernel alla
conclusione di un processo è quella di inviare questo segnale al padre. In
generale dunque, quando non interessa elaborare lo stato di uscita di un
processo, si può completare la gestione della terminazione installando un
gestore per \signal{SIGCHLD} il cui unico compito sia quello di chiamare
\func{waitpid} per completare la procedura di terminazione in modo da evitare
la formazione di \textit{zombie}.\footnote{si ricordi comunque che dal kernel
  2.6 seguendo lo standard POSIX.1-2001 per evitare di dover ricevere gli
  stati di uscita che non interessano basta impostare come azione predefinita
  quella di ignorare \signal{SIGCHLD}, nel qual caso viene assunta la
  semantica di System V, in cui il segnale non viene inviato, il sistema non
  genera \textit{zombie} e lo stato di terminazione viene scartato senza dover
  chiamare una \func{wait}.}

In fig.~\ref{fig:sig_sigchld_handl} è mostrato il codice contenente una
implementazione generica di una funzione di gestione per \signal{SIGCHLD},
(che si trova nei sorgenti allegati nel file \file{SigHand.c}); se ripetiamo i
test di sez.~\ref{sec:proc_termination}, invocando \cmd{forktest} con
l'opzione \cmd{-s} (che si limita ad effettuare l'installazione di questa
funzione come gestore di \signal{SIGCHLD}) potremo verificare che non si ha
più la creazione di \textit{zombie}.

\begin{figure}[!htbp]
  \footnotesize  \centering
  \begin{minipage}[c]{\codesamplewidth}
    \includecodesample{listati/hand_sigchild.c}
  \end{minipage}
  \normalsize 
  \caption{Codice di una funzione generica di gestione per il segnale
    \signal{SIGCHLD}.}
  \label{fig:sig_sigchld_handl}
\end{figure}

Il codice del gestore è di lettura immediata, come buona norma di
programmazione (si ricordi quanto accennato sez.~\ref{sec:sys_errno}) si
comincia (\texttt{\small 6--7}) con il salvare lo stato corrente di
\var{errno}, in modo da poterlo ripristinare prima del ritorno del gestore
(\texttt{\small 16--17}). In questo modo si preserva il valore della variabile
visto dal corso di esecuzione principale del processo, che altrimenti sarebbe
sovrascritto dal valore restituito nella successiva chiamata di
\func{waitpid}.

Il compito principale del gestore è quello di ricevere lo stato di
terminazione del processo, cosa che viene eseguita nel ciclo in
(\texttt{\small 9--15}).  Il ciclo è necessario a causa di una caratteristica
fondamentale della gestione dei segnali: abbiamo già accennato come fra la
generazione di un segnale e l'esecuzione del gestore possa passare un certo
lasso di tempo e niente ci assicura che il gestore venga eseguito prima della
generazione di ulteriori segnali dello stesso tipo. In questo caso normalmente
i segnali successivi vengono ``\textsl{fusi}'' col primo ed al processo ne
viene recapitato soltanto uno.

Questo può essere un caso comune proprio con \signal{SIGCHLD}, qualora capiti
che molti processi figli terminino in rapida successione. Esso inoltre si
presenta tutte le volte che un segnale viene bloccato: per quanti siano i
segnali emessi durante il periodo di blocco, una volta che quest'ultimo sarà
rimosso verrà recapitato un solo segnale.

Allora, nel caso della terminazione dei processi figli, se si chiamasse
\func{waitpid} una sola volta, essa leggerebbe lo stato di terminazione per un
solo processo, anche se i processi terminati sono più di uno, e gli altri
resterebbero in stato di \textit{zombie} per un tempo indefinito.

Per questo occorre ripetere la chiamata di \func{waitpid} fino a che essa non
ritorni un valore nullo, segno che non resta nessun processo di cui si debba
ancora ricevere lo stato di terminazione (si veda sez.~\ref{sec:proc_wait} per
la sintassi della funzione). Si noti anche come la funzione venga invocata con
il parametro \const{WNOHANG} che permette di evitare il suo blocco quando
tutti gli stati di terminazione sono stati ricevuti.



\section{La gestione avanzata dei segnali}
\label{sec:sig_adv_control}

Le funzioni esaminate finora fanno riferimento alle modalità più elementari
della gestione dei segnali; non si sono pertanto ancora prese in
considerazione le tematiche più complesse, collegate alle varie \textit{race
  condition} che i segnali possono generare e alla natura asincrona degli
stessi.

Affronteremo queste problematiche in questa sezione, partendo da un esempio
che le evidenzi, per poi prendere in esame le varie funzioni che permettono di
risolvere i problemi più complessi connessi alla programmazione con i segnali,
fino a trattare le caratteristiche generali della gestione dei medesimi nella
casistica ordinaria.


\subsection{Alcune problematiche aperte}
\label{sec:sig_example}

Come accennato in sez.~\ref{sec:sig_pause_sleep} è possibile implementare
\func{sleep} a partire dall'uso di \func{pause} e \func{alarm}. A prima vista
questo può sembrare di implementazione immediata; ad esempio una semplice
versione di \func{sleep} potrebbe essere quella illustrata in
fig.~\ref{fig:sig_sleep_wrong}.

\begin{figure}[!htb]
  \footnotesize \centering
  \begin{minipage}[c]{\codesamplewidth}
    \includecodesample{listati/sleep_danger.c}
  \end{minipage}
  \normalsize 
  \caption{Una implementazione pericolosa di \func{sleep}.} 
  \label{fig:sig_sleep_wrong}
\end{figure}

Dato che è nostra intenzione utilizzare \signal{SIGALRM} il primo passo della
nostra implementazione sarà quello di installare il relativo gestore salvando
il precedente (\texttt{\small 14--17}).  Si effettuerà poi una chiamata ad
\func{alarm} per specificare il tempo d'attesa per l'invio del segnale a cui
segue la chiamata a \func{pause} per fermare il programma (\texttt{\small
  18--20}) fino alla sua ricezione.  Al ritorno di \func{pause}, causato dal
ritorno del gestore (\texttt{\small 1--9}), si ripristina il gestore originario
(\texttt{\small 21--22}) restituendo l'eventuale tempo rimanente
(\texttt{\small 23--24}) che potrà essere diverso da zero qualora
l'interruzione di \func{pause} venisse causata da un altro segnale.

Questo codice però, a parte il non gestire il caso in cui si è avuta una
precedente chiamata a \func{alarm} (che si è tralasciato per brevità),
presenta una pericolosa \textit{race condition}.  Infatti, se il processo
viene interrotto fra la chiamata di \func{alarm} e \func{pause}, può capitare
(ad esempio se il sistema è molto carico) che il tempo di attesa scada prima
dell'esecuzione di quest'ultima, cosicché essa sarebbe eseguita dopo l'arrivo
di \signal{SIGALRM}. In questo caso ci si troverebbe di fronte ad un
\textit{deadlock}, in quanto \func{pause} non verrebbe mai più interrotta (se
non in caso di un altro segnale).

Questo problema può essere risolto (ed è la modalità con cui veniva fatto in
SVr2) usando la funzione \func{longjmp} (vedi sez.~\ref{sec:proc_longjmp}) per
uscire dal gestore. In questo modo, con una condizione sullo stato di
uscita di quest'ultima, si può evitare la chiamata a \func{pause}, usando un
codice del tipo di quello riportato in fig.~\ref{fig:sig_sleep_incomplete}.

\begin{figure}[!htb]
  \footnotesize \centering
  \begin{minipage}[c]{\codesamplewidth}
    \includecodesample{listati/sleep_defect.c}
  \end{minipage}
  \normalsize 
  \caption{Una implementazione ancora malfunzionante di \func{sleep}.} 
  \label{fig:sig_sleep_incomplete}
\end{figure}

In questo caso il gestore (\texttt{\small 18--27}) non ritorna come in
fig.~\ref{fig:sig_sleep_wrong}, ma usa la funzione \func{longjmp}
(\texttt{\small 25}) per rientrare direttamente nel corpo principale del
programma. Dato che in questo caso il valore di uscita che verrà restituito da
\func{setjmp} è 1, grazie alla condizione impostata in (\texttt{\small 9--12})
si potrà evitare comunque che \func{pause} sia chiamata a vuoto.

Ma anche questa implementazione comporta dei problemi, in questo caso infatti
non viene gestita correttamente l'interazione con gli altri segnali. Se
infatti il segnale di allarme interrompe un altro gestore, l'esecuzione non
riprenderà nel gestore in questione, ma nel ciclo principale, interrompendone
inopportunamente l'esecuzione.  Lo stesso tipo di problemi si presenterebbero
se si volesse usare questa implementazione di \func{alarm} per stabilire un
timeout su una qualunque \textit{system call} bloccante.

Un secondo esempio dei problemi a cui si può andare incontro è quello in cui
si usa un segnale per notificare una qualche forma di evento. In genere quello
che si fa in questo caso è impostare all'interno del gestore un opportuno flag
da controllare nel corpo principale del programma, con un codice del tipo di
quello riportato in fig.~\ref{fig:sig_event_wrong}.

La logica del programma è quella di far impostare al gestore (\texttt{\small
  14--19}) una variabile globale, preventivamente inizializzata nel programma
principale, ad un diverso valore. In questo modo dal corpo principale del
programma si potrà determinare, osservandone il contenuto di detta variabile,
l'occorrenza o meno del segnale, ed eseguire le azioni conseguenti
(\texttt{\small 6--11}) relative.

\begin{figure}[!htbp]
  \footnotesize\centering
  \begin{minipage}[c]{\codesamplewidth}
    \includecodesample{listati/sig_alarm.c}
  \end{minipage}
  \normalsize 
  \caption{Un esempio non funzionante del codice per il controllo di un
    evento generato da un segnale.}
  \label{fig:sig_event_wrong}
\end{figure}

Questo è il tipico esempio di caso, già citato in
sez.~\ref{sec:proc_race_cond}, in cui si genera una \textit{race
  condition}. Infatti, in una situazione in cui un segnale è già arrivato (e
quindi \var{flag} è già stata impostata ad 1 nel gestore) se un altro segnale
arriva immediatamente dopo l'esecuzione del controllo (\texttt{\small 6}) ma
prima della cancellazione di \var{flag} fatta subito dopo (\texttt{\small 7}),
la sua occorrenza sarà perduta.

Questi esempi ci mostrano come per poter eseguire una gestione effettiva dei
segnali occorrono delle funzioni più sofisticate di quelle finora
illustrate. La funzione \func{signal} infatti ha la sua origine nella
interfaccia alquanto primitiva che venne adottata nei primi sistemi Unix, ma
con questa funzione è sostanzialmente impossibile gestire in maniera adeguata
di tutti i possibili aspetti con cui un processo deve reagire alla ricezione
di un segnale.



\subsection{Gli \textsl{insiemi di segnali} o \textit{signal set}}
\label{sec:sig_sigset}

\itindbeg{signal~set} 

Come evidenziato nel paragrafo precedente, le funzioni di gestione dei segnali
originarie, nate con la semantica inaffidabile, hanno dei limiti non
superabili; in particolare non è prevista nessuna funzione che permetta di
gestire il blocco dei segnali o di verificare lo stato dei segnali pendenti.

Per questo motivo lo standard POSIX.1, insieme alla nuova semantica dei
segnali ha introdotto una interfaccia di gestione completamente nuova, che
permette di ottenere un controllo molto più dettagliato. In particolare lo
standard ha introdotto un nuovo tipo di dato \type{sigset\_t}, che permette di
rappresentare un \textsl{insieme di segnali} (un \textit{signal set}, come
viene usualmente chiamato), tale tipo di dato viene usato per gestire il
blocco dei segnali.

Inizialmente un \textsl{insieme di segnali} veniva rappresentato da un intero
di dimensione opportuna, di solito pari al numero di bit dell'architettura
della macchina, ciascun bit del quale era associato ad uno specifico
segnale. Nel caso di architetture a 32 bit questo comporta un massimo di 32
segnali distinti e dato che a lungo questi sono stati sufficienti non c'era
necessità di nessuna struttura più complicata, in questo modo era possibile
implementare le operazioni direttamente con istruzioni elementari del
processore. 

Oggi questo non è più vero, in particolare con l'introduzione dei segnali
\textit{real-rime} (che vedremo in sez.~\ref{sec:sig_real_time}).  Dato che in
generale non si può fare conto sulle caratteristiche di una implementazione,
perché non è detto che si disponga di un numero di bit sufficienti per mettere
tutti i segnali in un intero, o perché in \type{sigset\_t} possono essere
immagazzinate ulteriori informazioni, tutte le operazioni devono essere
effettuate tramite le opportune funzioni di libreria che si curano di
mascherare i dettagli di basso livello.

Lo standard POSIX.1 definisce cinque funzioni per la manipolazione degli
insiemi di segnali. Le prime quattro, che consentono di manipolare i contenuti
di un \textit{signal set}, sono \funcd{sigemptyset}, \funcd{sigfillset},
\funcd{sigaddset} e \funcd{sigdelset}; i rispettivi prototipi sono:

\begin{funcproto}{
\fhead{signal.h}
\fdecl{int sigemptyset(sigset\_t *set)}
\fdesc{Inizializza un insieme di segnali vuoto.}
\fdecl{int sigfillset(sigset\_t *set)}
\fdesc{Inizializza un insieme di segnali pieno.}
\fdecl{int sigaddset(sigset\_t *set, int signum)}
\fdesc{Aggiunge un segnale ad un insieme di segnali.}
\fdecl{int sigdelset(sigset\_t *set, int signum)}
\fdesc{Rimuove un segnale da un insieme di segnali.}
}

{Le funzioni ritornano $0$ in caso di successo, e $-1$ per un errore, nel qual
  caso \var{errno} assumerà il valore:
  \begin{errlist}
  \item[\errcode{EINVAL}] \param{signum} non è un segnale valido.
  \end{errlist}
}
\end{funcproto}

Le prime due funzioni inizializzano l'insieme di segnali indicato
dall'argomento \param{set} rispettivamente ad un contenuto vuoto (in cui cioè
non c'è nessun segnale) e pieno (in cui cioè ci sono tutti i segnali). Le
altre due funzioni consentono di inserire o rimuovere uno specifico segnale
indicato con l'argomento \param{signum} in un insieme. 

A queste funzioni si aggiunge l'ulteriore \funcd{sigismember}, che consente di
verificare la presenza di un segnale in un insieme, il suo prototipo è:

\begin{funcproto}{
\fhead{signal.h}
\fdecl{int sigismember(const sigset\_t *set, int signum)}
\fdesc{Controlla se un segnale è in un insieme di segnali.}
}

{La funzione ritorna $1$ il segnale è nell'insieme e $0$ altrimenti, e $-1$
  per un errore, nel qual caso \var{errno} assumerà il valore \errval{EINVAL}
  se si è specificato un puntatore \var{NULL}.}
\end{funcproto}

La \acr{glibc} prevede inoltre altre funzioni non standardizzate, accessibili
definendo la macro \macro{\_GNU\_SOURCE}. La prima di queste è
\funcd{sigisemptyset}, che consente di verificare un insieme è vuoto, il suo
prototipo è:

\begin{funcproto}{
\fhead{signal.h}
\fdecl{int sigisemptyset(sigset\_t *set)}
\fdesc{Controlla se un insieme di segnali è vuoto.}
}

{La funzione ritorna $1$ l'insieme è vuoto e $0$ altrimenti, non sono previste
  condizioni di errore.}
\end{funcproto}

Alla precedente si aggiungono altre due funzioni consentono di effettuare
delle operazioni logiche con gli insiemi di segnali, esse sono
\funcd{sigorset} e \funcd{sigandset}, ed i rispettivi prototipi sono:

\begin{funcproto}{
\fhead{signal.h}
\fdecl{sigorset(sigset\_t *dest, sigset\_t *left, sigset\_t *right)}
\fdesc{Crea l'unione di due insieme di segnali.}
\fdecl{sigandset(sigset\_t *dest, sigset\_t *left, sigset\_t *right)}
\fdesc{Crea l'intersezione di due insieme di segnali.} 
}

{Le funzioni ritornano $0$ in caso di successo e $-1$ per un errore, nel qual
  caso \var{errno} assumerà il valore \errcode{EINVAL}.}
\end{funcproto}


In genere si usa un insieme di segnali per specificare quali segnali si vuole
bloccare, o per riottenere dalle varie funzioni di gestione la maschera dei
segnali attivi (vedi sez.~\ref{sec:sig_sigmask}). La modalità più comune, che
è anche quella più portabile, prevede che possano essere definiti aggiungendo
i segnali voluti ad un insieme vuoto ottenuto con \func{sigemptyset} o
togliendo quelli che non servono da un insieme completo ottenuto con
\func{sigfillset}.

\itindend{signal~set} 


\subsection{La funzione \func{sigaction}}
\label{sec:sig_sigaction}

Abbiamo già accennato in sez.~\ref{sec:sig_signal} i problemi di compatibilità
relativi all'uso di \func{signal}. Per ovviare a tutto questo lo standard
POSIX.1 ha ridefinito completamente l'interfaccia per la gestione dei segnali,
rendendola molto più flessibile e robusta, anche se leggermente più complessa.

La principale funzione di sistema prevista dall'interfaccia POSIX.1 per la
gestione dei segnali è \funcd{sigaction}. Essa ha sostanzialmente lo stesso
uso di \func{signal}, permette cioè di specificare le modalità con cui un
segnale può essere gestito da un processo. Il suo prototipo è:

\begin{funcproto}{
\fhead{signal.h}
\fdecl{int sigaction(int signum, const struct sigaction *act, struct sigaction
  *oldact)}  
\fdesc{Installa una nuova azione per un segnale.} 
}

{La funzione ritorna $0$ in caso di successo e $-1$ per un errore, nel qual
  caso \var{errno} assumerà uno dei valori: 
  \begin{errlist}
  \item[\errcode{EFAULT}] si sono specificati indirizzi non validi.
  \item[\errcode{EINVAL}] si è specificato un numero di segnale invalido o si è
    cercato di installare il gestore per \signal{SIGKILL} o
    \signal{SIGSTOP}.
  \end{errlist}
}
\end{funcproto}

La funzione serve ad installare una nuova \textsl{azione} per il segnale
indicato dall'argomento \param{signum}. Si parla di \textsl{azione} e non di
\textsl{gestore} come nel caso di \func{signal}, in quanto la funzione
consente di specificare le varie caratteristiche della risposta al segnale,
non solo la funzione che verrà eseguita alla sua occorrenza.  

Per questo motivo lo standard POSIX.1 raccomanda di usare sempre questa
funzione al posto della precedente \func{signal}, che in genere viene
ridefinita in termini di \func{sigaction}, in quanto la nuova interfaccia
permette un controllo completo su tutti gli aspetti della gestione di un
segnale, sia pure al prezzo di una maggiore complessità d'uso.

Se il puntatore \param{act} non è nullo, la funzione installa la nuova azione
da esso specificata, se \param{oldact} non è nullo il valore dell'azione
corrente viene restituito indietro.  Questo permette (specificando \param{act}
nullo e \param{oldact} non nullo) di superare uno dei limiti di \func{signal},
che non consente di ottenere l'azione corrente senza installarne una nuova. Se
sia \param{act} che \param{oldact} la funzione può essere utilizzata per
verificare, se da luogo ad un errore, se il segnale indicato è valido per la
piattaforma che si sta usando.

Entrambi i puntatori fanno riferimento alla struttura \struct{sigaction},
tramite la quale si specificano tutte le caratteristiche dell'azione associata
ad un segnale.  Anch'essa è descritta dallo standard POSIX.1 ed in Linux è
definita secondo quanto riportato in fig.~\ref{fig:sig_sigaction}. Il campo
\var{sa\_restorer}, non previsto dallo standard, è obsoleto e non deve essere
più usato.

\begin{figure}[!htb]
  \footnotesize \centering
  \begin{minipage}[c]{0.8\textwidth}
    \includestruct{listati/sigaction.h}
  \end{minipage} 
  \normalsize 
  \caption{La struttura \structd{sigaction}.} 
  \label{fig:sig_sigaction}
\end{figure}

Il campo \var{sa\_mask} serve ad indicare l'insieme dei segnali che devono
essere bloccati durante l'esecuzione del gestore, ad essi viene comunque
sempre aggiunto il segnale che ne ha causato la chiamata, a meno che non si
sia specificato con \var{sa\_flag} un comportamento diverso. Quando il
gestore ritorna comunque la maschera dei segnali bloccati (vedi
sez.~\ref{sec:sig_sigmask}) viene ripristinata al valore precedente
l'invocazione.

L'uso di questo campo permette ad esempio di risolvere il problema residuo
dell'implementazione di \code{sleep} mostrata in
fig.~\ref{fig:sig_sleep_incomplete}. In quel caso infatti se il segnale di
allarme avesse interrotto un altro gestore questo non sarebbe stato eseguito
correttamente, la cosa poteva essere prevenuta installando gli altri gestori
usando \var{sa\_mask} per bloccare \signal{SIGALRM} durante la loro
esecuzione.  Il valore di \var{sa\_flag} permette di specificare vari aspetti
del comportamento di \func{sigaction}, e della reazione del processo ai vari
segnali; i valori possibili ed il relativo significato sono riportati in
tab.~\ref{tab:sig_sa_flag}.

\begin{table}[!htb]
  \footnotesize
  \centering
  \begin{tabular}[c]{|l|p{8cm}|}
    \hline
    \textbf{Valore} & \textbf{Significato} \\
    \hline
    \hline
    \constd{SA\_NOCLDSTOP}& Se il segnale è \signal{SIGCHLD} allora non deve
                            essere notificato quando il processo figlio viene
                            fermato da uno dei segnali \signal{SIGSTOP},
                            \signal{SIGTSTP}, \signal{SIGTTIN} o 
                            \signal{SIGTTOU}, questo flag ha significato solo
                            quando si imposta un gestore per \signal{SIGCHLD}.\\
    \constd{SA\_NOCLDWAIT}& Se il segnale è \signal{SIGCHLD} e si richiede di
                            ignorare il segnale con \const{SIG\_IGN} allora i
                            processi figli non diventano \textit{zombie} quando
                            terminano; questa funzionalità è stata introdotta
                            nel kernel 2.6 e va a modificare il comportamento
                            di \func{waitpid} come illustrato in
                            sez.~\ref{sec:proc_wait}, se si installa un gestore
                            con questo flag attivo il segnale \signal{SIGCHLD}
                            viene comunque generato.\\
    \constd{SA\_NODEFER}  & Evita che il segnale corrente sia bloccato durante
                            l'esecuzione del gestore.\\
    \constd{SA\_NOMASK}   & Nome obsoleto e sinonimo non standard di
                            \const{SA\_NODEFER}, non deve essere più
                            utilizzato.\\ 
    \constd{SA\_ONESHOT}  & Nome obsoleto e sinonimo non standard di
                            \const{SA\_RESETHAND}, non deve essere più
                            utilizzato.\\ 
    \constd{SA\_ONSTACK}  & Stabilisce l'uso di uno \textit{stack} alternativo
                            per l'esecuzione del gestore (vedi
                            sez.~\ref{sec:sig_specific_features}).\\  
    \constd{SA\_RESETHAND}& Ristabilisce l'azione per il segnale al valore 
                            predefinito una volta che il gestore è stato
                            lanciato, riproduce cioè il comportamento della
                            semantica inaffidabile.\\  
    \constd{SA\_RESTART}  & Riavvia automaticamente le \textit{slow system
                            call} quando vengono interrotte dal suddetto
                            segnale, riproduce cioè il comportamento standard
                            di BSD.\\ 
    \constd{SA\_SIGINFO}  & Deve essere specificato quando si vuole usare un
                            gestore in forma estesa usando
                            \var{sa\_sigaction} al posto di
                            \var{sa\_handler}.\\
    \hline
  \end{tabular}
  \caption{Valori del campo \var{sa\_flag} della struttura \struct{sigaction}.}
  \label{tab:sig_sa_flag}
\end{table}

Come si può notare in fig.~\ref{fig:sig_sigaction} \func{sigaction} permette
di utilizzare due forme diverse di gestore,\footnote{la possibilità è prevista
  dallo standard POSIX.1b, ed è stata aggiunta nei kernel della serie 2.1.x
  con l'introduzione dei segnali \textit{real-time} (vedi
  sez.~\ref{sec:sig_real_time}); in precedenza era possibile ottenere alcune
  informazioni addizionali usando \var{sa\_handler} con un secondo parametro
  addizionale di tipo \var{sigcontext}, che adesso è deprecato.}  da
specificare, a seconda dell'uso o meno del flag \const{SA\_SIGINFO},
rispettivamente attraverso i campi \var{sa\_sigaction} o \var{sa\_handler}.
Quest'ultima è quella classica usata anche con \func{signal}, mentre la prima
permette di usare un gestore più complesso, in grado di ricevere informazioni
più dettagliate dal sistema, attraverso la struttura \struct{siginfo\_t},
riportata in fig.~\ref{fig:sig_siginfo_t}.  I due campi devono essere usati in
maniera alternativa, in certe implementazioni questi campi vengono addirittura
definiti come una \direct{union}.\footnote{la direttiva \direct{union} del
  linguaggio C definisce una variabile complessa, analoga a una stuttura, i
  cui campi indicano i diversi tipi di valori che possono essere salvati, in
  maniera alternativa, all'interno della stessa.}

Installando un gestore di tipo \var{sa\_sigaction} diventa allora possibile
accedere alle informazioni restituite attraverso il puntatore a questa
struttura. Tutti i segnali impostano i campi \var{si\_signo}, che riporta il
numero del segnale ricevuto, \var{si\_errno}, che riporta, quando diverso da
zero, il codice dell'errore associato al segnale, e \var{si\_code}, che viene
usato dal kernel per specificare maggiori dettagli riguardo l'evento che ha
causato l'emissione del segnale.

\begin{figure}[!htb]
  \footnotesize \centering
  \begin{minipage}[c]{0.9\textwidth}
    \includestruct{listati/siginfo_t.h}
  \end{minipage} 
  \normalsize 
  \caption{La struttura \structd{siginfo\_t}.} 
  \label{fig:sig_siginfo_t}
\end{figure}
 
In generale \var{si\_code} contiene, per i segnali generici, per quelli
\textit{real-time} e per tutti quelli inviati tramite da un processo con
\func{kill} o affini, le informazioni circa l'origine del segnale stesso, ad
esempio se generato dal kernel, da un timer, da \func{kill}, ecc. Il valore
viene sempre espresso come una costante,\footnote{le definizioni di tutti i
  valori possibili si trovano in \file{bits/siginfo.h}.} ed i valori possibili
in questo caso sono riportati in tab.~\ref{tab:sig_si_code_generic}.

Nel caso di alcuni segnali però il valore di \var{si\_code} viene usato per
fornire una informazione specifica relativa alle motivazioni della ricezione
dello stesso; ad esempio i vari segnali di errore (\signal{SIGILL},
\signal{SIGFPE}, \signal{SIGSEGV} e \signal{SIGBUS}) lo usano per fornire
maggiori dettagli riguardo l'errore, come il tipo di errore aritmetico, di
istruzione illecita o di violazione di memoria; mentre alcuni segnali di
controllo (\signal{SIGCHLD}, \signal{SIGTRAP} e \signal{SIGPOLL}) forniscono
altre informazioni specifiche.

\begin{table}[!htb]
  \footnotesize
  \centering
  \begin{tabular}[c]{|l|p{10cm}|}
    \hline
    \textbf{Valore} & \textbf{Significato} \\
    \hline
    \hline
    \constd{SI\_USER}   & Generato da \func{kill} o \func{raise} o affini.\\
    \constd{SI\_KERNEL} & Inviato direttamente dal kernel.\\
    \constd{SI\_QUEUE}  & Inviato con \func{sigqueue} (vedi
                          sez.~\ref{sec:sig_real_time}).\\ 
    \constd{SI\_TIMER}  & Scadenza di un \textit{POSIX timer} (vedi
                          sez.~\ref{sec:sig_timer_adv}).\\ 
    \constd{SI\_MESGQ}  & Inviato al cambiamento di stato di una coda di
                          messaggi POSIX (vedi sez.~\ref{sec:ipc_posix_mq}),
                          introdotto con il kernel 2.6.6.\\ 
    \constd{SI\_ASYNCIO}& Una operazione di I/O asincrono (vedi
                          sez.~\ref{sec:file_asyncronous_io}) è stata
                          completata.\\
    \constd{SI\_SIGIO}  & Segnale di \signal{SIGIO} da una coda (vedi
                          sez.~\ref{sec:file_asyncronous_operation}).\\ 
    \constd{SI\_TKILL}  & Inviato da \func{tkill} o \func{tgkill} (vedi
                          sez.~\ref{cha:thread_xxx}), introdotto con il kernel
                          2.4.19.\\ 
    \hline
  \end{tabular}
  \caption{Valori del campo \var{si\_code} della struttura \struct{sigaction}
    per i segnali generici.}
  \label{tab:sig_si_code_generic}
\end{table}


In questo caso il valore del campo \var{si\_code} deve essere verificato nei
confronti delle diverse costanti previste per ciascuno di detti segnali; dato
che si tratta di costanti, e non di una maschera binaria, i valori numerici
vengono riutilizzati e ciascuno di essi avrà un significato diverso a seconda
del segnale a cui è associato. 

L'elenco dettagliato dei nomi di queste costanti è riportato nelle diverse
sezioni di tab.~\ref{tab:sig_si_code_special} che sono state ordinate nella
sequenza in cui si sono appena citati i rispettivi segnali, il prefisso del
nome indica comunque in maniera diretta il segnale a cui le costanti fanno
riferimento.

\begin{table}[!htb]
  \footnotesize
  \centering
  \begin{tabular}[c]{|l|p{8cm}|}
    \hline
    \textbf{Valore} & \textbf{Significato} \\
    \hline
    \hline
    \constd{ILL\_ILLOPC}  & Codice di operazione illegale.\\
    \constd{ILL\_ILLOPN}  & Operando illegale.\\
    \constd{ILL\_ILLADR}  & Modo di indirizzamento illegale.\\
    \constd{ILL\_ILLTRP}  & Trappola di processore illegale.\\
    \constd{ILL\_PRVOPC}  & Codice di operazione privilegiato.\\
    \constd{ILL\_PRVREG}  & Registro privilegiato.\\
    \constd{ILL\_COPROC}  & Errore del coprocessore.\\
    \constd{ILL\_BADSTK}  & Errore nello stack interno.\\
    \hline
    \constd{FPE\_INTDIV}  & Divisione per zero intera.\\
    \constd{FPE\_INTOVF}  & Overflow intero.\\
    \constd{FPE\_FLTDIV}  & Divisione per zero in virgola mobile.\\
    \constd{FPE\_FLTOVF}  & Overflow in virgola mobile.\\
    \constd{FPE\_FLTUND}  & Underflow in virgola mobile.\\
    \constd{FPE\_FLTRES}  & Risultato in virgola mobile non esatto.\\
    \constd{FPE\_FLTINV}  & Operazione in virgola mobile non valida.\\
    \constd{FPE\_FLTSUB}  & Mantissa? fuori intervallo.\\
    \hline
    \constd{SEGV\_MAPERR} & Indirizzo non mappato.\\
    \constd{SEGV\_ACCERR} & Permessi non validi per l'indirizzo.\\
    \hline
    \constd{BUS\_ADRALN}  & Allineamento dell'indirizzo non valido.\\
    \constd{BUS\_ADRERR}  & Indirizzo fisico inesistente.\\
    \constd{BUS\_OBJERR}  & Errore hardware sull'indirizzo.\\
    \hline
    \constd{TRAP\_BRKPT}  & Breakpoint sul processo.\\
    \constd{TRAP\_TRACE}  & Trappola di tracciamento del processo.\\
    \hline
    \constd{CLD\_EXITED}  & Il figlio è uscito.\\
    \constd{CLD\_KILLED}  & Il figlio è stato terminato.\\
    \constd{CLD\_DUMPED}  & Il figlio è terminato in modo anormale.\\
    \constd{CLD\_TRAPPED} & Un figlio tracciato ha raggiunto una trappola.\\
    \constd{CLD\_STOPPED} & Il figlio è stato fermato.\\
    \constd{CLD\_CONTINUED}& Il figlio è ripartito.\\
    \hline
    \constd{POLL\_IN}   & Disponibili dati in ingresso.\\
    \constd{POLL\_OUT}  & Spazio disponibile sul buffer di uscita.\\
    \constd{POLL\_MSG}  & Disponibili messaggi in ingresso.\\
    \constd{POLL\_ERR}  & Errore di I/O.\\
    \constd{POLL\_PRI}  & Disponibili dati di alta priorità in ingresso.\\
    \constd{POLL\_HUP}  & Il dispositivo è stato disconnesso.\\
    \hline
  \end{tabular}
  \caption{Valori del campo \var{si\_code} della struttura \struct{sigaction}
    impostati rispettivamente dai segnali \signal{SIGILL}, \signal{SIGFPE},
    \signal{SIGSEGV}, \signal{SIGBUS}, \signal{SIGCHLD}, \signal{SIGTRAP} e
    \signal{SIGPOLL}/\signal{SIGIO}.}
  \label{tab:sig_si_code_special}
\end{table}

Il resto della struttura \struct{siginfo\_t} è definito come una \dirct{union}
ed i valori eventualmente presenti dipendono dal segnale ricevuto, così
\signal{SIGCHLD} ed i segnali \textit{real-time} (vedi
sez.~\ref{sec:sig_real_time}) inviati tramite \func{kill} avvalorano
\var{si\_pid} e \var{si\_uid} coi valori corrispondenti al processo che ha
emesso il segnale, \signal{SIGCHLD} avvalora anche i campi \var{si\_status},
\var{si\_utime} e \var{si\_stime} che indicano rispettivamente lo stato di
uscita, l'\textit{user time} e il \textit{system time} (vedi
sez.~\ref{sec:sys_cpu_times}) usati dal processo; \signal{SIGILL},
\signal{SIGFPE}, \signal{SIGSEGV} e \signal{SIGBUS} avvalorano \var{si\_addr}
con l'indirizzo in cui è avvenuto l'errore, \signal{SIGIO} (vedi
sez.~\ref{sec:file_asyncronous_io}) avvalora \var{si\_fd} con il numero del
file descriptor e \var{si\_band} per i dati urgenti (vedi
sez.~\ref{sec:TCP_urgent_data}) su un socket, il segnale inviato alla scadenza
di un POSIX timer (vedi sez.~\ref{sec:sig_timer_adv}) avvalora i campi
\var{si\_timerid} e \var{si\_overrun}.

Benché sia possibile usare nello stesso programma sia \func{sigaction} che
\func{signal} occorre molta attenzione, in quanto le due funzioni possono
interagire in maniera anomala. Infatti l'azione specificata con
\struct{sigaction} contiene un maggior numero di informazioni rispetto al
semplice indirizzo del gestore restituito da \func{signal}.  Per questo motivo
se si usa quest'ultima per installare un gestore sostituendone uno
precedentemente installato con \func{sigaction}, non sarà possibile effettuare
un ripristino corretto dello stesso.

Per questo è sempre opportuno usare \func{sigaction}, che è in grado di
ripristinare correttamente un gestore precedente, anche se questo è stato
installato con \func{signal}. In generale poi non è il caso di usare il valore
di ritorno di \func{signal} come campo \var{sa\_handler}, o viceversa, dato
che in certi sistemi questi possono essere diversi. In definitiva dunque, a
meno che non si sia vincolati all'aderenza stretta allo standard ISO C, è
sempre il caso di evitare l'uso di \func{signal} a favore di \func{sigaction}.

\begin{figure}[!htbp]
  \footnotesize  \centering
  \begin{minipage}[c]{\codesamplewidth}
    \includecodesample{listati/Signal.c}
  \end{minipage} 
  \normalsize 
  \caption{La funzione \func{Signal}, equivalente a \func{signal}, definita
    attraverso \func{sigaction}.}
  \label{fig:sig_Signal_code}
\end{figure}

Per questo motivo si è provveduto, per mantenere un'interfaccia semplificata
che abbia le stesse caratteristiche di \func{signal}, a definire attraverso
\func{sigaction} una funzione equivalente \func{Signal}, il cui codice è
riportato in fig.~\ref{fig:sig_Signal_code} (il codice completo si trova nel
file \file{SigHand.c} nei sorgenti allegati). Anche in questo caso, per
semplificare la definizione si è poi definito un apposito tipo
\texttt{SigFunc} per esprimere in modo più comprensibile la forma di un
gestore di segnale.

Si noti come, essendo la funzione estremamente semplice, essa è definita come
\dirct{inline}. Questa direttiva viene usata per dire al compilatore di
trattare la funzione cui essa fa riferimento in maniera speciale inserendo il
codice direttamente nel testo del programma.  Anche se i compilatori più
moderni sono in grado di effettuare da soli queste manipolazioni (impostando
le opportune ottimizzazioni) questa è una tecnica usata per migliorare le
prestazioni per le funzioni piccole ed usate di frequente, in particolare nel
kernel, dove in certi casi le ottimizzazioni dal compilatore, tarate per l'uso
in \textit{user space}, non sono sempre adatte.

In tal caso infatti le istruzioni per creare un nuovo frame nello
\textit{stack} per chiamare la funzione costituirebbero una parte rilevante
del codice, appesantendo inutilmente il programma.  Originariamente questo
comportamento veniva ottenuto con delle macro, ma queste hanno tutta una serie
di problemi di sintassi nel passaggio degli argomenti (si veda ad esempio
\cite{PratC}) che in questo modo possono essere evitati.



\subsection{La gestione della \textsl{maschera dei segnali} o 
  \textit{signal mask}}
\label{sec:sig_sigmask}

\index{maschera dei segnali|(}

Come spiegato in sez.~\ref{sec:sig_semantics} tutti i moderni sistemi
unix-like permettono di bloccare temporaneamente (o di eliminare
completamente, impostando come azione \const{SIG\_IGN}) la consegna dei
segnali ad un processo. Questo è fatto specificando la cosiddetta
\textsl{maschera dei segnali} (o \textit{signal mask}) del
processo\footnote{nel caso di Linux essa è mantenuta dal campo \var{blocked}
  della \struct{task\_struct} del processo.} cioè l'insieme dei segnali la cui
consegna è bloccata.

Abbiamo accennato in sez.~\ref{sec:proc_fork} che la maschera dei segnali
viene ereditata dal padre alla creazione di un processo figlio, e abbiamo
visto al paragrafo precedente che essa può essere modificata durante
l'esecuzione di un gestore ed automaticamente ripristinata quando questo
ritorna, attraverso l'uso dal campo \var{sa\_mask} di \struct{sigaction}.

Uno dei problemi evidenziatisi con l'esempio di fig.~\ref{fig:sig_event_wrong}
è che in molti casi è necessario proteggere delle sezioni di codice, in modo
da essere sicuri che essi siano eseguite senza interruzioni da parte di un
segnale.  Nel caso in questione si trattava della sezione di codice fra il
controllo e la eventuale cancellazione del flag impostato dal gestore di un
segnale che testimoniava l'avvenuta occorrenza dello stesso.

Come illustrato in sez.~\ref{sec:proc_atom_oper} le operazioni più semplici,
come l'assegnazione o il controllo di una variabile, di norma sono atomiche, e
qualora si voglia essere sicuri si può usare il tipo \type{sig\_atomic\_t}. Ma
quando si devono eseguire più operazioni su delle variabili (nell'esempio
citato un controllo ed una assegnazione) o comunque eseguire una serie di
istruzioni, l'atomicità non è più possibile.

In questo caso, se si vuole essere sicuri di non poter essere interrotti da un
segnale durante l'esecuzione di una sezione di codice, lo si può bloccare
esplicitamente modificando la maschera dei segnali del processo con la
funzione di sistema \funcd{sigprocmask}, il cui prototipo è:

\begin{funcproto}{
\fhead{signal.h}
\fdecl{int sigprocmask(int how, const sigset\_t *set, sigset\_t *oldset)}
\fdesc{Imposta la maschera dei segnali del processo corrente.} 
}

{La funzione ritorna $0$ in caso di successo e $-1$ per un errore, nel qual
  caso \var{errno} assumerà uno dei valori: 
  \begin{errlist}
  \item[\errcode{EFAULT}] si sono specificati indirizzi non validi.
  \item[\errcode{EINVAL}] si è specificato un numero di segnale invalido.
  \end{errlist}
}
\end{funcproto}

La funzione usa l'insieme di segnali posto all'indirizzo passato
nell'argomento \param{set} per modificare la maschera dei segnali del processo
corrente. La modifica viene effettuata a seconda del valore
dell'argomento \param{how}, secondo le modalità specificate in
tab.~\ref{tab:sig_procmask_how}. Qualora si specifichi un valore non nullo
per \param{oldset} la maschera dei segnali corrente viene salvata a
quell'indirizzo.

\begin{table}[htb]
  \footnotesize
  \centering
  \begin{tabular}[c]{|l|p{8cm}|}
    \hline
    \textbf{Valore} & \textbf{Significato} \\
    \hline
    \hline
    \constd{SIG\_BLOCK}   & L'insieme dei segnali bloccati è l'unione fra
                            quello specificato e quello corrente.\\
    \constd{SIG\_UNBLOCK} & I segnali specificati in \param{set} sono rimossi
                            dalla maschera dei segnali, specificare la
                            cancellazione di un segnale non bloccato è legale.\\
    \constd{SIG\_SETMASK} & La maschera dei segnali è impostata al valore
                            specificato da \param{set}.\\
    \hline
  \end{tabular}
  \caption{Valori e significato dell'argomento \param{how} della funzione
    \func{sigprocmask}.}
  \label{tab:sig_procmask_how}
\end{table}

In questo modo diventa possibile proteggere delle sezioni di codice bloccando
l'insieme di segnali voluto per poi riabilitarli alla fine della sezione
critica. La funzione permette di risolvere problemi come quelli mostrati in
fig.~\ref{fig:sig_event_wrong}, proteggendo la sezione fra il controllo del
flag e la sua cancellazione.  La funzione può essere usata anche all'interno
di un gestore, ad esempio per riabilitare la consegna del segnale che l'ha
invocato, in questo caso però occorre ricordare che qualunque modifica alla
maschera dei segnali viene perduta al ritorno dallo stesso.

Benché con l'uso di \func{sigprocmask} si possano risolvere la maggior parte
dei casi di \textit{race condition} restano aperte alcune possibilità legate
all'uso di \func{pause}.  Il caso è simile a quello del problema illustrato
nell'esempio di fig.~\ref{fig:sig_sleep_incomplete}, e cioè la possibilità che
il processo riceva il segnale che si intende usare per uscire dallo stato di
attesa invocato con \func{pause} immediatamente prima dell'esecuzione di
quest'ultima. Per poter effettuare atomicamente la modifica della maschera dei
segnali (di solito attivandone uno specifico) insieme alla sospensione del
processo lo standard POSIX ha previsto la funzione di sistema
\funcd{sigsuspend}, il cui prototipo è:

\begin{funcproto}{
\fhead{signal.h}
\fdecl{int sigsuspend(const sigset\_t *mask)} 
\fdesc{Imposta la maschera dei segnali mettendo in attesa il processo.} 
}

{La funzione ritorna $0$ in caso di successo e $-1$ per un errore, nel qual
  caso \var{errno} assumerà uno dei valori: 
  \begin{errlist}
  \item[\errcode{EFAULT}] si sono specificati indirizzi non validi.
  \item[\errcode{EINVAL}] si è specificato un numero di segnale invalido.
  \end{errlist}
}
\end{funcproto}

Come esempio dell'uso di queste funzioni proviamo a riscrivere un'altra volta
l'esempio di implementazione di \code{sleep}. Abbiamo accennato in
sez.~\ref{sec:sig_sigaction} come con \func{sigaction} sia possibile bloccare
\signal{SIGALRM} nell'installazione dei gestori degli altri segnali, per poter
usare l'implementazione vista in fig.~\ref{fig:sig_sleep_incomplete} senza
interferenze.  Questo però comporta una precauzione ulteriore al semplice uso
della funzione, vediamo allora come usando la nuova interfaccia è possibile
ottenere un'implementazione, riportata in fig.~\ref{fig:sig_sleep_ok} che non
presenta neanche questa necessità.

\begin{figure}[!htbp]
  \footnotesize \centering
  \begin{minipage}[c]{\codesamplewidth}
    \includecodesample{listati/sleep.c}
  \end{minipage} 
  \normalsize 
  \caption{Una implementazione completa di \func{sleep}.} 
  \label{fig:sig_sleep_ok}
\end{figure}
 
Per evitare i problemi di interferenza con gli altri segnali in questo caso
non si è usato l'approccio di fig.~\ref{fig:sig_sleep_incomplete} evitando
l'uso di \func{longjmp}. Come in precedenza il gestore (\texttt{\small
  27--30}) non esegue nessuna operazione, limitandosi a ritornare per
interrompere il programma messo in attesa.

La prima parte della funzione (\texttt{\small 6--10}) provvede ad installare
l'opportuno gestore per \signal{SIGALRM}, salvando quello originario, che
sarà ripristinato alla conclusione della stessa (\texttt{\small 23}); il passo
successivo è quello di bloccare \signal{SIGALRM} (\texttt{\small 11--14}) per
evitare che esso possa essere ricevuto dal processo fra l'esecuzione di
\func{alarm} (\texttt{\small 16}) e la sospensione dello stesso. Nel fare
questo si salva la maschera corrente dei segnali, che sarà ripristinata alla
fine (\texttt{\small 22}), e al contempo si prepara la maschera dei segnali
\var{sleep\_mask} per riattivare \signal{SIGALRM} all'esecuzione di
\func{sigsuspend}.  

In questo modo non sono più possibili \textit{race condition} dato che
\signal{SIGALRM} viene disabilitato con \func{sigprocmask} fino alla chiamata
di \func{sigsuspend}.  Questo metodo è assolutamente generale e può essere
applicato a qualunque altra situazione in cui si deve attendere per un
segnale, i passi sono sempre i seguenti:
\begin{enumerate*}
\item leggere la maschera dei segnali corrente e bloccare il segnale voluto
  con \func{sigprocmask};
\item mandare il processo in attesa con \func{sigsuspend} abilitando la
  ricezione del segnale voluto;
\item ripristinare la maschera dei segnali originaria.
\end{enumerate*}
Per quanto possa sembrare strano bloccare la ricezione di un segnale per poi
riabilitarla immediatamente dopo, in questo modo si evita il \textit{deadlock}
dovuto all'arrivo del segnale prima dell'esecuzione di \func{sigsuspend}.

\index{maschera dei segnali|)}


\subsection{Criteri di programmazione per i gestori dei segnali}
\label{sec:sig_signal_handler}

Abbiamo finora parlato dei gestori dei segnali come funzioni chiamate in
corrispondenza della consegna di un segnale. In realtà un gestore non può
essere una funzione qualunque, in quanto esso può essere eseguito in
corrispondenza all'interruzione in un punto qualunque del programma
principale, cosa che ad esempio può rendere problematico chiamare all'interno
di un gestore di segnali la stessa funzione che dal segnale è stata
interrotta.

\index{funzioni!\textit{signal safe}|(}

Il concetto è comunque più generale e porta ad una distinzione fra quelle che
POSIX chiama \textsl{funzioni insicure} (\textit{signal unsafe function}) e
\textsl{funzioni sicure} (o più precisamente \textit{signal safe function}).
Quando un segnale interrompe una funzione insicura ed il gestore chiama al suo
interno una funzione insicura il sistema può dare luogo ad un comportamento
indefinito, la cosa non avviene invece per le funzioni sicure.

Tutto questo significa che la funzione che si usa come gestore di segnale deve
essere programmata con molta cura per evirare questa evenienza e che non è
possibile utilizzare al suo interno una qualunque funzione di sistema, se si
vogliono evitare questi problemi si può ricorrere soltanto all'uso delle
funzioni considerate sicure.

L'elenco delle funzioni considerate sicure varia a seconda della
implementazione utilizzata e dello standard a cui si fa riferimento. Non è
riportata una lista specifica delle funzioni sicure per Linux, e si suppone
pertanto che siano quelle richieste dallo standard. Secondo quanto richiesto
dallo standard POSIX 1003.1 nella revisione del 2003, le ``\textit{signal safe
  function}'' che possono essere chiamate anche all'interno di un gestore di
segnali sono tutte quelle della lista riportata in
fig.~\ref{fig:sig_safe_functions}.

\begin{figure}[!htb]
  \footnotesize \centering
  \begin{minipage}[c]{14cm}
    \func{\_exit}, \func{abort}, \func{accept}, \func{access},
    \func{aio\_error} \func{aio\_return}, \func{aio\_suspend}, \func{alarm},
    \func{bind}, \func{cfgetispeed}, \func{cfgetospeed}, \func{cfsetispeed},
    \func{cfsetospeed}, \func{chdir}, \func{chmod}, \func{chown},
    \func{clock\_gettime}, \func{close}, \func{connect}, \func{creat},
    \func{dup}, \func{dup2}, \func{execle}, \func{execve}, \func{fchmod},
    \func{fchown}, \func{fcntl}, \func{fdatasync}, \func{fork},
    \func{fpathconf}, \func{fstat}, \func{fsync}, \func{ftruncate},
    \func{getegid}, \func{geteuid}, \func{getgid}, \func{getgroups},
    \func{getpeername}, \func{getpgrp}, \func{getpid}, \func{getppid},
    \func{getsockname}, \func{getsockopt}, \func{getuid}, \func{kill},
    \func{link}, \func{listen}, \func{lseek}, \func{lstat}, \func{mkdir},
    \func{mkfifo}, \func{open}, \func{pathconf}, \func{pause}, \func{pipe},
    \func{poll}, \funcm{posix\_trace\_event}, \func{pselect}, \func{raise},
    \func{read}, \func{readlink}, \func{recv}, \func{recvfrom},
    \func{recvmsg}, \func{rename}, \func{rmdir}, \func{select},
    \func{sem\_post}, \func{send}, \func{sendmsg}, \func{sendto},
    \func{setgid}, \func{setpgid}, \func{setsid}, \func{setsockopt},
    \func{setuid}, \func{shutdown}, \func{sigaction}, \func{sigaddset},
    \func{sigdelset}, \func{sigemptyset}, \func{sigfillset},
    \func{sigismember}, \func{signal}, \func{sigpause}, \func{sigpending},
    \func{sigprocmask}, \func{sigqueue}, \funcm{sigset}, \func{sigsuspend},
    \func{sleep}, \func{socket}, \func{socketpair}, \func{stat},
    \func{symlink}, \func{sysconf}, \func{tcdrain}, \func{tcflow},
    \func{tcflush}, \func{tcgetattr}, \func{tcgetgrp}, \func{tcsendbreak},
    \func{tcsetattr}, \func{tcsetpgrp}, \func{time}, \func{timer\_getoverrun},
    \func{timer\_gettime}, \func{timer\_settime}, \func{times}, \func{umask},
    \func{uname}, \func{unlink}, \func{utime}, \func{wait}, \func{waitpid},
    \func{write}.
  \end{minipage} 
  \normalsize 
  \caption{Elenco delle funzioni sicure secondo lo standard POSIX
    1003.1-2003.}
  \label{fig:sig_safe_functions}
\end{figure}

\index{funzioni!\textit{signal safe}|)}

Lo standard POSIX.1-2004 modifica la lista di
fig.~\ref{fig:sig_safe_functions} aggiungendo le funzioni \func{\_Exit} e
\func{sockatmark}, mentre lo standard POSIX.1-2008 rimuove della lista le tre
funzioni \func{fpathconf}, \func{pathconf}, \func{sysconf} e vi aggiunge le
ulteriori funzioni in fig.~\ref{fig:sig_safe_functions_posix_2008}.

\begin{figure}[!htb]
  \footnotesize \centering
  \begin{minipage}[c]{14cm}
     \func{execl}, \func{execv}, \func{faccessat}, \func{fchmodat},
     \func{fchownat}, \func{fexecve}, \func{fstatat}, \func{futimens},
     \func{linkat}, \func{mkdirat}, \func{mkfifoat}, \func{mknod},
     \func{mknodat}, \func{openat}, \func{readlinkat}, \func{renameat},
     \func{symlinkat}, \func{unlinkat}, \func{utimensat}, \func{utimes}.
  \end{minipage} 
  \normalsize 
  \caption{Ulteriori funzioni sicure secondo lo standard POSIX.1-2008.}
  \label{fig:sig_safe_functions_posix_2008}
\end{figure}


Per questo motivo è opportuno mantenere al minimo indispensabile le operazioni
effettuate all'interno di un gestore di segnali, qualora si debbano compiere
operazioni complesse è sempre preferibile utilizzare la tecnica in cui si usa
il gestore per impostare il valore di una qualche variabile globale, e poi si
eseguono le operazioni complesse nel programma verificando (con tutti gli
accorgimenti visti in precedenza) il valore di questa variabile tutte le volte
che si è rilevata una interruzione dovuta ad un segnale.


\section{Funzionalità avanzate}
\label{sec:sig_advanced_signal}

Tratteremo in questa ultima sezione alcune funzionalità avanzate relativa ai
segnali ed in generale ai meccanismi di notifica, a partire dalla funzioni
introdotte per la gestione dei cosiddetti ``\textsl{segnali real-time}'', alla
gestione avanzata delle temporizzazioni e le nuove interfacce per la gestione
di segnali ed eventi attraverso l'uso di file descriptor.

\subsection{I segnali \textit{real-time}}
\label{sec:sig_real_time}

Lo standard POSIX.1b, nel definire una serie di nuove interfacce per i servizi
\textit{real-time}, ha introdotto una estensione del modello classico dei
segnali che presenta dei significativi miglioramenti,\footnote{questa
  estensione è stata introdotta in Linux a partire dal kernel 2.1.43, e dalla
  versione 2.1 della \acr{glibc}.} in particolare sono stati superati tre
limiti fondamentali dei segnali classici:
\begin{basedescript}{\desclabelwidth{1cm}\desclabelstyle{\nextlinelabel}}
\item[\textbf{I segnali non sono accumulati}] 
  se più segnali vengono generati prima dell'esecuzione di un gestore
  questo sarà eseguito una sola volta, ed il processo non sarà in grado di
  accorgersi di quante volte l'evento che ha generato il segnale è accaduto.
\item[\textbf{I segnali non trasportano informazione}]   
  i segnali classici non prevedono altra informazione sull'evento
  che li ha generati se non il fatto che sono stati emessi (tutta
  l'informazione che il kernel associa ad un segnale è il suo numero).
\item[\textbf{I segnali non hanno un ordine di consegna}] 
  l'ordine in cui diversi segnali vengono consegnati è casuale e non
  prevedibile. Non è possibile stabilire una priorità per cui la reazione a
  certi segnali ha la precedenza rispetto ad altri.
\end{basedescript}

Per poter superare queste limitazioni lo standard POSIX.1b ha introdotto delle
nuove caratteristiche, che sono state associate ad una nuova classe di
segnali, che vengono chiamati \textsl{segnali real-time}, in particolare le
funzionalità aggiunte sono:

\begin{enumerate}
\item i segnali sono inseriti in una coda che permette di consegnare istanze
  multiple dello stesso segnale qualora esso venga inviato più volte prima
  dell'esecuzione del gestore; si assicura così che il processo riceva un
  segnale per ogni occorrenza dell'evento che lo genera;
\item è stata introdotta una priorità nella consegna dei segnali: i segnali
  vengono consegnati in ordine a seconda del loro valore, partendo da quelli
  con un numero minore, che pertanto hanno una priorità maggiore;
\item è stata introdotta la possibilità di restituire dei dati al gestore,
  attraverso l'uso di un apposito campo \var{si\_value} nella struttura
  \struct{siginfo\_t}, accessibile tramite gestori di tipo
  \var{sa\_sigaction}.
\end{enumerate}

Tutte queste nuove funzionalità eccetto l'ultima, che, come illustrato in
sez.~\ref{sec:sig_sigaction}, è disponibile anche con i segnali ordinari, si
applicano solo ai nuovi segnali \textit{real-time}; questi ultimi sono
accessibili in un intervallo di valori specificati dalle due costanti
\constd{SIGRTMIN} e \constd{SIGRTMAX}, che specificano il numero minimo e
massimo associato ad un segnale \textit{real-time}.

Su Linux di solito il primo valore è 33, mentre il secondo è \code{\_NSIG-1},
che di norma (vale a dire sulla piattaforma i386) è 64. Questo dà un totale di
32 segnali disponibili, contro gli almeno 8 richiesti da POSIX.1b. Si tenga
presente però che i primi segnali \textit{real-time} disponibili vengono usati
dalla \acr{glibc} per l'implementazione dei \textit{thread} POSIX (vedi
sez.~\ref{sec:thread_posix_intro}), ed il valore di \const{SIGRTMIN} viene
modificato di conseguenza.\footnote{per la precisione vengono usati i primi
  tre per la vecchia implementazione dei \textit{LinuxThread} ed i primi due
  per la nuova NTPL (\textit{New Thread Posix Library}), il che comporta che
  \const{SIGRTMIN} a seconda dei casi può assumere i valori 34 o 35.}

Per questo motivo nei programmi che usano i segnali \textit{real-time} non si
deve mai usare un valore assoluto dato che si correrebbe il rischio di
utilizzare un segnale in uso alle librerie, ed il numero del segnale deve
invece essere sempre specificato in forma relativa a \const{SIGRTMIN} (come
\code{SIGRTMIN + n}) avendo inoltre cura di controllare di non aver mai
superato \const{SIGRTMAX}.

I segnali con un numero più basso hanno una priorità maggiore e vengono
consegnati per primi, inoltre i segnali \textit{real-time} non possono
interrompere l'esecuzione di un gestore di un segnale a priorità più alta; la
loro azione predefinita è quella di terminare il programma.  I segnali
ordinari hanno tutti la stessa priorità, che è più alta di quella di qualunque
segnale \textit{real-time}. Lo standard non definisce niente al riguardo ma
Linux, come molte altre implementazioni, adotta questa politica.

Si tenga presente che questi nuovi segnali non sono associati a nessun evento
specifico, a meno di non richiedere specificamente il loro utilizzo in
meccanismi di notifica come quelli per l'I/O asincrono (vedi
sez.~\ref{sec:file_asyncronous_io}) o per le code di messaggi POSIX (vedi
sez.~\ref{sec:ipc_posix_mq}), pertanto devono essere inviati esplicitamente.

Inoltre, per poter usufruire della capacità di restituire dei dati, i relativi
gestori devono essere installati con \func{sigaction}, specificando per
\var{sa\_flags} la modalità \const{SA\_SIGINFO} che permette di utilizzare la
forma estesa \var{sa\_sigaction} del gestore (vedi
sez.~\ref{sec:sig_sigaction}).  In questo modo tutti i segnali
\textit{real-time} possono restituire al gestore una serie di informazioni
aggiuntive attraverso l'argomento \struct{siginfo\_t}, la cui definizione è
stata già vista in fig.~\ref{fig:sig_siginfo_t}, nella trattazione dei gestori
in forma estesa.

In particolare i campi utilizzati dai segnali \textit{real-time} sono
\var{si\_pid} e \var{si\_uid} in cui vengono memorizzati rispettivamente il
\ids{PID} e l'\ids{UID} effettivo del processo che ha inviato il segnale,
mentre per la restituzione dei dati viene usato il campo \var{si\_value}.

\begin{figure}[!htb]
  \footnotesize \centering
  \begin{minipage}[c]{0.8\textwidth}
    \includestruct{listati/sigval_t.h}
  \end{minipage} 
  \normalsize 
  \caption{La definizione dell'unione \structd{sigval}, definita anche come
    tipo \typed{sigval\_t}.}
  \label{fig:sig_sigval}
\end{figure}

Detto campo, identificato con il tipo di dato \type{sigval\_t}, è una
\dirct{union} di tipo \struct{sigval} (la sua definizione è in
fig.~\ref{fig:sig_sigval}) in cui può essere memorizzato o un valore numerico,
se usata nella forma \var{sival\_int}, o un puntatore, se usata nella forma
\var{sival\_ptr}. L'unione viene usata dai segnali \textit{real-time} e da
vari meccanismi di notifica per restituire dati al gestore del segnale in
\var{si\_value}. Un campo di tipo \type{sigval\_t} è presente anche nella
struttura \struct{sigevent} (definita in fig.~\ref{fig:struct_sigevent}) che
viene usata dai meccanismi di notifica come quelli per i timer POSIX (vedi
sez.~\ref{sec:sig_timer_adv}), l'I/O asincrono (vedi
sez.~\ref{sec:file_asyncronous_io}) o le code di messaggi POSIX (vedi
sez.~\ref{sec:ipc_posix_mq}).

A causa delle loro caratteristiche, la funzione \func{kill} non è adatta ad
inviare segnali \textit{real-time}, poiché non è in grado di fornire alcun
valore per il campo \var{si\_value} restituito nella struttura
\struct{siginfo\_t} prevista da un gestore in forma estesa. Per questo motivo
lo standard ha previsto una nuova funzione, \funcd{sigqueue}, il cui prototipo
è:

\begin{funcproto}{
\fhead{signal.h}
\fdecl{int sigqueue(pid\_t pid, int signo, const union sigval value)}
\fdesc{Invia un segnale con un valore di informazione.} 
}

{La funzione ritorna $0$ in caso di successo e $-1$ per un errore, nel qual
  caso \var{errno} assumerà uno dei valori: 
  \begin{errlist}
  \item[\errcode{EAGAIN}] la coda è esaurita, ci sono già
    \const{SIGQUEUE\_MAX} segnali in attesa si consegna.
  \item[\errcode{EINVAL}] si è specificato un valore non valido per
    \param{signo}.
  \item[\errcode{EPERM}] non si hanno privilegi appropriati per inviare il
    segnale al processo specificato.
  \item[\errcode{ESRCH}] il processo \param{pid} non esiste.
  \end{errlist}
}
\end{funcproto}


La funzione invia il segnale indicato dall'argomento \param{signo} al processo
indicato dall'argomento \param{pid}. Per il resto il comportamento della
funzione è analogo a quello di \func{kill}, ed i privilegi occorrenti ad
inviare il segnale ad un determinato processo sono gli stessi; un valore nullo
di \param{signo} permette di verificare le condizioni di errore senza inviare
nessun segnale.

Se il segnale è bloccato la funzione ritorna immediatamente, se si è
installato un gestore con \const{SA\_SIGINFO} e ci sono risorse disponibili,
(vale a dire che c'è posto nella coda dei segnali \textit{real-time}) esso
viene inserito e diventa pendente. Una volta consegnato il segnale il gestore
otterrà nel campo \var{si\_code} di \struct{siginfo\_t} il valore
\const{SI\_QUEUE} e nel campo \var{si\_value} il valore indicato
nell'argomento \param{value}. Se invece si è installato un gestore nella forma
classica il segnale sarà generato, ma tutte le caratteristiche tipiche dei
segnali \textit{real-time} (priorità e coda) saranno perse.

Per lo standard POSIX la profondità della coda è indicata dalla costante
\constd{SIGQUEUE\_MAX}, una della tante costanti di sistema definite dallo
standard POSIX che non abbiamo riportato esplicitamente in
sez.~\ref{sec:sys_limits}. Il suo valore minimo secondo lo standard,
\macrod{\_POSIX\_SIGQUEUE\_MAX}, è pari a 32. Nel caso di Linux la coda ha una
dimensione variabile; fino alla versione 2.6.7 c'era un limite massimo globale
che poteva essere impostato come parametro del kernel in
\sysctlfiled{kernel/rtsig-max} ed il valore predefinito era pari a 1024. A
partire dal kernel 2.6.8 il valore globale è stato rimosso e sostituito dalla
risorsa \const{RLIMIT\_SIGPENDING} associata al singolo utente, che può essere
modificata con \func{setrlimit} come illustrato in
sez.~\ref{sec:sys_resource_limit}.

Lo standard POSIX.1b definisce inoltre delle nuove funzioni di sistema che
permettono di gestire l'attesa di segnali specifici su una coda, esse servono
in particolar modo nel caso dei \textit{thread}, in cui si possono usare i
segnali \textit{real-time} come meccanismi di comunicazione elementare; la
prima di queste è \funcd{sigwait}, il cui prototipo è:

\begin{funcproto}{
\fhead{signal.h}
\fdecl{int sigwait(const sigset\_t *set, int *sig)}
\fdesc{Attende la ricezione di un segnale.} 
}
{La funzione ritorna $0$ in caso di successo e $-1$ per un errore, nel qual
  caso \var{errno} assumerà uno dei valori: 
  \begin{errlist}
  \item[\errcode{EINTR}] la funzione è stata interrotta.
  \item[\errcode{EINVAL}] si è specificato un valore non valido per
  \end{errlist}
  ed inoltre \errval{EFAULT} nel suo significato generico.}
\end{funcproto}

La funzione estrae dall'insieme dei segnali pendenti uno qualunque fra quelli
indicati nel \textit{signal set} specificato in \param{set}, il cui valore
viene restituito nella variabile puntata da \param{sig}.  Se sono pendenti più
segnali, viene estratto quello a priorità più alta, cioè quello con il numero
più basso. Se, nel caso di segnali \textit{real-time}, c'è più di un segnale
pendente, ne verrà estratto solo uno. Una volta estratto il segnale non verrà
più consegnato, e se era in una coda il suo posto sarà liberato. Se non c'è
nessun segnale pendente il processo viene bloccato fintanto che non ne arriva
uno.

Per un funzionamento corretto la funzione richiede che alla sua chiamata i
segnali di \param{set} siano bloccati. In caso contrario si avrebbe un
conflitto con gli eventuali gestori: pertanto non si deve utilizzare per
lo stesso segnale questa funzione e \func{sigaction}. Se questo non avviene il
comportamento del sistema è indeterminato: il segnale può sia essere
consegnato che essere ricevuto da \func{sigwait}, il tutto in maniera non
prevedibile.

Lo standard POSIX.1b definisce altre due funzioni di sistema, anch'esse usate
prevalentemente con i \textit{thread}; \funcd{sigwaitinfo} e
\funcd{sigtimedwait}, i relativi prototipi sono:

\begin{funcproto}{
\fhead{signal.h}
\fdecl{int sigwaitinfo(const sigset\_t *set, siginfo\_t *info)}
\fdesc{Attende un segnale con le relative informazioni.}
\fdecl{int sigtimedwait(const sigset\_t *set, siginfo\_t *info, const
  struct timespec *timeout)}
\fdesc{Attende un segnale con le relative informazioni per un tempo massimo.}
}

{Le funzioni ritornano $0$ in caso di successo e $-1$ per un errore, nel qual
  caso \var{errno} assumerà uno gli stessi valori di \func{sigwait} ai quali
  si aggiunge per \func{sigtimedwait}:
  \begin{errlist}
  \item[\errcode{EAGAIN}] si è superato il timeout senza che un segnale atteso
    sia stato ricevuto.
  \end{errlist}
}
\end{funcproto}


Entrambe le funzioni sono estensioni di \func{sigwait}. La prima permette di
ricevere, oltre al numero del segnale, anche le informazioni ad esso associate
tramite l'argomento \param{info}; in particolare viene restituito il numero
del segnale nel campo \var{si\_signo}, la sua causa in \var{si\_code}, e se il
segnale è stato immesso sulla coda con \func{sigqueue}, il valore di ritorno
ad esso associato viene riportato in \var{si\_value}, che altrimenti è
indefinito.

La seconda è identica alla prima ma in più permette di specificare un timeout
con l'argomento omonimo, scaduto il quale ritornerà con un errore. Se si
specifica per \param{timeout} un puntatore nullo il comportamento sarà
identico a \func{sigwaitinfo}. Se si specifica un tempo di timeout nullo e non
ci sono segnali pendenti la funzione ritornerà immediatamente, in questo modo
si può eliminare un segnale dalla coda senza dover essere bloccati qualora
esso non sia presente.

L'uso di queste funzioni è principalmente associato alla gestione dei segnali
con i \textit{thread}. In genere esse vengono chiamate dal \textit{thread}
incaricato della gestione, che al ritorno della funzione esegue il codice che
usualmente sarebbe messo nel gestore, per poi ripetere la chiamata per
mettersi in attesa del segnale successivo. Questo ovviamente comporta che non
devono essere installati gestori, che solo il \textit{thread} di gestione deve
usare \func{sigwait} e che i segnali gestiti in questa maniera, per evitare
che venga eseguita l'azione predefinita, devono essere mascherati per tutti i
\textit{thread}, compreso quello dedicato alla gestione, che potrebbe
riceverlo fra due chiamate successive.


\subsection{La gestione avanzata delle temporizzazioni}
\label{sec:sig_timer_adv}

Sia le funzioni per la gestione dei tempi viste in
sez.~\ref{sec:sys_cpu_times} che quelle per la gestione dei timer di
sez.~\ref{sec:sig_alarm_abort} sono state a lungo limitate dalla risoluzione
massima dei tempi dell'orologio interno del kernel, che era quella ottenibile
dal timer di sistema che governa lo \textit{scheduler}, e quindi limitate
dalla frequenza dello stesso che si ricordi, come già illustrato in
sez.~\ref{sec:proc_hierarchy}, è data dal valore della costante \texttt{HZ}. 

I contatori usati per il calcolo dei tempi infatti erano basati sul numero di
\textit{jiffies} che vengono incrementati ad ogni \textit{clock tick} del
timer di sistema, il che comportava anche, come accennato in
sez.~\ref{sec:sig_alarm_abort} per \func{setitimer}, problemi per il massimo
periodo di tempo copribile da alcuni di questi orologi, come quelli associati
al \textit{process time} almeno fino a quando, con il kernel 2.6.16, non è
stato rimosso il limite di un valore a 32 bit per i \textit{jiffies}.

\itindbeg{POSIX~Timer~API}

Nelle architetture moderne però tutti i computer sono dotati di temporizzatori
hardware che possono supportare risoluzioni molto elevate, ed in maniera del
tutto indipendente dalla frequenza scelta per il timer di sistema che governa
lo \textit{scheduler}, normalmente si possono ottenere precisioni fino al
microsecondo, andando molto oltre in caso di hardware dedicato. 

Per questo lo standard POSIX.1-2001 ha previsto una serie di nuove funzioni
relative a quelli che vengono chiamati ``\textsl{orologi}
\textit{real-time}'', in grado di supportare risoluzioni fino al
nanosecondo. Inoltre le CPU più moderne sono dotate a loro volta di contatori
ad alta definizione che consentono una grande accuratezza nella misura del
tempo da esse dedicato all'esecuzione di un processo.

Per usare queste funzionalità ed ottenere risoluzioni temporali più accurate,
occorre però un opportuno supporto da parte del kernel, ed i cosiddetti
\itindex{High~Resolution~Timer~(HRT)} \textit{high resolution timer} che
consentono di fare ciò sono stati introdotti nel kernel ufficiale solo a
partire dalla versione 2.6.21.\footnote{per il supporto deve essere stata
  abilitata l'opzione di compilazione \texttt{CONFIG\_HIGH\_RES\_TIMERS}, il
  supporto era però disponibile anche in precedenza nei patch facenti parte
  dello sviluppo delle estensioni \textit{real-time} del kernel, per cui
  alcune distribuzioni possono averlo anche con versioni precedenti del
  kernel.} Le funzioni definite dallo standard POSIX per gestire orologi ad
alta definizione però erano già presenti, essendo stata introdotte insieme ad
altre funzioni per il supporto delle estensioni \textit{real-time} con il
rilascio del kernel 2.6, ma la risoluzione effettiva era nominale.

A tutte le implementazioni che si rifanno a queste estensioni è richiesto di
disporre di una versione \textit{real-time} almeno per l'orologio generale di
sistema, quello che mantiene il \textit{calendar time} (vedi
sez.~\ref{sec:sys_time_base}), che in questa forma deve indicare il numero di
secondi e nanosecondi passati a partire dal primo gennaio 1970 (\textit{The
  Epoch}). Si ricordi infatti che l'orologio ordinario usato dal
\textit{calendar time} riporta solo un numero di secondi, e che la risoluzione
effettiva normalmente non raggiunge il nanosecondo (a meno di hardware
specializzato).  Oltre all'orologio generale di sistema possono essere
presenti altri tipi di orologi \textit{real-time}, ciascuno dei quali viene
identificato da un opportuno valore di una variabile di tipo
\type{clockid\_t}; un elenco di quelli disponibili su Linux è riportato in
tab.~\ref{tab:sig_timer_clockid_types}.

\begin{table}[htb]
  \footnotesize
  \centering
  \begin{tabular}[c]{|l|p{8cm}|}
    \hline
    \textbf{Valore} & \textbf{Significato} \\
    \hline
    \hline
    \constd{CLOCK\_REALTIME}    & Orologio \textit{real-time} di sistema, può
                                  essere impostato solo con privilegi
                                  amministrativi.\\ 
    \constd{CLOCK\_MONOTONIC}   & Orologio che indica un tempo monotono
                                  crescente (a partire da un tempo iniziale non
                                  specificato) che non può essere modificato e
                                  non cambia neanche in caso di reimpostazione
                                  dell'orologio di sistema.\\
    \constd{CLOCK\_PROCESS\_CPUTIME\_ID}& Contatore del tempo di CPU usato 
                                  da un processo (il \textit{process time} di
                                  sez.~\ref{sec:sys_cpu_times}, nel totale di
                                  \textit{system time} e \textit{user time})
                                  comprensivo di tutto il tempo di CPU usato
                                  da eventuali \textit{thread}.\\
    \constd{CLOCK\_THREAD\_CPUTIME\_ID}& Contatore del tempo di CPU
                                  (\textit{user time} e \textit{system time})
                                  usato da un singolo \textit{thread}.\\
    \hline
    \constd{CLOCK\_MONOTONIC\_RAW}&Simile al precedente, ma non subisce gli
                                  aggiustamenti dovuti all'uso di NTP (viene
                                  usato per fare riferimento ad una fonte
                                  hardware). Questo orologio è specifico di
                                  Linux, ed è disponibile a partire dal kernel
                                  2.6.28.\\
    \constd{CLOCK\_BOOTTIME}    & Identico a \const{CLOCK\_MONOTONIC} ma tiene
                                  conto anche del tempo durante il quale il
                                  sistema è stato sospeso (nel caso di
                                  sospensione in RAM o \textsl{ibernazione} su
                                  disco. Questo orologio è specifico di Linux,
                                  ed è disponibile a partire dal kernel
                                  2.6.39.\\
    \constd{CLOCK\_REALTIME\_ALARM}&Identico a \const{CLOCK\_REALTIME}, ma se
                                  usato per un timer il sistema sarà riattivato 
                                  anche se è in sospensione. Questo orologio è
                                  specifico di Linux, ed è disponibile a
                                  partire dal kernel 3.0.\\
    \constd{CLOCK\_BOOTTIME\_ALARM}&Identico a \const{CLOCK\_BOOTTIME}, ma se
                                  usato per un timer il sistema sarà riattivato 
                                  anche se è in sospensione. Questo orologio è
                                  specifico di Linux, ed è disponibile a
                                  partire dal kernel 3.0.\\
%    \const{}   & .\\
    \hline
  \end{tabular}
  \caption{Valori possibili per una variabile di tipo \typed{clockid\_t} 
    usata per indicare a quale tipo di orologio si vuole fare riferimento.}
  \label{tab:sig_timer_clockid_types}
\end{table}


% TODO: dal 4.17 CLOCK_MONOTONIC e CLOCK_BOOTTIME sono identici vedi
% https://lwn.net/Articles/751651/ e
% https://git.kernel.org/linus/d6ed449afdb38f89a7b38ec50e367559e1b8f71f
% change reverted, vedi: https://lwn.net/Articles/752757/

% NOTE: dal 3.0 anche i cosiddetti Posix Alarm Timers, con
% CLOCK_REALTIME_ALARM vedi http://lwn.net/Articles/429925/
% TODO: dal 3.10 anche CLOCK_TAI 

% TODO seguire l'evoluzione delle nuove syscall per il problema del 2038,
% iniziate ad entrare nel kernel dal 5.1, vedi
% https://lwn.net/Articles/776435/, https://lwn.net/Articles/782511/,
% https://git.kernel.org/linus/b1b988a6a035 

Per poter utilizzare queste funzionalità la \acr{glibc} richiede che la
macro \macro{\_POSIX\_C\_SOURCE} sia definita ad un valore maggiore o uguale
di \texttt{199309L} (vedi sez.~\ref{sec:intro_gcc_glibc_std}), inoltre i
programmi che le usano devono essere collegati con la libreria delle
estensioni \textit{real-time} usando esplicitamente l'opzione \texttt{-lrt}.

Si tenga presente inoltre che la disponibilità di queste funzionalità avanzate
può essere controllato dalla definizione della macro \macrod{\_POSIX\_TIMERS}
ad un valore maggiore di 0, e che le ulteriori macro
\macrod{\_POSIX\_MONOTONIC\_CLOCK}, \macrod{\_POSIX\_CPUTIME} e
\macrod{\_POSIX\_THREAD\_CPUTIME} indicano la presenza dei rispettivi orologi
di tipo \const{CLOCK\_MONOTONIC}, \const{CLOCK\_PROCESS\_CPUTIME\_ID} e
\const{CLOCK\_THREAD\_CPUTIME\_ID}; tutte queste macro sono definite in
\headfile{unistd.h}, che pertanto deve essere incluso per poterle
controllarle. Infine se il kernel ha il supporto per gli \textit{high
  resolution timer} un elenco degli orologi e dei timer può essere ottenuto
tramite il file \procfile{/proc/timer\_list}.

Le due funzioni che ci consentono rispettivamente di modificare o leggere il
valore per uno degli orologi \textit{real-time} sono \funcd{clock\_settime} e
\funcd{clock\_gettime}; i rispettivi prototipi sono:

\begin{funcproto}{
\fhead{time.h}
\fdecl{int clock\_settime(clockid\_t clockid, const struct timespec *tp)}
\fdesc{Imposta un orologio \textit{real-time}.} 
\fdecl{int clock\_gettime(clockid\_t clockid, struct timespec *tp)}
\fdesc{Legge un orologio \textit{real-time}.} 
}

{Le funzioni ritornano $0$ in caso di successo e $-1$ per un errore, nel qual
  caso \var{errno} assumerà uno dei valori: 
  \begin{errlist}
  \item[\errcode{EFAULT}] l'indirizzo \param{tp} non è valido.
  \item[\errcode{EINVAL}] il valore specificato per \param{clockid} non è
    valido o il relativo orologio \textit{real-time} non è supportato dal
    sistema.
  \item[\errcode{EPERM}] non si ha il permesso di impostare l'orologio
    indicato (solo per \func{clock\_settime}).
  \end{errlist}
}
\end{funcproto}

Entrambe le funzioni richiedono che si specifichi come primo argomento il tipo
di orologio su cui si vuole operare con uno dei valori di
tab.~\ref{tab:sig_timer_clockid_types} o con il risultato di una chiamata a
\func{clock\_getcpuclockid} (che tratteremo a breve), il secondo argomento
invece è sempre il puntatore \param{tp} ad una struttura \struct{timespec}
(vedi fig.~\ref{fig:sys_timespec_struct}) che deve essere stata
precedentemente allocata.  Per \func{clock\_settime} questa dovrà anche essere
stata inizializzata con il valore che si vuole impostare sull'orologio, mentre
per \func{clock\_gettime} verrà restituito al suo interno il valore corrente
dello stesso.

Si tenga presente inoltre che per eseguire un cambiamento sull'orologio
generale di sistema \const{CLOCK\_REALTIME} occorrono i privilegi
amministrativi;\footnote{ed in particolare la \textit{capability}
  \const{CAP\_SYS\_TIME}.} inoltre ogni cambiamento ad esso apportato non avrà
nessun effetto sulle temporizzazioni effettuate in forma relativa, come quelle
impostate sulle quantità di \textit{process time} o per un intervallo di tempo
da trascorrere, ma solo su quelle che hanno richiesto una temporizzazione ad
un istante preciso (in termini di \textit{calendar time}). Si tenga inoltre
presente che nel caso di Linux \const{CLOCK\_REALTIME} è l'unico orologio per
cui si può effettuare una modifica, infatti nonostante lo standard preveda la
possibilità di modifiche anche per \const{CLOCK\_PROCESS\_CPUTIME\_ID} e
\const{CLOCK\_THREAD\_CPUTIME\_ID}, il kernel non le consente.

Oltre alle due funzioni precedenti, lo standard POSIX prevede una terza
funzione di sistema che consenta di ottenere la risoluzione effettiva fornita
da un certo orologio, la funzione è \funcd{clock\_getres} ed il suo prototipo
è:

\begin{funcproto}{
\fhead{time.h}
\fdecl{int clock\_getres(clockid\_t clockid, struct timespec *res)}
\fdesc{Legge la risoluzione di un orologio \textit{real-time}.} 
}

{La funzione ritorna $0$ in caso di successo e $-1$ per un errore, nel qual
  caso \var{errno} assumerà uno dei valori: 
  \begin{errlist}
  \item[\errcode{EFAULT}] l'indirizzo di \param{res} non è valido.
  \item[\errcode{EINVAL}] il valore specificato per \param{clockid} non è
    valido.
  \end{errlist}
}
\end{funcproto}

La funzione richiede come primo argomento l'indicazione dell'orologio di cui
si vuole conoscere la risoluzione (effettuata allo stesso modo delle due
precedenti) e questa verrà restituita in una struttura \struct{timespec}
all'indirizzo puntato dall'argomento \param{res}.

Come accennato il valore di questa risoluzione dipende sia dall'hardware
disponibile che dalla implementazione delle funzioni, e costituisce il limite
minimo di un intervallo di tempo che si può indicare. Qualunque valore si
voglia utilizzare nelle funzioni di impostazione che non corrisponda ad un
multiplo intero di questa risoluzione, sarà troncato in maniera automatica. 

Gli orologi elencati nella seconda sezione di
tab.~\ref{tab:sig_timer_clockid_types} sono delle estensioni specifiche di
Linux, create per rispondere ad alcune esigenze specifiche, come quella di
tener conto di eventuali periodi di sospensione del sistema, e presenti solo
nelle versioni più recenti del kernel. In particolare gli ultimi due,
contraddistinti dal suffisso \texttt{\_ALARM}, hanno un impiego particolare,
derivato dalle esigenze emerse con Android per l'uso di Linux sui cellulari,
che consente di creare timer che possono scattare, riattivando il sistema,
anche quando questo è in sospensione. Per il loro utilizzo è prevista la
necessità di una capacità specifica, \const{CAP\_WAKE\_ALARM} (vedi
sez.~\ref{sec:proc_capabilities}).

Si tenga presente inoltre che con l'introduzione degli \textit{high resolution
  timer} i due orologi \const{CLOCK\_PROCESS\_CPUTIME\_ID} e
\const{CLOCK\_THREAD\_CPUTIME\_ID} fanno riferimento ai contatori presenti in
opportuni registri interni del processore; questo sui sistemi multiprocessore
può avere delle ripercussioni sulla precisione delle misure di tempo che vanno
al di là della risoluzione teorica ottenibile con \func{clock\_getres}, che
può essere ottenuta soltanto quando si è sicuri che un processo (o un
\textit{thread}) sia sempre stato eseguito sullo stesso processore.

Con i sistemi multiprocessore infatti ogni singola CPU ha i suoi registri
interni, e se ciascuna di esse utilizza una base di tempo diversa (se cioè il
segnale di temporizzazione inviato ai processori non ha una sola provenienza)
in genere ciascuna di queste potrà avere delle frequenze leggermente diverse,
e si otterranno pertanto dei valori dei contatori scorrelati fra loro, senza
nessuna possibilità di sincronizzazione.

Il problema si presenta, in forma più lieve, anche se la base di tempo è la
stessa, dato che un sistema multiprocessore non avvia mai tutte le CPU allo
stesso istante, si potrà così avere di nuovo una differenza fra i contatori,
soggetta però soltanto ad uno sfasamento costante. Per questo caso il kernel
per alcune architetture ha del codice che consente di ridurre al minimo la
differenza, ma non può essere comunque garantito che questa si annulli (anche
se in genere risulta molto piccola e trascurabile nella gran parte dei casi).

Per poter gestire questo tipo di problematiche lo standard ha previsto una
apposita funzione che sia in grado di ottenere l'identificativo dell'orologio
associato al \textit{process time} di un processo, la funzione è
\funcd{clock\_getcpuclockid} ed il suo prototipo è:

\begin{funcproto}{
\fhead{time.h}
\fdecl{int clock\_getcpuclockid(pid\_t pid, clockid\_t *clockid)}
\fdesc{Ottiene l'identificatore dell'orologio di CPU usato da un processo.} 
}

{La funzione ritorna $0$ in caso di successo ed un numero positivo per un
  errore, nel qual caso \var{errno} assumerà uno dei valori:
  \begin{errlist}
  \item[\errcode{ENOSYS}] non c'è il supporto per ottenere l'orologio relativo
    al \textit{process time} di un altro processo, e \param{pid} non
    corrisponde al processo corrente.
  \item[\errcode{EPERM}] il chiamante non ha il permesso di accedere alle
    informazioni relative al processo \param{pid}, avviene solo se è
    disponibile il supporto per leggere l'orologio relativo ad un altro
    processo.
  \item[\errcode{ESRCH}] non esiste il processo \param{pid}.
  \end{errlist}
}
\end{funcproto}

La funzione ritorna l'identificativo di un orologio di sistema associato ad un
processo indicato tramite l'argomento \param{pid}. Un utente normale, posto
che il kernel sia sufficientemente recente da supportare questa funzionalità,
può accedere soltanto ai dati relativi ai propri processi.

Del tutto analoga a \func{clock\_getcpuclockid}, ma da utilizzare per ottenere
l'orologio associato ad un \textit{thread} invece che a un processo, è
\funcd{pthread\_getcpuclockid},\footnote{per poterla utilizzare, come per
  qualunque funzione che faccia riferimento ai \textit{thread}, occorre
  effettuare il collegamento alla relativa libreria di gestione compilando il
  programma con \texttt{-lpthread}.} il cui prototipo è:

\begin{funcproto}{
\fhead{pthread.h}
\fhead{time.h}
\fdecl{int pthread\_getcpuclockid(pthread\_t thread, clockid\_t *clockid)}
\fdesc{Ottiene l'identificatore dell'orologio di CPU associato ad un
  \textit{thread}.} 
}

{La funzione ritorna $0$ in caso di successo ed un numero positivo per un
  errore, nel qual caso \var{errno} assumerà uno dei valori:
  \begin{errlist}
  \item[\errcode{ENOENT}] la funzione non è supportata dal sistema.
  \item[\errcode{ESRCH}] non esiste il \textit{thread} identificato
  \end{errlist}
 }
\end{funcproto}


% TODO, dal 2.6.39 aggiunta clock_adjtime 

Con l'introduzione degli orologi ad alta risoluzione è divenuto possibile
ottenere anche una gestione più avanzata degli allarmi; abbiamo già visto in
sez.~\ref{sec:sig_alarm_abort} come l'interfaccia di \func{setitimer} derivata
da BSD presenti delle serie limitazioni, come la possibilità di perdere un
segnale sotto carico, tanto che nello standard POSIX.1-2008 questa viene
marcata come obsoleta, e ne viene fortemente consigliata la sostituzione con
nuova interfaccia definita dallo standard POSIX.1-2001 che va sotto il nome di
\textit{POSIX Timer API}. Questa interfaccia è stata introdotta a partire dal
kernel 2.6, anche se il supporto di varie funzionalità da essa previste è
stato aggiunto solo in un secondo tempo.

Una delle principali differenze della nuova interfaccia è che un processo può
utilizzare un numero arbitrario di timer; questi vengono creati (ma non
avviati) tramite la funzione di sistema \funcd{timer\_create}, il cui
prototipo è:

\begin{funcproto}{
\fhead{signal.h}
\fhead{time.h}
\fdecl{int timer\_create(clockid\_t clockid, struct sigevent *evp,
    timer\_t *timerid)}
\fdesc{Crea un nuovo timer POSIX.} 
}

{La funzione ritorna $0$ in caso di successo e $-1$ per un errore, nel qual
  caso \var{errno} assumerà uno dei valori: 
  \begin{errlist}
  \item[\errcode{EAGAIN}] fallimento nel tentativo di allocare le strutture
    dei timer.
  \item[\errcode{EINVAL}] uno dei valori specificati per \param{clockid} o per
    i campi \var{sigev\_notify}, \var{sigev\_signo} o
    \var{sigev\_notify\_thread\_id} di \param{evp} non è valido.
  \item[\errcode{ENOMEM}] errore di allocazione della memoria.
  \end{errlist}
}
\end{funcproto}

La funzione richiede tre argomenti: il primo argomento serve ad indicare quale
tipo di orologio si vuole utilizzare e prende uno dei valori di
tab.~\ref{tab:sig_timer_clockid_types}; di detti valori però non è previsto
l'uso di \const{CLOCK\_MONOTONIC\_RAW} mentre
\const{CLOCK\_PROCESS\_CPUTIME\_ID} e \const{CLOCK\_THREAD\_CPUTIME\_ID} sono
disponibili solo a partire dal kernel 2.6.12. Si può così fare riferimento sia
ad un tempo assoluto che al tempo utilizzato dal processo (o \textit{thread})
stesso. Si possono inoltre utilizzare, posto di avere un kernel che li
supporti, gli orologi aggiuntivi della seconda parte di
tab.~\ref{tab:sig_timer_clockid_types}. 

Il secondo argomento richiede una trattazione più dettagliata, in quanto
introduce una struttura di uso generale, \struct{sigevent}, che viene
utilizzata anche da altre funzioni, come quelle per l'I/O asincrono (vedi
sez.~\ref{sec:file_asyncronous_io}) o le code di messaggi POSIX (vedi
sez.~\ref{sec:ipc_posix_mq}) e che serve ad indicare in maniera generica un
meccanismo di notifica.

\begin{figure}[!htb]
  \footnotesize \centering
  \begin{minipage}[c]{0.8\textwidth}
    \includestruct{listati/sigevent.h}
  \end{minipage} 
  \normalsize 
  \caption{La struttura \structd{sigevent}, usata per specificare in maniera
    generica diverse modalità di notifica degli eventi.}
  \label{fig:struct_sigevent}
\end{figure}

La struttura \struct{sigevent} (accessibile includendo \headfile{time.h}) è
riportata in fig.~\ref{fig:struct_sigevent}, la definizione effettiva dipende
dall'implementazione, quella mostrata è la versione descritta nella pagina di
manuale di \func{timer\_create}. Il campo \var{sigev\_notify} è il più
importante essendo quello che indica le modalità della notifica, gli altri
dipendono dal valore che si è specificato per \var{sigev\_notify}, si sono
riportati in tab.~\ref{tab:sigevent_sigev_notify}. La scelta del meccanismo di
notifica viene fatta impostando uno dei valori di
tab.~\ref{tab:sigevent_sigev_notify} per \var{sigev\_notify}, e fornendo gli
eventuali ulteriori argomenti necessari a secondo della scelta
effettuata. Diventa così possibile indicare l'uso di un segnale o l'esecuzione
(nel caso di uso dei \textit{thread}) di una funzione di modifica in un
\textit{thread} dedicato.

\begin{table}[htb]
  \footnotesize
  \centering
  \begin{tabular}[c]{|l|p{10cm}|}
    \hline
    \textbf{Valore} & \textbf{Significato} \\
    \hline
    \hline
    \constd{SIGEV\_NONE}   & Non viene inviata nessuna notifica.\\
    \constd{SIGEV\_SIGNAL} & La notifica viene effettuata inviando al processo
                             chiamante il segnale specificato dal campo
                             \var{sigev\_signo}; se il gestore di questo
                             segnale è stato installato con
                             \const{SA\_SIGINFO} gli verrà restituito il
                             valore specificato con \var{sigev\_value} (una
                             \dirct{union} \texttt{sigval}, la cui definizione
                             è in fig.~\ref{fig:sig_sigval}) come valore del
                             campo \var{si\_value} di \struct{siginfo\_t}.\\
    \constd{SIGEV\_THREAD} & La notifica viene effettuata creando un nuovo
                             \textit{thread} che esegue la funzione di
                             notifica specificata da 
                             \var{sigev\_notify\_function} con argomento
                             \var{sigev\_value}. Se questo è diverso da
                             \val{NULL}, il \textit{thread} viene creato con
                             gli attributi specificati da
                             \var{sigev\_notify\_attribute}.\footnotemark\\
    \constd{SIGEV\_THREAD\_ID}& Invia la notifica come segnale (con le stesse
                             modalità di \const{SIGEV\_SIGNAL}) che però viene
                             recapitato al \textit{thread} indicato dal campo
                             \var{sigev\_notify\_thread\_id}. Questa modalità
                             è una estensione specifica di Linux, creata come
                             supporto per le librerie di gestione dei
                             \textit{thread}, pertanto non deve essere usata
                             da codice normale.\\
    \hline
  \end{tabular}
  \caption{Valori possibili per il campo \var{sigev\_notify} in una struttura
    \struct{sigevent}.} 
  \label{tab:sigevent_sigev_notify}
\end{table}

\footnotetext{nel caso dei \textit{timer} questa funzionalità è considerata un
  esempio di pessima implementazione di una interfaccia, richiesta dallo
  standard POSIX, ma da evitare totalmente nell'uso ordinario, a causa della
  possibilità di creare disservizi generando una gran quantità di processi,
  tanto che ne è stata richiesta addirittura la rimozione.}

Nel caso di \func{timer\_create} occorrerà passare alla funzione come secondo
argomento l'indirizzo di una di queste strutture per indicare le modalità con
cui si vuole essere notificati della scadenza del timer, se non si specifica
nulla (passando un valore \val{NULL}) verrà inviato il segnale
\signal{SIGALRM} al processo corrente, o per essere più precisi verrà
utilizzato un valore equivalente all'aver specificato \const{SIGEV\_SIGNAL}
per \var{sigev\_notify}, \signal{SIGALRM} per \var{sigev\_signo} e
l'identificatore del timer come valore per \var{sigev\_value.sival\_int}.

Il terzo argomento deve essere l'indirizzo di una variabile di tipo
\typed{timer\_t} dove sarà scritto l'identificativo associato al timer appena
creato, da usare in tutte le successive funzioni di gestione. Una volta creato
questo identificativo resterà univoco all'interno del processo stesso fintanto
che il timer non viene cancellato.

Si tenga presente che eventuali POSIX timer creati da un processo non vengono
ereditati dai processi figli creati con \func{fork} e che vengono cancellati
nella esecuzione di un programma diverso attraverso una delle funzioni
\func{exec}. Si tenga presente inoltre che il kernel prealloca l'uso di un
segnale \textit{real-time} per ciascun timer che viene creato con
\func{timer\_create}; dato che ciascuno di essi richiede un posto nella coda
dei segnali \textit{real-time}, il numero massimo di timer utilizzabili da un
processo è limitato dalle dimensioni di detta coda, ed anche, qualora questo
sia stato impostato, dal limite \const{RLIMIT\_SIGPENDING}.

Una volta creato il timer \func{timer\_create} ed ottenuto il relativo
identificatore, si può attivare o disattivare un allarme (in gergo
\textsl{armare} o \textsl{disarmare} il timer) con la funzione di sistema
\funcd{timer\_settime}, il cui prototipo è:

\begin{funcproto}{
\fhead{signal.h}
\fhead{time.h}
\fdecl{int timer\_settime(timer\_t timerid, int flags, const struct
    itimerspec *new\_value, struct itimerspec *old\_value)}
\fdesc{Arma o disarma un timer POSIX.} 
}

{La funzione ritorna $0$ in caso di successo e $-1$ per un errore, nel qual
  caso \var{errno} assumerà uno dei valori: 
  \begin{errlist}
  \item[\errcode{EFAULT}] si è specificato un indirizzo non valido
    per \param{new\_value} o \param{old\_value}.
  \item[\errcode{EINVAL}] all'interno di \param{new\_value.value} si è
    specificato un tempo negativo o un numero di nanosecondi maggiore di
    999999999.
  \end{errlist}
}
\end{funcproto}

La funzione richiede che si indichi la scadenza del timer con
l'argomento \param{new\_value}, che deve essere specificato come puntatore ad
una struttura di tipo \struct{itimerspec}, la cui definizione è riportata in
fig.~\ref{fig:struct_itimerspec}; se il puntatore \param{old\_value} è diverso
da \val{NULL} il valore corrente della scadenza verrà restituito in una
analoga struttura, ovviamente in entrambi i casi le strutture devono essere
state allocate.

\begin{figure}[!htb]
  \footnotesize \centering
  \begin{minipage}[c]{0.8\textwidth}
    \includestruct{listati/itimerspec.h}
  \end{minipage} 
  \normalsize 
  \caption{La struttura \structd{itimerspec}, usata per specificare la
    scadenza di un allarme.}
  \label{fig:struct_itimerspec}
\end{figure}

Ciascuno dei due campi di \struct{itimerspec} indica un tempo, da specificare
con una precisione fino al nanosecondo tramite una struttura \struct{timespec}
(la cui definizione è riportata fig.~\ref{fig:sys_timespec_struct}). Il campo
\var{it\_value} indica la prima scadenza dell'allarme. Di default, quando il
valore di \param{flags} è nullo, questo valore viene considerato come un
intervallo relativo al tempo corrente, il primo allarme scatterà cioè dopo il
numero di secondi e nanosecondi indicati da questo campo. Se invece si usa
per \param{flags} il valore \constd{TIMER\_ABSTIME}, che al momento è l'unico
valore valido per \param{flags}, allora \var{it\_value} viene considerato come
un valore assoluto rispetto al valore usato dall'orologio a cui è associato il
timer. 

Quindi a seconda dei casi si potrà impostare un timer o con un tempo assoluto,
quando si opera rispetto all'orologio di sistema (nel qual caso il valore deve
essere in secondi e nanosecondi dalla \textit{epoch}) o con un numero di
secondi o nanosecondi rispetto alla partenza di un orologio di CPU, quando si
opera su uno di questi.  Infine un valore nullo di \var{it\_value}, dove per
nullo si intende con valori nulli per entrambi i campi \var{tv\_sec} e
\var{tv\_nsec}, può essere utilizzato, indipendentemente dal tipo di orologio
utilizzato, per disarmare l'allarme.

Il campo \var{it\_interval} di \struct{itimerspec} viene invece utilizzato per
impostare un allarme periodico.  Se il suo valore è nullo, se cioè sono nulli
tutti e due i due campi \var{tv\_sec} e \var{tv\_nsec} di detta struttura
\struct{timespec}, l'allarme scatterà una sola volta secondo quando indicato
con \var{it\_value}, altrimenti il valore specificato nella struttura verrà
preso come l'estensione del periodo di ripetizione della generazione
dell'allarme, che proseguirà indefinitamente fintanto che non si disarmi il
timer.

Se il timer era già stato armato la funzione sovrascrive la precedente
impostazione, se invece si indica come prima scadenza un tempo già passato,
l'allarme verrà notificato immediatamente e al contempo verrà incrementato il
contatore dei superamenti. Questo contatore serve a fornire una indicazione al
programma che riceve l'allarme su un eventuale numero di scadenze che sono
passate prima della ricezione della notifica dell'allarme. 

É infatti possibile, qualunque sia il meccanismo di notifica scelto, che
quest'ultima venga ricevuta dopo che il timer è scaduto più di una volta,
specialmente se si imposta un timer con una ripetizione a frequenza
elevata. Nel caso dell'uso di un segnale infatti il sistema mette in coda un
solo segnale per timer,\footnote{questo indipendentemente che si tratti di un
  segnale ordinario o \textit{real-time}, per questi ultimi sarebbe anche
  possibile inviare un segnale per ogni scadenza, questo però non viene fatto
  per evitare il rischio, tutt'altro che remoto, di riempire la coda.} e se il
sistema è sotto carico o se il segnale è bloccato, prima della sua ricezione
può passare un intervallo di tempo sufficientemente lungo ad avere scadenze
multiple, e lo stesso può accadere anche se si usa un \textit{thread} di
notifica.

Per questo motivo il gestore del segnale o il \textit{thread} di notifica può
ottenere una indicazione di quante volte il timer è scaduto dall'invio della
notifica utilizzando la funzione di sistema \funcd{timer\_getoverrun}, il cui
prototipo è:

\begin{funcproto}{
\fhead{time.h}
\fdecl{int timer\_getoverrun(timer\_t timerid)}
\fdesc{Ottiene il numero di scadenze di un timer POSIX.} 
}

{La funzione ritorna il numero di scadenze di un timer in caso di successo e
  $-1$ per un errore, nel qual caso \var{errno} assumerà uno dei valori:
  \begin{errlist}
  \item[\errcode{EINVAL}] \param{timerid} non indica un timer valido.
  \end{errlist}
}
\end{funcproto}

La funzione ritorna il numero delle scadenze avvenute, che può anche essere
nullo se non ve ne sono state. Come estensione specifica di Linux,\footnote{in
  realtà lo standard POSIX.1-2001 prevede gli \textit{overrun} solo per i
  segnali e non ne parla affatto in riferimento ai \textit{thread}.}  quando
si usa un segnale come meccanismo di notifica, si può ottenere direttamente
questo valore nel campo \var{si\_overrun} della struttura \struct{siginfo\_t}
(illustrata in fig.~\ref{fig:sig_siginfo_t}) restituita al gestore del segnale
installato con \func{sigaction}; in questo modo non è più necessario eseguire
successivamente una chiamata a questa funzione per ottenere il numero delle
scadenze. Al gestore del segnale viene anche restituito, come ulteriore
informazione, l'identificativo del timer, in questo caso nel campo
\var{si\_timerid}.

Qualora si voglia rileggere lo stato corrente di un timer, ed ottenere il
tempo mancante ad una sua eventuale scadenza, si deve utilizzare la funzione
di sistema \funcd{timer\_gettime}, il cui prototipo è:

\begin{funcproto}{
\fhead{time.h}
\fdecl{int timer\_gettime(timer\_t timerid, int flags, struct
    itimerspec *curr\_value)}
\fdesc{Legge lo stato di un timer POSIX.} 
}

{La funzione ritorna $0$ in caso di successo e $-1$ per un errore, nel qual
  caso \var{errno} assumerà uno dei valori: 
  \begin{errlist}
  \item[\errcode{EFAULT}] si è specificato un indirizzo non valido
    per \param{curr\_value}.
  \item[\errcode{EINVAL}] \param{timerid} non indica un timer valido.
  \end{errlist}
}
\end{funcproto}

La funzione restituisce nella struttura \struct{itimerspec} puntata
da \param{curr\_value} il tempo restante alla prossima scadenza nel campo
\var{it\_value}. Questo tempo viene sempre indicato in forma relativa, anche
nei casi in cui il timer era stato precedentemente impostato con
\const{TIMER\_ABSTIME} indicando un tempo assoluto.  Il ritorno di un valore
nullo nel campo \var{it\_value} significa che il timer è disarmato o è
definitivamente scaduto. 

Nel campo \var{it\_interval} di \param{curr\_value} viene invece restituito,
se questo era stato impostato, il periodo di ripetizione del timer.  Anche in
questo caso il ritorno di un valore nullo significa che il timer non era stato
impostato per una ripetizione e doveva operare, come suol dirsi, a colpo
singolo (in gergo \textit{one shot}).

Infine, quando un timer non viene più utilizzato, lo si può cancellare,
rimuovendolo dal sistema e recuperando le relative risorse, effettuando in
sostanza l'operazione inversa rispetto a \func{timer\_create}. Per questo
compito lo standard prevede una apposita funzione di sistema,
\funcd{timer\_delete}, il cui prototipo è:

\begin{funcproto}{
\fhead{time.h}
\fdecl{int timer\_delete(timer\_t timerid)}
\fdesc{Cancella un timer POSIX.} 
}

{La funzione ritorna $0$ in caso di successo e $-1$ per un errore, nel qual
  caso \var{errno} assumerà uno dei valori: 
  \begin{errlist}
    \item[\errcode{EINVAL}] \param{timerid} non indica un timer valido.
  \end{errlist}
}
\end{funcproto}

La funzione elimina il timer identificato da \param{timerid}, disarmandolo se
questo era stato attivato. Nel caso, poco probabile ma comunque possibile, che
un timer venga cancellato prima della ricezione del segnale pendente per la
notifica di una scadenza, il comportamento del sistema è indefinito.

Infine a partire dal kernel 2.6 e per le versioni della \acr{libc} superiori
alla 2.1, si può utilizzare la nuova interfaccia dei timer POSIX anche per le
funzioni di attesa, per questo è disponibile la funzione di sistema
\funcd{clock\_nanosleep}, il cui prototipo è:

\begin{funcproto}{
\fhead{time.h}
\fdecl{int clock\_nanosleep(clockid\_t clock\_id, int flags, const struct
    timespec *request,\\
\phantom{int clock\_nanosleep(}struct timespec *remain)}
\fdesc{Pone il processo in pausa per un tempo specificato.}
}

{La funzione ritorna $0$ in caso di successo ed un valore positivo per un
  errore, espresso dai valori:
  \begin{errlist}
    \item[\errcode{EINTR}] la funzione è stata interrotta da un segnale.
    \item[\errcode{EINVAL}] si è specificato un numero di secondi negativo o
      un numero di nanosecondi maggiore di 999.999.999 o indicato un orologio
      non valido.
  \end{errlist}
  ed inoltre \errval{EFAULT} nel suo significato generico.}
\end{funcproto}

I due argomenti \param{request} e \param{remain} sono identici agli analoghi di
\func{nanosleep} che abbiamo visto in sez.~\ref{sec:sig_pause_sleep}, ed hanno
lo stesso significato.  L'argomento \param{clock\_id} consente di indicare
quale orologio si intende utilizzare per l'attesa con uno dei valori della
prima parte di tab.~\ref{tab:sig_timer_clockid_types} (eccetto
\const{CLOCK\_THREAD\_CPUTIME\_ID}).  L'argomento \param{flags} consente di
modificare il comportamento della funzione, il suo unico valore valido al
momento è \const{TIMER\_ABSTIME} che, come per \func{timer\_settime} indica di
considerare il tempo indicato in \param{request} come assoluto anziché
relativo.

Il comportamento della funzione è analogo a \func{nanosleep}, se la chiamata
viene interrotta il tempo rimanente viene restituito in \param{remain}.
Utilizzata normalmente con attese relative può soffrire degli stessi problemi
di deriva di cui si è parlato in sez.~\ref{sec:sig_pause_sleep} dovuti ad
interruzioni ripetute per via degli arrotondamenti fatti a questo tempo.  Ma
grazie alla possibilità di specificare tempi assoluti con \param{flags} si può
ovviare a questo problema ricavando il tempo corrente con
\func{clock\_gettime}, aggiungendovi l'intervallo di attesa, ed impostando
questa come tempo assoluto.

Si tenga presente che se si è usato il valore \const{TIMER\_ABSTIME}
per \param{flags} e si è indicato un tempo assoluto che è già passato la
funzione ritorna immediatamente senza nessuna sospensione. In caso di
interruzione da parte di un segnale il tempo rimanente viene restituito
in \param{remain} soltanto se questo non è un puntatore \val{NULL} e non si è
specificato \const{TIMER\_ABSTIME} per \param{flags}.


\itindend{POSIX~Timer~API}



\subsection{Ulteriori funzioni di gestione}
\label{sec:sig_specific_features}

In questo ultimo paragrafo esamineremo le rimanenti funzioni di gestione dei
segnali non descritte finora, relative agli aspetti meno utilizzati e più
``\textsl{esoterici}'' della interfaccia.

% TODO: trattare (qui?) pidfd_send_signal() introdotta con il kernel 5.1 vedi
% https://lwn.net/Articles/784831/ e https://lwn.net/Articles/773459/

La prima di queste funzioni è la funzione di sistema \funcd{sigpending},
anch'essa introdotta dallo standard POSIX.1, il suo prototipo è:

\begin{funcproto}{
\fhead{signal.h}
\fdecl{int sigpending(sigset\_t *set)}
\fdesc{Legge l'insieme dei segnali pendenti.} 
}

{La funzione ritorna $0$ in caso di successo e $-1$ per un errore, nel qual
  caso \var{errno} assumerà solo il valore \errcode{EFAULT} nel suo
  significato generico.}
\end{funcproto}

La funzione permette di ricavare quali sono i segnali pendenti per il processo
in corso, cioè i segnali che sono stati inviati dal kernel ma non sono stati
ancora ricevuti dal processo in quanto bloccati. Non esiste una funzione
equivalente nella vecchia interfaccia, ma essa è tutto sommato poco utile,
dato che essa può solo assicurare che un segnale è stato inviato, dato che
escluderne l'avvenuto invio al momento della chiamata non significa nulla
rispetto a quanto potrebbe essere in un qualunque momento successivo.

Una delle caratteristiche di BSD, disponibile anche in Linux, è la possibilità
di usare uno \textit{stack} alternativo per i segnali; è cioè possibile fare
usare al sistema un altro \textit{stack} (invece di quello relativo al
processo, vedi sez.~\ref{sec:proc_mem_layout}) solo durante l'esecuzione di un
gestore.  L'uso di uno \textit{stack} alternativo è del tutto trasparente ai
gestori, occorre però seguire una certa procedura:
\begin{enumerate*}
\item allocare un'area di memoria di dimensione sufficiente da usare come
  \textit{stack} alternativo;
\item usare la funzione \func{sigaltstack} per rendere noto al sistema
  l'esistenza e la locazione dello \textit{stack} alternativo;
\item quando si installa un gestore occorre usare \func{sigaction}
  specificando il flag \const{SA\_ONSTACK} (vedi tab.~\ref{tab:sig_sa_flag})
  per dire al sistema di usare lo \textit{stack} alternativo durante
  l'esecuzione del gestore.
\end{enumerate*}

In genere il primo passo viene effettuato allocando un'opportuna area di
memoria con \code{malloc}; in \headfile{signal.h} sono definite due costanti,
\constd{SIGSTKSZ} e \constd{MINSIGSTKSZ}, che possono essere utilizzate per
allocare una quantità di spazio opportuna, in modo da evitare overflow. La
prima delle due è la dimensione canonica per uno \textit{stack} di segnali e
di norma è sufficiente per tutti gli usi normali.

La seconda è lo spazio che occorre al sistema per essere in grado di lanciare
il gestore e la dimensione di uno \textit{stack} alternativo deve essere
sempre maggiore di questo valore. Quando si conosce esattamente quanto è lo
spazio necessario al gestore gli si può aggiungere questo valore per allocare
uno \textit{stack} di dimensione sufficiente.

Come accennato, per poter essere usato, lo \textit{stack} per i segnali deve
essere indicato al sistema attraverso la funzione \funcd{sigaltstack}; il suo
prototipo è:

\begin{funcproto}{
\fhead{signal.h}
\fdecl{int sigaltstack(const stack\_t *ss, stack\_t *oss)}
\fdesc{Installa uno \textit{stack} alternativo per i gestori di segnali.} 
}

{La funzione ritorna $0$ in caso di successo e $-1$ per un errore, nel qual
  caso \var{errno} assumerà uno dei valori: 
  \begin{errlist}
  \item[\errcode{EFAULT}] uno degli indirizzi degli argomenti non è valido.
  \item[\errcode{EINVAL}] \param{ss} non è nullo e \var{ss\_flags} contiene un
  valore diverso da zero che non è \const{SS\_DISABLE}.
  \item[\errcode{ENOMEM}] la dimensione specificata per il nuovo
    \textit{stack} è minore di \const{MINSIGSTKSZ}.
  \item[\errcode{EPERM}] si è cercato di cambiare lo \textit{stack}
    alternativo mentre questo è attivo (cioè il processo è in esecuzione su di
    esso).
  \end{errlist}
}
\end{funcproto}

La funzione prende come argomenti puntatori ad una struttura di tipo
\var{stack\_t}, definita in fig.~\ref{fig:sig_stack_t}. I due valori
\param{ss} e \param{oss}, se non nulli, indicano rispettivamente il nuovo
\textit{stack} da installare e quello corrente (che viene restituito dalla
funzione per un successivo ripristino).

\begin{figure}[!htb]
  \footnotesize \centering
  \begin{minipage}[c]{0.8\textwidth}
    \includestruct{listati/stack_t.h}
  \end{minipage} 
  \normalsize 
  \caption{La struttura \structd{stack\_t}.} 
  \label{fig:sig_stack_t}
\end{figure}

Il campo \var{ss\_sp} di \struct{stack\_t} indica l'indirizzo base dello
\textit{stack}, mentre \var{ss\_size} ne indica la dimensione; il campo
\var{ss\_flags} invece indica lo stato dello \textit{stack}.  Nell'indicare un
nuovo \textit{stack} occorre inizializzare \var{ss\_sp} e \var{ss\_size}
rispettivamente al puntatore e alla dimensione della memoria allocata, mentre
\var{ss\_flags} deve essere nullo.  Se invece si vuole disabilitare uno
\textit{stack} occorre indicare \constd{SS\_DISABLE} come valore di
\var{ss\_flags} e gli altri valori saranno ignorati.

Se \param{oss} non è nullo verrà restituito dalla funzione indirizzo e
dimensione dello \textit{stack} corrente nei relativi campi, mentre
\var{ss\_flags} potrà assumere il valore \constd{SS\_ONSTACK} se il processo è
in esecuzione sullo \textit{stack} alternativo (nel qual caso non è possibile
cambiarlo) e \const{SS\_DISABLE} se questo non è abilitato.

In genere si installa uno \textit{stack} alternativo per i segnali quando si
teme di avere problemi di esaurimento dello \textit{stack} standard o di
superamento di un limite (vedi sez.~\ref{sec:sys_resource_limit}) imposto con
chiamate del tipo \code{setrlimit(RLIMIT\_STACK, \&rlim)}.  In tal caso
infatti si avrebbe un segnale di \signal{SIGSEGV}, che potrebbe essere gestito
soltanto avendo abilitato uno \textit{stack} alternativo.

Si tenga presente che le funzioni chiamate durante l'esecuzione sullo
\textit{stack} alternativo continueranno ad usare quest'ultimo, che, al
contrario di quanto avviene per lo \textit{stack} ordinario dei processi, non
si accresce automaticamente (ed infatti eccederne le dimensioni può portare a
conseguenze imprevedibili).  Si ricordi infine che una chiamata ad una
funzione della famiglia \func{exec} cancella ogni \textit{stack} alternativo.

Abbiamo visto in fig.~\ref{fig:sig_sleep_incomplete} come si possa usare
\func{longjmp} per uscire da un gestore rientrando direttamente nel corpo
del programma, sappiamo però che nell'esecuzione di un gestore il segnale
che l'ha invocato viene bloccato, e abbiamo detto che possiamo ulteriormente
modificarlo con \func{sigprocmask}. 

Resta quindi il problema di cosa succede alla maschera dei segnali quando si
esce da un gestore usando questa funzione. Il comportamento dipende
dall'implementazione. In particolare la semantica usata da BSD prevede che sia
ripristinata la maschera dei segnali precedente l'invocazione, come per un
normale ritorno, mentre quella usata da System V no.

Lo standard POSIX.1 non specifica questo comportamento per \func{setjmp} e
\func{longjmp}, ed il comportamento della \acr{glibc} dipende da quale delle
caratteristiche si sono abilitate con le macro viste in
sez.~\ref{sec:intro_gcc_glibc_std}.

Lo standard POSIX però prevede anche la presenza di altre due funzioni
\funcd{sigsetjmp} e \funcd{siglongjmp}, che permettono di decidere in maniera
esplicita quale dei due comportamenti il programma deve assumere; i loro
prototipi sono:

\begin{funcproto}{
\fhead{setjmp.h}
\fdecl{int sigsetjmp(sigjmp\_buf env, int savesigs)}
\fdesc{Salva il contesto dello \textit{stack} e la maschera dei segnali.}  
\fdecl{void siglongjmp(sigjmp\_buf env, int val)}
\fdesc{Ripristina il contesto dello \textit{stack} e la maschera dei segnali.} 
}

{La funzioni sono identiche alle analoghe \func{setjmp} e \func{longjmp} di
  sez.~\ref{sec:proc_longjmp} ed hanno gli stessi errori e valori di uscita.}
\end{funcproto}

Le due funzioni prendono come primo argomento la variabile su cui viene
salvato il contesto dello \textit{stack} per permettere il salto non-locale;
nel caso specifico essa è di tipo \typed{sigjmp\_buf}, e non \type{jmp\_buf}
come per le analoghe di sez.~\ref{sec:proc_longjmp} in quanto in questo caso
viene salvata anche la maschera dei segnali.

Nel caso di \func{sigsetjmp}, se si specifica un valore di \param{savesigs}
diverso da zero la maschera dei valori verrà salvata in \param{env} insieme al
contesto dello \textit{stack} usato per il salto non locale. Se così si è
fatto la maschera dei segnali verrà ripristinata in una successiva chiamata a
\func{siglongjmp}. Quest'ultima, a parte l'uso di un valore di \param{env} di
tipo \type{sigjmp\_buf}, è assolutamente identica a \func{longjmp}.


% TODO: se e quando si troverà un argomento adeguato inserire altre funzioni
% sparse attinenti ai segnali, al momento sono note solo:
% * sigreturn (funzione interna, scarsamente interessante)
% argomento a priorità IDLE (fare quando non resta niente altro da trattare)


% LocalWords:  kernel POSIX timer shell control ctrl kill raise signal handler
% LocalWords:  reliable unreliable fig race condition sez struct process table
% LocalWords:  delivered pending scheduler sigpending l'I suspend SIGKILL wait
% LocalWords:  SIGSTOP sigaction waitpid dump stack debugger nell'header NSIG
% LocalWords:  tab BSD SUSv SIGHUP PL Hangup SIGINT Interrupt SIGQUIT Quit AEF
% LocalWords:  SIGILL SIGABRT abort SIGFPE SIGSEGV SIGPIPE SIGALRM alarm SIGUSR
% LocalWords:  SIGTERM SIGCHLD SIGCONT SIGTSTP SIGTTIN SIGTTOU SIGBUS bad SL of
% LocalWords:  memory access SIGPOLL Pollable event Sys SIGIO SIGPROF profiling
% LocalWords:  SIGSYS SVID SIGTRAP breakpoint SIGURG urgent socket Virtual IOT
% LocalWords:  clock SIGXCPU SIGXFSZ SIGIOT trap SIGEMT SIGSTKFLT SIGCLD SIGPWR
% LocalWords:  SIGINFO SIGLOST lock NFS SIGWINCH Sun SIGUNUSED fault point heap
% LocalWords:  exception l'overflow illegal instruction overflow segment error
% LocalWords:  violation system call interrupt INTR hang SIGVTALRM virtual SUSP
% LocalWords:  profilazione fcntl descriptor sleep interactive Broken FIFO lost
% LocalWords:  EPIPE Resource advisory client limit exceeded size window change
% LocalWords:  strsignal psignal SOURCE strerror string char int signum perror
% LocalWords:  void sig const sys siglist L'array decr fork exec DFL IGN ioctl
% LocalWords:  EINTR glibc TEMP FAILURE RETRY expr multitasking SVr sighandler
% LocalWords:  ERR libc bsd sysv XOPEN EINVAL pid errno ESRCH EPERM getpid init
% LocalWords:  killpg pidgrp group unistd unsigned seconds all' setitimer which
% LocalWords:  itimerval value ovalue EFAULT ITIMER it interval timeval ms VIRT
% LocalWords:  getitimer stdlib stream atexit exit usleep long usec nanosleep
% LocalWords:  timespec req rem HZ scheduling SCHED RR SigHand forktest WNOHANG
% LocalWords:  deadlock longjmp setjmp sigset sigemptyset sigfillset sigaddset
% LocalWords:  sigdelset sigismember act oldact restorer mask NOCLDSTOP ONESHOT
% LocalWords:  RESETHAND RESTART NOMASK NODEFER ONSTACK sigcontext union signo
% LocalWords:  siginfo bits uid addr fd inline like blocked atomic sigprocmask
% LocalWords:  how oldset BLOCK UNBLOCK SETMASK sigsuspend sigaltstack malloc
% LocalWords:  SIGSTKSZ MINSIGSTKSZ ss oss ENOMEM flags DISABLE sp setrlimit LB
% LocalWords:  RLIMIT rlim sigsetjmp siglongjmp sigjmp buf env savesigs jmp ptr
% LocalWords:  SIGRTMIN SIGRTMAX sigval sigevent sigqueue EAGAIN sysctl safe tp
% LocalWords:  QUEUE thread sigwait sigwaitinfo sigtimedwait info DEF SLB bind
% LocalWords:  function accept return cfgetispeed cfgetospeed cfsetispeed chdir
% LocalWords:  cfsetospeed chmod chown gettime close connect creat dup execle
% LocalWords:  execve fchmod fchown fdatasync fpathconf fstat fsync ftruncate
% LocalWords:  getegid geteuid getgid getgroups getpeername getpgrp getppid sem
% LocalWords:  getsockname getsockopt getuid listen lseek lstat mkdir mkfifo tv
% LocalWords:  pathconf poll posix pselect read readlink recv recvfrom recvmsg
% LocalWords:  rename rmdir select send sendmsg sendto setgid setpgid setsid
% LocalWords:  setsockopt setuid shutdown sigpause socketpair stat symlink page
% LocalWords:  sysconf tcdrain tcflow tcflush tcgetattr tcgetgrp tcsendbreak
% LocalWords:  tcsetattr tcsetpgrp getoverrun times umask uname unlink utime
% LocalWords:  write sival SIVGTALRM NOCLDWAIT MESGQ ASYNCIO TKILL tkill tgkill
% LocalWords:  ILL ILLOPC ILLOPN ILLADR ILLTRP PRVOPC PRVREG COPROC BADSTK FPE
% LocalWords:  INTDIV INTOVF FLTDIV FLTOVF FLTUND underflow FLTRES FLTINV SEGV
% LocalWords:  FLTSUB MAPERR ACCERR ADRALN ADRERR OBJERR BRKPT CLD EXITED MSG
% LocalWords:  KILLED DUMPED TRAPPED STOPPED CONTINUED PRI HUP SigFunc jiffies
% LocalWords:  SEC unsafe sockatmark execl execv faccessat fchmodat fchownat
% LocalWords:  fexecve fstatat futimens linkat mkdirat mkfifoat mknod mknodat
% LocalWords:  openat readlinkat renameat symlinkat unlinkat utimensat utimes
% LocalWords:  LinuxThread NTPL Library clockid evp timerid sigev notify high
% LocalWords:  resolution CONFIG RES patch REALTIME MONOTONIC RAW NTP CPUTIME
% LocalWords:  tick calendar The Epoch list getcpuclockid capability CAP getres
% LocalWords:  ENOSYS pthread ENOENT NULL attribute itimerspec new old ABSTIME
% LocalWords:  epoch multiplexing overrun res lpthread sec nsec curr one shot
% LocalWords:  delete stopped gdb alpha mips emulation locking ppoll epoll PGID
% LocalWords:  pwait msgrcv msgsnd semop semtimedop runnable sigisemptyset HRT
% LocalWords:  sigorset sigandset BOOTTIME Android request remain


%%% Local Variables: 
%%% mode: latex
%%% TeX-master: "gapil"
%%% End: 
